% Fragmento maestro: incluir desde el anfitrion, con \HMTFGRoot fijada a la raiz del paquete.
% Requiere amsmath, amssymb/mathrsfs o unicode-math, amsthm, booktabs, longtable, hyperref, needspace.
% El anfitrion conserva sus entornos theorem, proposition, lemma y proof.
\providecommand{\HMTFGRoot}{.}
\begingroup
\providecommand{\ZZ}{}\renewcommand{\ZZ}{\mathbb Z}
\providecommand{\RR}{}\renewcommand{\RR}{\mathbb R}
\providecommand{\FF}{}\renewcommand{\FF}{\mathbb F}
\providecommand{\APP}{}\renewcommand{\APP}{\mathrm{APP}}
\providecommand{\TPK}{}\renewcommand{\TPK}{\mathrm{TPK}}
\providecommand{\TRIT}{}\renewcommand{\TRIT}{\{-1,0,+1\}}
% unicode-math usa lBrack/rBrack; newtxmath ya proporciona llbracket/rrbracket.
\providecommand{\llbracket}{\lBrack}
\providecommand{\rrbracket}{\rBrack}
% Fragmento de inserción; se incluye desde la raíz de este paquete editorial.
% Las fuentes completas y los cortes materiales constan en ANTECEDENTES_Y_DEPENDENCIAS.md.
% Procedencia: XI REV11 ES, sections/nucleo.tex, tpk_desarrollo_integrado.tex,
% ch_tres_vueltas_ext.tex, ch_elevacion_catalogal.tex, ch_cierre_cardinal.tex,
% ch_descenso_por_fibras.tex; perfil integral public_final/feigenbaum/c37_feigenbaum.tex;
% Grassmannianos REV07 ES, sections/jacobi_spectral.tex;
% TRANSVERSALIDAD_Y_ESCALADO_LOCAL_FEIGENBAUM.md, sección 14.
% Ampliación expositiva de demostraciones existentes; no introduce resultados nuevos.
\providecommand{\ZZ}{\mathbb Z}
\providecommand{\RR}{\mathbb R}
\providecommand{\FF}{\mathbb F}
\providecommand{\APP}{\mathrm{APP}}
\providecommand{\TPK}{\mathrm{TPK}}
\providecommand{\TRIT}{\{-1,0,+1\}}

\section{Operator antecedents of the criticality family}
\label{hmtfg:antecedentes-operatorios}

The criticality family is constructed from a readout of the TPK calendar. The two arithmetic sheets, tritic orientation and transport with memory first determine the states and their prolongations. The ordered reading of the continuum then allows analytic functions to be evaluated, parameters to be compared and roots to be selected. This development preserves the genealogical state and its two readouts: the Archimedean chart provides the evaluation space, while the incidence readout preserves boundary relations. The joint discrete structure of the continuum retains its five consubstantial constructions. The present exposition uses its ordered reader and preserves the common antecedent from which it originates.

\subsection{Arithmetic cells, orientation and transport}
\label{hmtfg:app-trit-transporte}

% XI REV11, sections/nucleo.tex: definiciones de APP y TRIT;
% tpk_desarrollo_integrado.tex:22--165, transporte levantado y calendario.
Let \(D_9=\{1,\ldots,9\}\) and \(X_9=(\ZZ/9\ZZ)^2\). For
\(n\in\ZZ\), division with a positive representative is
\[
 \rho_9(n)=1+((n-1)\bmod9),\qquad
 q_9(n)=\frac{n-\rho_9(n)}9,\qquad
 n=\rho_9(n)+9q_9(n).
\]
The APP cell associated with \((p,q)\in D_9^2\) is
\[
 a_{p,q}=(p,q;R^\Sigma,q^\Sigma,R^\Pi,q^\Pi),
 \quad
 p+q=R^\Sigma+9q^\Sigma,\qquad
 pq=R^\Pi+9q^\Pi.
\]
The two pairs \((R^\Sigma,q^\Sigma)\) and \((R^\Pi,q^\Pi)\) retain the additive and multiplicative evaluations, respectively. The division formula reconstructs both evaluations before any subsequent reduction. Selection of an active sheet is recorded in a state coordinate and preserves the complete cell.

The balanced tritic representative \(r_3(n)\in\{-1,0,+1\}\) satisfies
\(n\equiv r_3(n)\pmod3\). Its carry is
\[
 \chi(a,b)=\frac{a+b-r_3(a+b)}3,\qquad
 a,b\in\{-1,0,+1\}.
\]
The equality
\[
 \chi(a,b)+\chi(r_3(a+b),c)
 =\chi(b,c)+\chi(a,r_3(b+c))
\]
is obtained by writing both sides as
\((a+b+c-r_3(a+b+c))/3\). This identity preserves the carry when regrouping a tritic sum. The operative selector is
\[
 M_{\rm ph}(r)=
 \begin{cases}
 +1,&r\in\{1,4,7\},\\
 -1,&r\in\{2,5,8\},\\
 0,&r\in\{3,6,9\}.
 \end{cases}
\]
The value \(+1\) activates the additive sheet, the value \(-1\) the multiplicative sheet, and the value \(0\) preserves the position of both cursors during the neutral event. All three events belong to the complete chronology and are recorded in the route.

Let \(\mathcal D=\{N,E,S,O\}\), with vectors
\(\nu_N=(-1,0)\), \(\nu_E=(0,1)\), \(\nu_S=(1,0)\) and
\(\nu_O=(0,-1)\), and \(\Theta_{108}=\ZZ/108\ZZ\). The observable configuration is
\[
 q=(x_+,d_+,x_\times,d_\times,\theta)
 \in\mathcal Q_{\rm obs}=(X_9\times\mathcal D)^2\times\Theta_{108}.
\]
For the representative \(\bar\theta\in\{0,\ldots,107\}\), define
\[
 r(\theta)=1+(\bar\theta\bmod9),\qquad
 \beta(\theta)=\left[\left\lfloor\frac{\bar\theta}{9}\right\rfloor\right]_{12},
\]
\[
 a_+(\theta)=\mathbf1_{\{1,4,7\}}(r(\theta)),\qquad
 a_\times(\theta)=\mathbf1_{\{2,5,8\}}(r(\theta)).
\]
The position transition is
\[
 x_+'=x_++a_+(\theta)\nu_{d_+},\qquad
 x_\times'=x_\times+a_\times(\theta)\nu_{d_\times}.
\]
The involution \(\jmath\) exchanges \(N,S\) and \(E,O\). With
\[
 h(\theta)=\mathbf1_{\{0,27,54,81\}}\bigl((\bar\theta+1)\bmod108\bigr),
\]
the complete observable operator is fixed by
\[
 F_{\rm obs}(q)=
 (x_+',\jmath^{h(\theta)}d_+,x_\times',
   \jmath^{h(\theta)}d_\times,\theta+1).
\]

\begin{proposition}[Conservation of lifted transport and observable return]
\label{hmtfg:transporte-calendario}
Transport preserves the integer displacement through the record of boundary crossings. The operator \(F_{\rm obs}\) is invertible and satisfies
\(F_{\rm obs}^{108}=I\).
\end{proposition}
\begin{proof}
Let \(s:X_9\to D_9^2\) be the section of positive representatives. For
\(x'=x+a\nu_d\), \(a\in\{0,1\}\), define
\[
 b_d(x,a)=\frac{s(x)+a\nu_d-s(x')}{9}\in\ZZ^2.
\]
If \(z'=z+b_d(x,a)\), then
\[
 s(x')+9z'=s(x)+9z+a\nu_d.
\]
Summing this equality along a route gives
\[
 s(x_n)+9z_n=s(x_0)+9z_0+\sum_{t=0}^{n-1}a_t\nu_{d_t}.
\]
Partitioning the sum preserves the identity under concatenation. To invert \(F_{\rm obs}\), first subtract one unit from the phase; that phase determines the reversal of directions and the displaced cursor. The previous directions are recovered and the corresponding vector is subtracted. Each block of \(27\) steps contains nine activations of each cursor. The next block reverses directions and cancels their integer displacements. In \(54\) steps, positions and directions are restored; in \(108\), the calendar coordinate is also restored.
\end{proof}

In addition to this configuration, the complete state preserves both arithmetic evaluations, carries, orientation, trajectory, prefix, compatible cylinder, its boundary, memory degree and survival record. Its update factors as
\[
 U_t=\mathsf{Mem}_t\circ\mathsf{Tra}_t\circ\mathsf{Sel}_t.
\]
The holonomy \(\Gamma_9\) composes nine complete transitions. Phase return and memory advance are different coordinates of this composition. The following construction fixes the transitions and their update field by field.

% Recuperación íntegra del propietario de 855 líneas, con etiquetas prefijadas.
% Procedencia literal: output/PREPARACION_ENTREGA_USB_20260928/ENTREGA_HMT_USB/03_FUENTES_Y_REPRODUCCION/ACTUALIZACIONES_20260928/FINAL01/PAQUETES/ARTICULO_XI_REV11_ES_EN_FINAL_BN_20260928/ARTICULO_XI_REV11_ES_EN_FINAL_BN/ES/source/sections/ch_tres_vueltas_ext.tex
% Recuperación íntegra; única transformación mecánica: prefijo hmtfg: en label/ref/eqref.
% No se modifican operadores, pruebas, tablas ni afirmaciones del propietario.
\subsubsection{Edge-by-edge realization of the three nonadic extensions}
\label{hmtfg:subsubsec:tres-vueltas-ext-ch}

This subsection specifies the autonomous model generated by the TPK extension
relation entering the cardinal proof. The existence of an edge is not inferred
from a coincidence of types: its initial and final states, APP relation, fiber
transport, and unary update of the ledger are constructed. The construction
uses exclusively the two APP sheets, the TRIT selector, and the TPK
composition \(\mathsf{Sel}\)--\(\mathsf{Tra}\)--\(\operatorname{Upd}\).

\paragraph{Complete cells and three balanced reductions.}
In the chart \(D_9=\{1,\ldots,9\}\), we will always retain the complete APP
cell
\[
 a_{p,q}:=(p,q;R^\Sigma_{p,q},q^\Sigma_{p,q},
                 R^\Pi_{p,q},q^\Pi_{p,q}),
\]
with
\[
 p+q=R^\Sigma_{p,q}+9q^\Sigma_{p,q},\qquad
 pq=R^\Pi_{p,q}+9q^\Pi_{p,q}.
\]
The active sheet is recorded in a separate coordinate; \(a_{p,q}\) is not
split into a purported additive cell and a separate multiplicative cell. For
\(r\in D_9\), let
\begin{equation}
 a(r):=
 \begin{cases}
  a_{1,9},&r=1,\\
  a_{1,r-1},&2\le r\le9,
 \end{cases}
 \qquad
 \tau(r):=M_{\rm ph}(r),
 \qquad
 m(r):=\tau(r)\in\{-1,0,+1\}.
 \label{hmtfg:eq:celula-completa-fase-ch}
\end{equation}
The additive residue of \(a(r)\) supplies an APP representative of phase
here; this does not make \(\Sigma\) the universally active sheet. The dynamic
selector activates \(\Sigma\) for \(r\in\{1,4,7\}\), activates \(\Pi\) for
\(r\in\{2,5,8\}\), and preserves both sheets without moving the cursor for
\(r\in\{3,6,9\}\). In all three cases,
\(R^\Sigma,q^\Sigma,R^\Pi,q^\Pi\) are transported jointly. We have
\[
 (m(1),\ldots,m(9))=(1,-1,0,1,-1,0,1,-1,0).
\]
The three pairs ordered by the nonadic base are
\begin{equation}
 \mathfrak p_{+1}:=(1,4),\qquad
 \mathfrak p_{-1}:=(2,5),\qquad
 \mathfrak p_{0}:=(3,6).
 \label{hmtfg:eq:tres-pares-internos-ch}
\end{equation}
These pairs form the normalized section that takes the first two occurrences of each TRIT type under the shift \(r\mapsto r+3\); the third phase block \(7,8,9\) preserves control of turn closure. Their two members belong to the same fiber of \(M_{\rm ph}\). The balanced reduction defined in these antecedents gives, without introducing a target coefficient,
\begin{equation}
 \begin{array}{c|c|c|c}
 \varepsilon&m(\mathfrak p_\varepsilon^{(1)})
                  +m(\mathfrak p_\varepsilon^{(2)})
   &r_3\bigl(2\varepsilon\bigr)&
   \chi(\varepsilon,\varepsilon)\\ \hline
 +1&1+1=(-1)+3\cdot1&-1&+1\\
 0&0+0=0+3\cdot0&0&0\\
 -1&(-1)+(-1)=1+3\cdot(-1)&+1&-1
 \end{array}
 \label{hmtfg:eq:tres-carry-derivados-ch}
\end{equation}
Therefore,
\begin{equation}
 \chi(\varepsilon,\varepsilon)
 =\frac{2\varepsilon-r_3(2\varepsilon)}3=\varepsilon.
 \label{hmtfg:eq:carry-no-parametrico-ch}
\end{equation}
The branch sign is thus the output of the balanced cocycle applied to two TRIT
emissions produced by APP. Each pair supplies one of the three carry outputs.
For clarity, we subsequently index the pair by
\(\varepsilon(\mathfrak p):=
\chi(m(\mathfrak p^{(1)}),m(\mathfrak p^{(2)}))\); the symbol
\(\varepsilon\) does not supply a free carry to \(\operatorname{Upd}^{\rm cal}\).

The phase architecture, balanced lift, and memory transport are realized here
through twenty-seven local incidences and their extension graphs. The
normalized section fixes representatives of the fibers of \(M_{\rm ph}\); the
covering to be proved uses the existence of these representatives and remains
invariant when they are replaced by another compatible section.

\paragraph{Typed data of the nine edges.}
Let \(K_{\rm Carr}:=\langle\kappa_{\rm Carr}\rangle_{\rm Ab}\) be the carry
module. The actions are typed by
\[
 H_{\rm Carr}:\mathscr P_{\rm TPK}^{\rm op}
   \longrightarrow\operatorname{End}(K_{\rm Carr}),\qquad
 Y_{\rm Carr}:\mathscr P_{\rm TPK}
   \longrightarrow\operatorname{End}(K_{\rm Carr}).
\]
In this generated model, take the trivial actions
\(H_{\rm Carr}(e)=Y_{\rm Carr}(e)=\operatorname{id}_{K_{\rm Carr}}\): the
active sheet and the route are retained as coordinates of the ledger itself
and do not act on the free carry summand. Also let
\((\mathsf M_d,\jmath_{d',d})\) be the graded direct system of memory. To
avoid confusing two representatives identified by the colimit, the state also
retains the degree as a coordinate, and the labeled carrier
\[
 \widetilde{\mathsf M}:=
 \bigsqcup_{d\ge1}\bigl(\{d\}\times\mathsf M_d\bigr),
 \qquad
 \operatorname{gr}_{\mathsf M}(d,\mu)=d.
\]
Publication in the limit memory is given by
\[
\begin{aligned}
 \mathsf M&:=\varinjlim_d\mathsf M_d,&
 \jmath_d&:\mathsf M_d\longrightarrow\mathsf M,\\
 u&:=\jmath_d(u_d),&
 \iota_{\mathsf M}:\mathbb Z&\longrightarrow\mathsf M,
 \qquad n\longmapsto nu .
\end{aligned}
\]
Here \(\jmath_d\) is the canonical colimit map and the \(u_d\) form a
compatible family. Its free component permits
the reader \(\deg_u(nu)=n\).
The configuration, phase, memory, carry, route and
survival projections will be typed on the calibrated carrier constructed
recursively below; no ambient space will be imported for them.
For each constructed edge \(e\), the
calibrated component
\(\operatorname{Upd}^{\rm cal}_e:=\mathsf{Mem}^{\rm cal}_e\) will denote the
third map in the factorization; its law is fixed field by field
below and is not assumed as the restriction of an ambient relation.

Fix the initial depth \(k_0=1\), a complete configuration
\(q_\star\in\mathcal Q_{\rm obs}\) of phase one, and
\[
 q_{n,r}:=F_{\rm obs}^{9n+r-1}(q_\star),\qquad
 q_{n,10}:=q_{n+1,1}.
\]
Because the observable calendar is periodic, its value \(q_{n,r}\) does not
by itself distinguish occurrences separated by twelve turns. The calibrated
model therefore uses the labeled lift
\[
 \widetilde q_{n,r}:=(q_{n,r},n,r),\qquad
 \widetilde q_{n,10}:=\widetilde q_{n+1,1},\qquad
 \pi_{\rm obs}(\widetilde q_{n,r})=q_{n,r}.
\]
The lifted dynamics
\(\widetilde F_{\rm obs}(\widetilde q_{n,r})
=\widetilde q_{n,r+1}\) projects to \(F_{\rm obs}\); the label does not alter the physical phase, but makes its
levels of occurrence distinct. The copy of the cell at each level is labeled by
\[
 a_{n,r}:=(n,r;a(r)),\qquad
 (n,r)^+:=
 \begin{cases}
  (n,r+1),&1\le r<9,\\
  (n+1,1),&r=9.
 \end{cases}
\]
The label only prevents two occurrences of the same cell from being
identified; the two Euclidean divisions and their four coordinates are those
of \(a(r)\). For turn \(n\ge0\) and phase \(r\in D_9\), the symbol
\[
 e_{\varepsilon,n,r}:\widetilde q_{n,r}
 \longrightarrow\widetilde q_{n,r+1},
 \qquad r^+:=1+(r\bmod9),
\]
occupies level \(d_{n,r}:=k_n+r-1\), where \(k_n:=k_0+9n\). Its local ledger
contains the complete cell \(a(r)\), the active sheet determined by
\(M_{\rm ph}(r)\), the type \(\tau(r)\), the membership mark in
\(\mathfrak p_\varepsilon\), and, upon reaching the second member of that
pair, the certified identity \eqref{hmtfg:eq:tres-carry-derivados-ch}. We define its
operators solely from that ledger:
\begin{align}
 \mathsf C_{\mathsf M}(e_{\varepsilon,n,r})
   &=\jmath_{d_{n,r}+1,d_{n,r}},&
 \kappa_{\mathsf M}(e_{\varepsilon,n,r})
   &=\mathbf1_{\{r=9\}}u_{d_{n,r}+1},
 \label{hmtfg:eq:memoria-arista-ch}\\
       \operatorname{Carr}(e_{\varepsilon,n,r})
   &=d_{\varepsilon,r}\kappa_{\rm Carr},&
 d_{\varepsilon,r}
   &:=\mathbf1_{\{r=\mathfrak p_\varepsilon^{(2)}\}}
       \chi\!\left(m(\mathfrak p_\varepsilon^{(1)}),
                    m(\mathfrak p_\varepsilon^{(2)})\right).
 \label{hmtfg:eq:carry-arista-derivado-ch}
\end{align}
For identities and every admissible composition, also define
\begin{align}
 \mathsf C_{\mathsf M}(\operatorname{id}_{\widetilde q};d)
   &:=\operatorname{id}_{\mathsf M_d},&
 \kappa_{\mathsf M}(\operatorname{id}_{\widetilde q};d)&:=0,
 \label{hmtfg:eq:memoria-identidad-ch}\\
 \mathsf C_{\mathsf M}(e_t\cdots e_1)
   &:=\mathsf C_{\mathsf M}(e_t)\circ\cdots\circ
      \mathsf C_{\mathsf M}(e_1),&
 \kappa_{\mathsf M}(e_2e_1)
   &:=\mathsf C_{\mathsf M}(e_2)
        \bigl(\kappa_{\mathsf M}(e_1)\bigr)
      +\kappa_{\mathsf M}(e_2).
 \label{hmtfg:eq:cociclo-memoria-calibrado-ch}
\end{align}
The identity carries degree \(d\) because the same observable configuration
may recur at different depths. Iterated, the final equality defines
\(\kappa_{\mathsf M}(e_t\cdots e_1)\) for every length and is the cocycle law
of the direct system; the inclusions satisfy
\(\jmath_{d+2,d+1}\jmath_{d+1,d}=\jmath_{d+2,d}\).
In addition, set \(s(e_{\varepsilon,n,r})=\top\). The following twenty-seven
entries certify the coordinates distinguishing the three extensions; the
universal transport fields are specified immediately afterward. The symbols
\(I\) and \(II\) designate, respectively, the first retained emission and the
second emission performing the reduction; a dash denotes a zero carry
increment, not the absence of an edge.
\begin{center}
\small
\setlength{\tabcolsep}{3.2pt}
\begin{tabular}{c c cc cc cc c}
\toprule
\(r\)&\(M_{\rm ph}(r)\)&\multicolumn{2}{c}{\(\mathfrak p_{+1}\)}&
\multicolumn{2}{c}{\(\mathfrak p_0\)}&
\multicolumn{2}{c}{\(\mathfrak p_{-1}\)}&
\(\Delta\operatorname{mem}\)\\
& &mark&\(d_{+1,r}\)&mark&\(d_{0,r}\)&mark&\(d_{-1,r}\)&\\
\midrule
1&\(+1\)&I&0&--&0&--&0&0\\
2&\(-1\)&--&0&--&0&I&0&0\\
3&0&--&0&I&0&--&0&0\\
4&\(+1\)&II&\(+1\)&--&0&--&0&0\\
5&\(-1\)&--&0&--&0&II&\(-1\)&0\\
6&0&--&0&II&0&--&0&0\\
7&\(+1\)&--&0&--&0&--&0&0\\
8&\(-1\)&--&0&--&0&--&0&0\\
9&0&--&0&--&0&--&0&\(u\)\\
\bottomrule
\end{tabular}
\end{center}
Thus, \(d_{\varepsilon,r}\) is computed from two outputs of the selector and
is distinct from the cell quotients \(q^\Sigma,q^\Pi\). The ledger also
preserves the marked pair and the balanced residue; in particular, the zero
branch is not confused with an absence of event.

The carry composition law used by the TPK category is
\begin{equation}
 \operatorname{Carr}(\operatorname{id}_{\widetilde q})=0,\qquad
 \operatorname{Carr}(e_2e_1)
 =H_{\rm Carr}(e_1)\bigl(\operatorname{Carr}(e_2)\bigr)
  +Y_{\rm Carr}(e_2)\bigl(\operatorname{Carr}(e_1)\bigr).
 \label{hmtfg:eq:ley-carry-correcta-ch}
\end{equation}
On identities and paths, they are extended functorially:
\[
\begin{aligned}
 H_{\rm Carr}(\operatorname{id})
   &=Y_{\rm Carr}(\operatorname{id})=\operatorname{id},\\
 H_{\rm Carr}(e_t\cdots e_1)
   &=H_{\rm Carr}(e_1)\circ\cdots\circ H_{\rm Carr}(e_t),\\
 Y_{\rm Carr}(e_t\cdots e_1)
   &=Y_{\rm Carr}(e_t)\circ\cdots\circ Y_{\rm Carr}(e_1).
\end{aligned}
\]
In this model all these actions remain the identity, but
their typing makes explicit which endomorphism acts on each summand of
\eqref{hmtfg:eq:ley-carry-correcta-ch}.
On calibrated edges, \(H_{\rm Carr}=Y_{\rm Carr}=\operatorname{id}\),
so composition sums the increments already produced. Equation
\eqref{hmtfg:eq:ley-carry-correcta-ch}, rather than an identification between orientation and
quotient, governs the accumulated register.

\paragraph{Independent construction of the transition tuples.}
A fiber state is written, in the order of its relevant fields, as
\[
 x=(a;q;\tau,o,h;b_{\rm ref};
       R^\Sigma,q^\Sigma,R^\Pi,q^\Pi;c;
       \operatorname{Sig};C,\partial C;p_{\APP},p_{\TRIT};
       \lambda;\operatorname{gr}_{\mathsf M};\operatorname{mem};
       s;\operatorname{Led}).
\]
The update notation \(x[f\leftarrow v]\) changes only the indicated field.
For \(h\in\{(\Sigma,\Pi),(\Pi,\Sigma)\}\), define
\[
 h_r(h)=
 \begin{cases}
  (\Sigma,\Pi),&r\in\{1,4,7\},\\
  (\Pi,\Sigma),&r\in\{2,5,8\},\\
 h,&r\in\{3,6,9\}.
 \end{cases}
\]
The complete sheet-and-arithmetic coordinate is abbreviated as
\[
 \widehat h(x):=
 \bigl(h(x),R^\Sigma(x),q^\Sigma(x),R^\Pi(x),q^\Pi(x)\bigr).
\]
Let \(\widehat{\mathcal H}^{\rm cal}_{\widetilde q_{n,r}}\) be the set of
quintuples
\[
 \bigl(h,R^\Sigma_{a_{n,r}},q^\Sigma_{a_{n,r}},
          R^\Pi_{a_{n,r}},q^\Pi_{a_{n,r}}\bigr),
 \qquad
 h\in\{(\Sigma,\Pi),(\Pi,\Sigma)\},
\]
whose four arithmetic data are the two Euclidean divisions of the labeled APP
cell \(a_{n,r}\). This direct definition does not presuppose a state in a
carrier not yet constructed. Let
\(\mathsf{Hist}^{\rm cal}_{\widetilde q}\) be the set of admissible finite
paths in the graph of labeled edges ending at \(\widetilde q\), including the
identity, and let \(\mathsf B^{\rm cal}_{\widetilde q}\) be the set of their
ordered endpoint sequences. Define
\[
 \mathsf S^{\rm cal}_{\widetilde q}:=
 \widehat{\mathcal H}^{\rm cal}_{\widetilde q}
 \times\{o_{\rightarrow},o_{\circ},o_{\leftarrow}\}
 \times K_{\rm Carr}\times
 \mathsf{Hist}^{\rm cal}_{\widetilde q}\times
 \mathsf B^{\rm cal}_{\widetilde q}\times\mathsf M .
\]
The calibrated signature is not postulated: it is the typed assembly map
\[\begin{gathered}
 \operatorname{Sig}^{\rm cal}_{\widetilde q}:
 \widehat{\mathcal H}^{\rm cal}_{\widetilde q}
 \times\{o_{\rightarrow},o_{\circ},o_{\leftarrow}\}
 \times K_{\rm Carr}\times
 \mathsf{Hist}^{\rm cal}_{\widetilde q}\times
 \mathsf B^{\rm cal}_{\widetilde q}\times\mathsf M
 \longrightarrow\mathsf S^{\rm cal}_{\widetilde q},
\\
 (\widehat h,o,c,\lambda,\partial C,\mu)
 \longmapsto(\widehat h,o,c,\lambda,\partial C,\mu).
\end{gathered}\]
Likewise, \(\pi_{\mathcal F}\) denotes the projection that removes only the
first two coordinates \((a;q)\) of the state tuple; it preserves
\(\tau,o,h,b_{\rm ref}\), the four Euclidean data, carry, signature, cylinder,
boundary, prefixes, route, memory degree, memory, survival, and ledger. The
corresponding orientation is fixed by
\(o(+1)=o_{\rightarrow}\), \(o(0)=o_{\circ}\), and
\(o(-1)=o_{\leftarrow}\).
The calibrated APP fiber over \(\widetilde q_{n,r}\) contains the already
defined labeled copy \(a_{n,r}\). The initial state is fixed without free
coordinates:
\[
 \widehat h_\star:=
 \bigl((\Sigma,\Pi),R^\Sigma_{1,9},q^\Sigma_{1,9},
       R^\Pi_{1,9},q^\Pi_{1,9}\bigr).
\]
\begin{equation}
\begin{aligned}
 x_\star:=\bigl(&a_{0,1};\widetilde q_{0,1};
 +1,o_{\rightarrow},(\Sigma,\Pi);\bot;
 R^\Sigma_{1,9},q^\Sigma_{1,9},R^\Pi_{1,9},q^\Pi_{1,9};0;\\
 &\operatorname{Sig}^{\rm cal}_{\widetilde q_{0,1}}
       (\widehat h_\star,o_{\rightarrow},0,
        \operatorname{id}_{\widetilde q_{0,1}},\varnothing,
        \jmath_{k_0}(0_{\mathsf M_{k_0}}));
 \varnothing,\varnothing;
 (a_{0,1}),(+1);\operatorname{id}_{\widetilde q_{0,1}};
 k_0;0_{\mathsf M_{k_0}};\top;\varnothing\bigr).
\end{aligned}
\label{hmtfg:eq:estado-inicial-calibrado-ch}
\end{equation}
The tuple contains no indeterminate coordinates and satisfies
\(s(x_\star)=\top\). Let \(G_0:=\{x_\star\}\). Assuming \(G_n\) has been
constructed, fix \(x\in G_n\), \(\varepsilon\in\{-1,0,+1\}\), and
\(x_{\varepsilon,0}(x):=x\). For \(r=1,\ldots,9\), first construct the
elementary ledger
\begin{equation}
\begin{aligned}
 \ell_{\varepsilon,n,r}:=\bigl(&e_{\varepsilon,n,r},a_{n,r},a_{(n,r)^+},
 \widetilde q_{n,r},\widetilde q_{n,r+1},M_{\rm ph}(r),\\
 &h_r(h(x_{\varepsilon,r-1}(x))),o(M_{\rm ph}(r)),\mathfrak p_\varepsilon,
 d_{\varepsilon,r},\\
 &\mathsf C_{\mathsf M}(e_{\varepsilon,n,r}),
 \kappa_{\mathsf M}(e_{\varepsilon,n,r}),
 \operatorname{Carr}(e_{\varepsilon,n,r})\bigr),
\end{aligned}
 \label{hmtfg:eq:registro-elemental-ext-ch}
\end{equation}
and set
\(\operatorname{Led}^{\rm new}_{\varepsilon,n,r}:=
\operatorname{Led}(x_{\varepsilon,r-1}(x))^\frown
\ell_{\varepsilon,n,r}\). Then construct the tuple
\begin{equation}
 y_{\varepsilon,r}(x):=
 x_{\varepsilon,r-1}(x)
 [\tau\leftarrow M_{\rm ph}(r),\
  h\leftarrow h_r(h(x_{\varepsilon,r-1}(x))),\
  o\leftarrow o(M_{\rm ph}(r)),\
  b_{\rm ref}\leftarrow\mathfrak p_\varepsilon],
 \label{hmtfg:eq:tupla-sel-cal-ch}
\end{equation}
where selection leaves the cell, quotients, carry, route, memory, cylinder,
boundary, survival, and ledger unchanged; the coordinate \(b_{\rm ref}\)
makes the chosen balanced reduction explicit. Membership in an ambient
relation is not presupposed here.

Next, define \(z_{\varepsilon,r}(x)\) field by field: its configuration is
\(\widetilde q_{n,r+1}\); its complete cell is \(a_{(n,r)^+}\), with the four
residues and quotients recomputed by the two Euclidean divisions; it preserves
the selected sheet and orientation; and it performs the following six updates,
in which only the common argument \(x\) is omitted:
\begin{equation}
\begin{aligned}
 C(z_{\varepsilon,r})&:=C(y_{\varepsilon,r})^\frown
       e_{\varepsilon,n,r},\\
 \partial C(z_{\varepsilon,r})&:=\partial C(y_{\varepsilon,r})^\frown
       (\widetilde q_{n,r},\widetilde q_{n,r+1}),\\
 p_{\APP}(z_{\varepsilon,r})&:=p_{\APP}(y_{\varepsilon,r})^\frown
       a_{(n,r)^+},\\
 p_{\TRIT}(z_{\varepsilon,r})&:=p_{\TRIT}(y_{\varepsilon,r})^\frown
       \tau(r),\\
 \operatorname{Led}(z_{\varepsilon,r})&:=
       \operatorname{Led}^{\rm new}_{\varepsilon,n,r},\\
 \lambda(z_{\varepsilon,r})&:=e_{\varepsilon,n,r}
       \lambda(y_{\varepsilon,r}),
\end{aligned}
 \label{hmtfg:eq:tupla-tra-cal-ch}
\end{equation}
and still retains the fields
\(c,\operatorname{gr}_{\mathsf M},\operatorname{mem},s\) of
\(y_{\varepsilon,r}(x)\). Its signature, by contrast, is retyped at the target
configuration through the transported presignature
\begin{equation}
\begin{aligned}
 \operatorname{Sig}(z_{\varepsilon,r}(x))
  :=\operatorname{Sig}^{\rm cal}_{\widetilde q_{n,r+1}}
       \Bigl(&\widehat h(z_{\varepsilon,r}(x)),
             o(z_{\varepsilon,r}(x)),c(y_{\varepsilon,r}(x)),
             \lambda(z_{\varepsilon,r}(x)),\\[-2pt]
             &\partial C(z_{\varepsilon,r}(x)),
             \jmath_{d_{n,r}}
                (\operatorname{mem}_{d_{n,r}}
                   (y_{\varepsilon,r}(x)))\Bigr).
\end{aligned}
\label{hmtfg:eq:prefirma-transporte-cal-ch}
\end{equation}
Thus, and only in \(\mathsf{Tra}^{\rm cal}\), the edge, its endpoints, the
mark \((\mathfrak p_\varepsilon,r,d_{\varepsilon,r})\), and its cocycles are
added to the incidence ledger.

The initial state satisfies
\[
 c(x_\star)=0=\operatorname{Carr}
   (\operatorname{id}_{\widetilde q_{0,1}})
   =\operatorname{Carr}(\lambda(x_\star)).
\]
In the complete states \(x_{\varepsilon,r-1}(x)\), and also after the
selection \(y_{\varepsilon,r}(x)\), the typed invariant
\(c(v)=\operatorname{Carr}(\lambda(v))\) is preserved, since
\(\mathsf{Sel}^{\rm cal}\) modifies neither field. Transport advances the
route but leaves the carry pending; hence the pre-update tuple satisfies
exactly
\[
 c(z_{\varepsilon,r}(x))=c(y_{\varepsilon,r}(x))
 =\operatorname{Carr}(\lambda(y_{\varepsilon,r}(x))).
\]
The map \(\operatorname{Upd}^{\rm cal}\) then restores
\(c(x_{\varepsilon,r}(x))=
\operatorname{Carr}(\lambda(x_{\varepsilon,r}(x)))\) by means of the
composition law, without attributing that invariant to the intermediate state \(z\).

Finally, apply the unary function
\(\operatorname{Upd}^{\rm cal}_{e_{\varepsilon,n,r}}
=\mathsf{Mem}^{\rm cal}_{e_{\varepsilon,n,r}}\) to that tuple. Its result
\(x_{\varepsilon,r}(x)\) is given by
\begin{equation}
\begin{aligned}
 c\bigl(x_{\varepsilon,r}(x)\bigr)
   &=\operatorname{Carr}\bigl(\lambda(z_{\varepsilon,r}(x))\bigr)\\
   &=H_{\rm Carr}(\lambda(y_{\varepsilon,r}(x)))
       \bigl(\operatorname{Carr}(e_{\varepsilon,n,r})\bigr)
     +Y_{\rm Carr}(e_{\varepsilon,n,r})
       \bigl(c(y_{\varepsilon,r}(x))\bigr),\\
 \lambda\bigl(x_{\varepsilon,r}(x)\bigr)
   &=\lambda(z_{\varepsilon,r}(x)),\\
 \operatorname{gr}_{\mathsf M}
      \bigl(x_{\varepsilon,r}(x)\bigr)
   &=d_{n,r}+1,\\
 \operatorname{mem}_{d_{n,r}+1}
       \bigl(x_{\varepsilon,r}(x)\bigr)
   &=\mathsf C_{\mathsf M}(e_{\varepsilon,n,r})
       \operatorname{mem}_{d_{n,r}}(z_{\varepsilon,r}(x))
     +\kappa_{\mathsf M}(e_{\varepsilon,n,r}),\\
 s\bigl(x_{\varepsilon,r}(x)\bigr)
   &=s(e_{\varepsilon,n,r})\wedge s(z_{\varepsilon,r}(x)),\\
 \operatorname{Sig}\bigl(x_{\varepsilon,r}(x)\bigr)
   &=\operatorname{Sig}^{\rm cal}_{\widetilde q_{n,r+1}}
       \Bigl(\widehat h(z_{\varepsilon,r}(x)),
             o(z_{\varepsilon,r}(x)),
             c(x_{\varepsilon,r}(x)),
             \lambda(z_{\varepsilon,r}(x)),
             \partial C(z_{\varepsilon,r}(x)),\\[-2pt]
   &\hspace{112pt}
             \jmath_{d_{n,r}+1}
              \bigl(\operatorname{mem}_{d_{n,r}+1}
                 (x_{\varepsilon,r}(x))\bigr)\Bigr).
\end{aligned}
\label{hmtfg:eq:upd-unaria-tres-vueltas-ch}
\end{equation}
All fields not enumerated in \eqref{hmtfg:eq:upd-unaria-tres-vueltas-ch} coincide
with those of \(z_{\varepsilon,r}(x)\). In particular, the signature is
computed from the sheet, orientation, boundary, and route already fixed by
\(\mathsf{Sel}^{\rm cal}\) and \(\mathsf{Tra}^{\rm cal}\), and from the carry
and memory computed in the preceding lines; it is not circularly defined from
the final state. Neither \(\varepsilon\), nor \(d_{\varepsilon,r}\), nor
\(u\) is an additional argument of \(\operatorname{Upd}^{\rm cal}\): the
function reads all three data from the already constructed edge.

\paragraph{Simultaneous induction of the levels.}
Once the preceding nine stages have been completed for an already constructed
\(G_n\), define immediately
\begin{equation}
 \gamma_{\varepsilon,k_n}(x):=x_{\varepsilon,9}(x),\qquad
 G_{n+1}:=\{\gamma_{\varepsilon,k_n}(x):
       x\in G_n,\ \varepsilon\in\TRIT\}.
 \label{hmtfg:eq:injerto-total-tres-vueltas-ch}
\end{equation}
\label{hmtfg:eq:tres-vueltas-tpk-ch}
The base \(G_0=\{x_\star\}\), the intermediate tuples, and
\eqref{hmtfg:eq:injerto-total-tres-vueltas-ch} constitute a single definition by
recursion on \(n\). Therefore, no subsequent carrier is used to assume the
existence of \(G_n\): the next level is introduced only after constructing all
its edges from the preceding level.

\paragraph{TPK extension relation and effective membership.}
Let \(\mathcal E^{\rm cal}_{\widetilde q}\) be the smallest labelled family containing
\(x_\star\) and, at each step of the preceding recurrence, the four tuples
\(x_{\varepsilon,r-1}(x),y_{\varepsilon,r}(x),
z_{\varepsilon,r}(x),x_{\varepsilon,r}(x)\) in their source
and target configurations. Also let
\(\mathcal A^{\rm cal}_{\widetilde q_{n,r}}:=\{a_{n,r}\}\), and set
\[
 \widetilde{\mathcal Q}_{\rm obs}:=
   \{\widetilde q_{n,r}:n\ge0,\ 1\le r\le9\},
 \qquad
 \mathcal E^{\rm cal}:=\bigsqcup_{\widetilde q}
     \mathcal E^{\rm cal}_{\widetilde q},
 \qquad
 \mathcal F^{\rm cal}_{\widetilde q}:=
     \pi_{\mathcal F}(\mathcal E^{\rm cal}_{\widetilde q}).
\]
For each depth \(d\), let
\(\mathcal E^{\rm cal}_d:=
\operatorname{gr}_{\mathsf M}^{-1}(d)\subset\mathcal E^{\rm cal}\). The
explicit degree coordinate, rather than the mere membership of a
representative in some image of the direct system, makes these carriers
disjoint. On them, the projections are typed as
\[
\begin{aligned}
 \operatorname{cfg}&:\mathcal E^{\rm cal}\longrightarrow
     \widetilde{\mathcal Q}_{\rm obs},\\
 g(x)&:=\operatorname{fase}\!
       \bigl(\pi_{\rm obs}(\operatorname{cfg}(x))\bigr)\in C_9,\\
 c&:\mathcal E^{\rm cal}\longrightarrow K_{\rm Carr},\\
 s&:\mathcal E^{\rm cal}\longrightarrow\{\bot,\top\},\\
 \lambda&:\mathcal E^{\rm cal}_{\widetilde q}
       \longrightarrow\mathsf{Hist}^{\rm cal}_{\widetilde q},\\
 \operatorname{gr}_{\mathsf M}&:\mathcal E^{\rm cal}
       \longrightarrow\mathbb N_{\ge1},\\
 \operatorname{mem}_d&:\mathcal E^{\rm cal}_d
       \longrightarrow\mathsf M_d,\\
 \operatorname{mem}(x)&:=
       \jmath_d\bigl(\operatorname{mem}_d(x)\bigr)\in\mathsf M
       \quad(x\text{ of depth }d).
\end{aligned}
\]
Here \(\operatorname{cfg}(x)\) is the second coordinate of the tuple, and
\(\operatorname{fase}\) reads the phase of the observable calendar. For the
three carrier classes, fix the diagonals
\[
\begin{aligned}
 \Delta_{\APP,\widetilde q}
   &:=\{(a,a):a\in\mathcal A^{\rm cal}_{\widetilde q}\},\\
 \Delta_{{\rm fib},\widetilde q}
   &:=\{(f,f):f\in\mathcal F^{\rm cal}_{\widetilde q}\},\\
 \Delta_{{\rm enr},\widetilde q}
   &:=\{(x,x):x\in\mathcal E^{\rm cal}_{\widetilde q}\}.
\end{aligned}
\]
In particular,
\[
 a_{0,1}\in\mathcal A^{\rm cal}_{\widetilde q_{0,1}},
 \qquad
 x_\star\in\mathcal E^{\rm cal}_{\widetilde q_{0,1}}.
\]
On these carriers, the following relations are defined by their complete
graphs, and not through necessary implications:
\begin{align}
 \mathsf{Sel}^{\rm cal}_{e_{\varepsilon,n,r}}
   &:=\{(x_{\varepsilon,r-1}(x),y_{\varepsilon,r}(x)):
          x\in G_n\},\\
 \mathsf{Tra}^{\rm cal}_{e_{\varepsilon,n,r}}
   &:=\{(y_{\varepsilon,r}(x),z_{\varepsilon,r}(x)):
          x\in G_n\},\\
 \operatorname{Ext}^{\APP,\rm cal}_{e_{\varepsilon,n,r}}
   &:=\{(a_{n,r},a_{(n,r)^+})\},
 \label{hmtfg:eq:ext-app-calibrada-ch}\\
 \operatorname{Ext}^{\rm fib,cal}_{e_{\varepsilon,n,r}}
   &:=\{(\pi_{\mathcal F}x_{\varepsilon,r-1}(x),
          \pi_{\mathcal F}x_{\varepsilon,r}(x)):x\in G_n\}.
 \label{hmtfg:eq:ext-fibra-calibrada-ch}
\end{align}
The synchronized enriched extension is
\begin{equation}
 \operatorname{Ext}^{\rm enr,cal}_{e_{\varepsilon,n,r}}
 :=\{(x_{\varepsilon,r-1}(x),x_{\varepsilon,r}(x)):
       x\in G_n\}.
 \label{hmtfg:eq:ext-enriquecida-calibrada-ch}
\end{equation}
The third map is defined by
\begin{equation}
 \operatorname{graph}
 \bigl(\operatorname{Upd}^{\rm cal}_{e_{\varepsilon,n,r}}\bigr)
 :=\{(z_{\varepsilon,r}(x),x_{\varepsilon,r}(x)):
       x\in G_n\},
 \label{hmtfg:eq:grafo-upd-calibrado-ch}
\end{equation}
where the second component is exactly the one in
\eqref{hmtfg:eq:upd-unaria-tres-vueltas-ch}. Consequently,
\begin{equation}
 \mathcal U^{(1),\rm cal}_{e_{\varepsilon,n,r}}
 :=\operatorname{Upd}^{\rm cal}_{e_{\varepsilon,n,r}}\circ
   \mathsf{Tra}^{\rm cal}_{e_{\varepsilon,n,r}}\circ
   \mathsf{Sel}^{\rm cal}_{e_{\varepsilon,n,r}}
 =\{(x_{\varepsilon,r-1}(x),x_{\varepsilon,r}(x)):
      x\in G_n\}.
 \label{hmtfg:eq:transicion-cal-en-tpk-ch}
\end{equation}
Each operator is indexed by the edge \(e_{\varepsilon,n,r}\) produced by the
reduction \(\mathfrak p_\varepsilon\), rather than by a trit supplied from
outside. The coordinate \(b_{\rm ref}\) also distinguishes their codomains.
Therefore,
\(\mathsf{Sel}^{\rm cal}_{e}\),
\(\mathsf{Tra}^{\rm cal}\) and \(\operatorname{Upd}^{\rm cal}\) are functions
on their respective domains, whereas the union
\(\bigcup_{\varepsilon\in\TRIT}
\mathcal U^{(1),\rm cal}_{e_{\varepsilon,n,r}}\) deliberately preserves the
three outputs of the trisection.

For \(\bullet\in\{\APP,{\rm fib},{\rm enr}\}\), identities and compositions
are not left implicit. Define
\begin{align}
 \operatorname{Ext}^{\bullet,\rm cal}_{\operatorname{id}_{\widetilde q}}
   &:=\Delta_{\bullet,\widetilde q},\\
 \operatorname{Ext}^{\bullet,\rm cal}_{e_t\cdots e_1}
   &:=\operatorname{Ext}^{\bullet,\rm cal}_{e_t}\circ\cdots\circ
      \operatorname{Ext}^{\bullet,\rm cal}_{e_1}
 \label{hmtfg:eq:ext-palabra-calibrada-ch}
\end{align}
for every admissible edge word. The two APP and fiber path closures are formed
separately,
\begin{equation}
 \operatorname{Path}(\operatorname{Ext}^{\APP,\rm cal}),\qquad
 \operatorname{Path}(\operatorname{Ext}^{\rm fib,cal}),
 \label{hmtfg:eq:dos-clausuras-ext-ch}
\end{equation}
and \(\mathsf{Tra}^{\rm cal}\) synchronizes them edge by edge.

Finally, the levels of complete states and their truncations are fixed by the
recursion itself:
\begin{align}
 H^{\rm cal}_{d_{n,r}}
   &:=\{x_{\varepsilon,r-1}(x):
          x\in G_n,\ \varepsilon\in\TRIT\},\\
 H^{\rm cal}_{d_{n,r}+1}
   &:=\{x_{\varepsilon,r}(x):
          x\in G_n,\ \varepsilon\in\TRIT\},\\
 \rho^{\rm cal}_{d_{n,r}+1,d_{n,r}}
   \bigl(x_{\varepsilon,r}(x)\bigr)
   &:=x_{\varepsilon,r-1}(x).
 \label{hmtfg:eq:truncamiento-calibrado-definido-ch}
\end{align}
The final ledger preserves \((n,r,\mathfrak p_\varepsilon)\), so the
predecessor of the left-hand side is unique and \(\rho^{\rm cal}\) is
well-defined. For \(b>a\), set
\(\rho^{\rm cal}_{b,a}:=\rho^{\rm cal}_{a+1,a}\circ\cdots\circ
\rho^{\rm cal}_{b,b-1}\). This map recovers the complete preceding state,
including its signature; it does not purport to invert an overwritten
coordinate without consulting the ledger.

\begin{proposition}[Autonomous realization of the TPK through explicit transitions]
\label{hmtfg:prop:modelo-tpk-calibrado-ch}
The family
\[
\begin{aligned}
 \mathfrak T_{\rm TPK}^{\rm cal}:=\bigl(
 &(\widetilde q_{n,r})_{n,r},\pi_{\rm obs},
   \widetilde F_{\rm obs},M_{\rm ph},\\
 &(\mathcal A_{\widetilde q}^{\rm cal})_{\widetilde q},
   (\mathcal E_{\widetilde q}^{\rm cal})_{\widetilde q},
   (\mathcal F_{\widetilde q}^{\rm cal})_{\widetilde q},
   \operatorname{Ext}^{\APP,\rm cal},
   \operatorname{Ext}^{\rm fib,cal},
   \operatorname{Ext}^{\rm enr,cal},\rho^{\rm cal},\\
 &\mathsf{Sel}^{\rm cal},\mathsf{Tra}^{\rm cal},
   \operatorname{Upd}^{\rm cal},\mathsf C_{\mathsf M},
   \kappa_{\mathsf M},\operatorname{gr}_{\mathsf M},
   H_{\rm Carr},Y_{\rm Carr},\\
 &\operatorname{Sig}^{\rm cal},\pi_{\mathcal F}\bigr).
\end{aligned}
\]
is an enriched model of the TPK. In particular, every pair on the right-hand
side of \eqref{hmtfg:eq:transicion-cal-en-tpk-ch} is an effective TPK edge, not a
consequence inferred from a merely necessary condition.
\end{proposition}

\begin{proof}
The labels \((n,r)\) make the carrier copies disjoint and unambiguously fix
each source and target. The relation \(\operatorname{Ext}^{\APP,\rm cal}\)
transports the complete cell and preserves the two Euclidean identities; the
relation \(\operatorname{Ext}^{\rm fib,cal}\) relates the fibers of the
preceding and subsequent complete states: it transports sheet, orientation,
prefix, cylinder, boundary, and ledger, and incorporates the updated carry,
memory, and signature. Selection does not modify the ledger, and transport
concatenates exactly one entry and produces the target presignature
\eqref{hmtfg:eq:prefirma-transporte-cal-ch}. The update recomputes the signature by
means of the typed map \(\operatorname{Sig}^{\rm cal}_{\widetilde q}\) from
the already determined sheet, orientation, carry, route, boundary, and memory.
The map \(\operatorname{Upd}^{\rm cal}\) is unary because the edge contained
in \(z_{\varepsilon,r}(x)\) determines its cocycles.

The identities are the diagonals, and the word extensions are the compositions
in \eqref{hmtfg:eq:ext-palabra-calibrada-ch}. Concatenation of prefixes and ledgers
is associative; memory satisfies the cocycle law, and carry satisfies
\eqref{hmtfg:eq:ley-carry-correcta-ch}. Therefore, the two parenthesizations of three
edges produce the same state. Formula
\eqref{hmtfg:eq:truncamiento-calibrado-definido-ch} defines truncation on complete
states, and uniqueness of the final ledger proves that it recovers the source
in every field. These are the identity, source, target, composition, memory,
carry, and truncation laws of the enriched TPK. In particular, the extension
relations are determined as graphs of effective maps rather than as
consequences of merely necessary conditions.
\end{proof}

Each \(\gamma_{\varepsilon,k_n}\) comprises exactly nine applications
of \(\operatorname{Upd}^{\rm cal}_{e_{\varepsilon,n,r}}\circ
\mathsf{Tra}^{\rm cal}_{e_{\varepsilon,n,r}}\circ
\mathsf{Sel}^{\rm cal}_{e_{\varepsilon,n,r}}\). Survival will be proved by
means of a constructed tail and will not be used to define \(G_n\).

\begin{lemma}[Effective internal trisection of one turn]
\label{hmtfg:lem:elevacion-catalogal-nonadica}
For every \(n\ge0\), every \(x\in G_n\), and every
\(\varepsilon\in\{-1,0,+1\}\), the expression
\(\gamma_{\varepsilon,k_n}(x)\) is defined by nine TPK transitions, with its
synchronized relations
\(\operatorname{Ext}^{\APP,\rm cal}\) and
\(\operatorname{Ext}^{\rm fib,cal}\), and satisfies
\begin{align}
 \rho^{\rm cal}_{k_n+9,k_n}\gamma_{\varepsilon,k_n}(x)&=x,
 \label{hmtfg:eq:truncamiento-injerto-ch}\\
 g\bigl(\gamma_{\varepsilon,k_n}(x)\bigr)
   &=g(x),&
 \jmath_{k_n+9}\!\left(\operatorname{mem}_{k_n+9}
        (\gamma_{\varepsilon,k_n}(x))\right)
   &=\jmath_{k_n}\!\left(\operatorname{mem}_{k_n}(x)\right)+u,
 \label{hmtfg:eq:fase-memoria-injerto-ch}\\
 c\bigl(\gamma_{\varepsilon,k_n}(x)\bigr)
   &=c(x)+\varepsilon\kappa_{\rm Carr}.&&
 \label{hmtfg:eq:carry-injerto-ch}
\end{align}
\label{hmtfg:eq:tres-vueltas-propiedades-ch}
\label{hmtfg:eq:tres-vueltas-acarreo-ch}
The ordered ledger of the nine edges determines \(\varepsilon\) uniquely.
The three images are therefore disjoint. Each endpoint belongs to
\(G_{n+1}\) and again admits the three maps
\(\gamma_{-1,k_{n+1}},\gamma_{0,k_{n+1}},
\gamma_{+1,k_{n+1}}\).
\end{lemma}

\begin{proof}
Formulas \eqref{hmtfg:eq:celula-completa-fase-ch} traverse the nine phases once and
return the copy \(a_{n+1,1}\) of \(a(1)\). At each step, the APP relation
transports the complete sextuple; \(\mathsf{Sel}^{\rm cal}\) fixes the TRIT
type from the phase mark \(r\) preserved in the state; and
\(\mathsf{Tra}^{\rm cal}\) adds exactly one symbol to the prefix, one segment
to the cylinder, and its ordered pair of endpoints to the boundary. By
induction on \(r\), the explicit tuples \(y_{\varepsilon,r}(x)\) and
\(z_{\varepsilon,r}(x)\) satisfy the predicates of the indicated relations,
and the unary update produces the state at the next level. The identity
\(\rho^{\rm cal}_{d_{n,r}+1,d_{n,r}}
(x_{\varepsilon,r}(x))=x_{\varepsilon,r-1}(x)\) is the definition
\eqref{hmtfg:eq:truncamiento-calibrado-definido-ch}; its well-definedness follows
from uniqueness of the final ledger. The composition of the nine truncation
maps recovers \(x\), proving \eqref{hmtfg:eq:truncamiento-injerto-ch} without
separately inverting the overwritten fields.

The phase \(9\to1\) occurs exactly once. By
\eqref{hmtfg:eq:memoria-arista-ch} and the cocycle law, the linear memory component
transports the preceding memory, and its translational component adds exactly
\(u\) in the direct limit. The compatibility
\(\jmath_{d+1}\jmath_{d+1,d}=\jmath_d\) proves
\eqref{hmtfg:eq:fase-memoria-injerto-ch}; the return concerns the phase, not the
enriched state.

By \eqref{hmtfg:eq:carry-arista-derivado-ch}, eight edges have zero increment, and
the edge completing \(\mathfrak p_\varepsilon\) has the computed increment
\(\chi(\varepsilon,\varepsilon)\kappa_{\rm Carr}\). The laws
\eqref{hmtfg:eq:ley-carry-correcta-ch} and
\eqref{hmtfg:eq:carry-no-parametrico-ch}
give \eqref{hmtfg:eq:carry-injerto-ch}. The ledger preserves
\((\mathfrak p_\varepsilon,2\varepsilon,r_3(2\varepsilon),
\chi(\varepsilon,\varepsilon))\); the three values of this ledger are
distinct even when the scalar increment is zero. Hence the images are not
identified, and the block decoder recovers a unique \(\varepsilon\).

The edge definitions depend only on the new phase and concatenate, rather than
replace, the prefix, cylinder, boundary, and ledger. They are also invariant
under
\(\jmath_d\operatorname{mem}_d\mapsto
  \jmath_d\operatorname{mem}_d+nu\)
and under translation of the input carry. Consequently, the same construction
applies at the endpoint of any turn. For every finite state, the explicit
sequence
\[
 x,\quad\gamma_{0,k_n}(x),\quad
 \gamma_{0,k_{n+1}}\gamma_{0,k_n}(x),\quad\ldots
\]
is a compatible infinite tail. This proves the survival of all constructed
states and the repeated availability of the three branches, without
postulating either property. Consequently, in accordance with the preceding
level definitions, we may finally identify
\begin{equation}
 H_{k_n}^{\rm cal}:=G_n
 \label{hmtfg:eq:subarbol-calibrado-superviviente-ch}
\end{equation}
for every \(n\ge0\).
\end{proof}

The preceding closure distinguishes the three clocks involved. Nine emissions
return the phase in \(C_9\) and advance memory. Twelve turns, that is,
\(108\) emissions, also return the position of the observable calendar; even
then, neither the ledger nor memory is erased. The operator
\(\mathcal K_{\rm ph}\) acts only afterward, as a recognition of reduced
memory from the already constructed turns; it never generates the edges of
\(\operatorname{Ext}^{\rm enr,cal}\).


\subsection{Cylinder refinement and ordered readout}
\label{hmtfg:refinamiento-ordenado}

% XI REV11, ch_elevacion_catalogal.tex: bloques de 54 aristas, beta_blk,
% partición efectiva de 729 subcilindros; ch_cierre_cardinal.tex:375--497.
For \(\mathbf b=(b_1,\ldots,b_6)\in\FF_3^6\), let
\(\upsilon:\FF_3\to\{-1,0,+1\}\) be the balanced representative. The three prolongations constructed edge by edge define
\[
 \Gamma^{[54]}_{\mathbf b,k}
 =\gamma_{\upsilon(b_6),k+45}\circ\cdots\circ
   \gamma_{\upsilon(b_1),k}.
\]
Starting from \(Y_0=\{x_\star\}\), put
\[
 Y_{N+1}
 =\{\Gamma^{[54]}_{\mathbf b,1+54N}(x):
          x\in Y_N,\ \mathbf b\in\FF_3^6\}.
\]
Truncation \(\rho^{\rm blk}_{N+1,N}\) removes the last fifty-four edges, recovering the previous state through the transition record.

\begin{proposition}[Effective coding of prolongation blocks]
\label{hmtfg:bloques-729}
Reading the six tritic pairs of each block defines bijections
\[
 \beta_N:Y_N\xrightarrow{\ \cong\ }(\FF_3^6)^N
\]
that commute with truncation. Each state of \(Y_N\) has exactly \(729\) block prolongations in \(Y_{N+1}\).
\end{proposition}
\begin{proof}
On each turn, the record retains the tritic pair and its carry. The three pairs produce the three distinct outputs \(\chi(\varepsilon,\varepsilon)=\varepsilon\), and the record allows recovery of the branch used, including the zero-carry branch. The six concatenations therefore produce all words of \(\FF_3^6\). Two different words are distinguished at the first recorded pair in which they differ. Recovery of the prefix proves compatibility with truncation. By induction on \(N\), every word of length \(N\) has a unique antecedent state; the next word can be any of the \(3^6=729\) possibilities.
\end{proof}

Let \(\Omega_{\Gamma}^{\rm cal}=\varprojlim_NY_N\). The preceding bijections induce
\[
 \Omega_{\Gamma}^{\rm cal}\cong(\FF_3^6)^{\mathbb N}.
\]
The correspondence and its inverse are continuous: each finite coordinate depends on a finite prefix. The cylinders are open and closed. Each cylinder of depth \(N\) is partitioned into its \(729\) subcylinders of depth \(N+1\); each contains an infinite prolongation, obtained, for example, by continuing with the zero branch in all subsequent blocks.

% Antecedentes efectivos del continuo conjunto y del universo interno.
% Se imprimen en ambos destinos; la etiqueta de universo no reemplaza su regla.
% Punto de inclusión común: después de APP--TRIT--TPK y los bloques 729,
% antes de la cobertura arquimediana y de la carta interna de Feigenbaum.
% No duplica 00a ni incluye clasificación cardinal posterior.
% La regla semántica transfinita se declara como tal, no como consecuencia
% de la sola existencia de prolongaciones finitas.
\begingroup
\let\HMTFGBaseSection\subsection
\let\HMTFGBaseSubsection\subsubsection
\let\HMTFGBaseSubsubsection\paragraph
\let\section\HMTFGBaseSection
\let\subsection\HMTFGBaseSubsection
\let\subsubsection\HMTFGBaseSubsubsection
% Propietario reproducido por unidad demostrativa; véase PROPIETARIOS_CONTINUO_UNIVERSO.json.
% Origen: XI ES CUMPLIMIENTO_EDITORIAL_20260928, sections/continuo_conjunto.tex:1--629.
% Cuerpos y pruebas conservados; sólo namespace, rutas y remisiones contextuales declaradas.
\section{Joint construction of the continuum through compatible histories and prolongations}
\label{hmtfg:base:iv:continuo-conjunto}
The arithmetic expressed by a normal form preserves a precise readout of the state that produces it. To situate that readout within HMT, the histories, sheets and operations acting on them must also be retained. This chapter brings together that common substrate. Its different constructions arise from the same tree of APP–TRIT–TPK executions; the expository divisions allow their properties to be studied without turning them into independent domains chosen by a subsequent application.

The starting point is the formal kernel presented above. An admissible emission contains the APP operation, TRIT orientation, sheet change, remainder and quotient, together with the effective TPK memory update. When only a remainder is expressed, the execution history still belongs to the starting object. Two executions with identical final readout are not identified on that basis alone.

\input{inserciones_20260930/feigenbaum/antecedentes_comunes/supervivencia_punto_fijo.tex}

\subsection{Admissible histories and memory by composition}
Let \(H_0:=\{\ast\}\). Level \(H_1\) is the finite nonempty set of admissible APP--TRIT seeds, and, for \(k\geq1\), \(H_k\) is the set of enriched histories of length \(k\) that admit compatible TPK prolongations to every finite horizon. Its elements are constructed by successive emissions from the initial domain of the kernel. We write \(\zeta\varepsilon\) for the prolongation of \(\zeta\) by an admissible emission \(\varepsilon\), and \(\rho_k(\zeta\varepsilon)=\zeta\) for deletion of the last emission. The label produced by the emission is retained, even when different labels express the same scalar state. The space of complete histories is defined by
\[
 H_\infty^{\rm enr}:=\varprojlim_k(H_k,\rho_k).
\]
The cofinal subsystem \(X_N:=H_{6N}\), with \(X_0=H_0\), agrees with the sextet clock used in the cardinality proof; cofinality gives a canonical identification \(H_\infty^{\rm enr}\cong\varprojlim_N(X_N,\rho_{6(N+1),6N})\).

On the two arithmetic sheets, the local data satisfy the exact divisions

\[
\delta_\diamond=r_\diamond+9q_\diamond,
\qquad 0\leq r_\diamond<9,
\qquad \diamond\in\{\Sigma,\Pi\}.
\]
The oriented TRIT coordinate likewise admits an expression \(\Delta=o_r+3\kappa\), with the oriented representative declared by the kernel. Preserving the quotient allows those representations to be composed without replacing an emission by its isolated remainder.

In a memory coordinatization, the update of an emission is written

\[
U_\varepsilon(m)=C_\varepsilon m+\kappa_\varepsilon.
\]
For a trajectory \(\alpha=(\varepsilon_1,\ldots,\varepsilon_k)\), the effective operation is the ordered composition of these updates. Its linear part and its translation are

\[
C_\alpha=C_{\varepsilon_k}\cdots C_{\varepsilon_1},
\qquad
\kappa_\alpha=
\sum_{j=1}^{k}C_{\varepsilon_k}\cdots
C_{\varepsilon_{j+1}}\kappa_{\varepsilon_j}.
\]
The empty product in the last sum is the identity.

\begin{proposition}[Affine compatibility of concatenation]
\label{hmtfg:base:prop:continuo-compatibilidad-afin}
If \(\alpha\) precedes \(\beta\), then

\[
C_{\beta\alpha}=C_\beta C_\alpha,
\qquad
\kappa_{\beta\alpha}=C_\beta\kappa_\alpha+\kappa_\beta,
\qquad
U_{\beta\alpha}=U_\beta U_\alpha.
\]
\end{proposition}
\begin{proof}
Substituting \(U_\alpha(m)\) into \(U_\beta\) produces \(C_\beta C_\alpha m+C_\beta\kappa_\alpha+\kappa_\beta\). Decomposition into linear part and translation gives the three identities. Associativity of composition proves that the result does not depend on an auxiliary subdivision of the trajectory.
\end{proof}

This cocycle identity is the precise content of accumulated memory. The quantity transported by a second trajectory depends on the coordinatization update it performs; in general, it is not the untransported sum of two records.

The nonadic return preserves a second distinction. In phase and turn coordinates, with \(\bar r\in\{0,\ldots,8\}\), its representation is

\[
\sigma_9(w,r)=
\left(w+\left\lfloor\frac{\bar r+1}{9}\right\rfloor,
r+1\pmod9\right),
\qquad
\sigma_9^9(w,r)=(w+1,r).
\]
During nine steps, exactly one crossing from \(8\) to \(0\) occurs. The phase returns and the turn coordinate increases. The enriched holonomy also retains the memory updates accompanying that path.

\subsection{Tree refinement and measure of histories}
The initial measure is the canonical census measure: \(\mu(\ast)=1\) and \(\mu(c)=1/|H_1|\) for each seed \(c\in H_1\). In this section each history has a finite positive number of prolongations. Deletion of emissions defines the refinement maps. Admissibility of the kernel supplies the prolongations; when a realization uses a neutral emission, this explicitly gives a section of \(\rho_k\). Existence of that section refers to the tree containing that emission, not to every arbitrary restriction of routes.

Let \(d(\zeta)\geq1\) be the number of admissible emissions at \(\zeta\). The APP order and the census of prolongations determine the canonical law \(p_{\rm can}(\varepsilon\mid\zeta)=1/d(\zeta)\). Without additional parameters, this law defines

\[
\mu(\zeta\varepsilon)=\frac{\mu(\zeta)}{d(\zeta)}.
\]
A nonuniform weighting may replace \(p_{\rm can}\) only when it has already been produced and added to the enriched record. The correlative constructor of this chapter uses only \(p_{\rm can}\); its measure is therefore a function of the APP--TRIT--TPK tree and not of an external choice.

\begin{proposition}[Cylindrical compatibility and conditional expectation]
\label{hmtfg:base:prop:continuo-esperanza-condicional}
At each level we have

\[
\sum_{\rho_k(y)=x}\mu(y)=\mu(x).
\]
In the spaces \(L^2(H_k,\mu_k)\), the maps

\[
(\mathsf J_kf)(y)=f(\rho_k(y)),
\qquad
(E_kg)(x)=\sum_{\rho_k(y)=x}\frac{\mu(y)}{\mu(x)}g(y)
\]
satisfy \(E_k=\mathsf J_k^*\), \(E_k\mathsf J_k=I\) and \(\|\mathsf J_kf\|=\|f\|\).
\end{proposition}
\begin{proof}
The first identity is the normalization of local probabilities. Applying it to \(|f(x)|^2\) gives the isometry. Grouping the sum by fibers of \(\rho_k\) gives the formula for \(E_k\); on a function constant on each fiber, that formula returns its value.
\end{proof}

The compatible masses define a probability measure on the space of infinite histories: the cylinder of a finite history receives its mass \(\mu(\zeta)\). The measure extends from the cylinder algebra, and its uniqueness follows because cylinders generate the \(\sigma\)-algebra. If the tree is finitely branching and every history has a prolongation, the inverse limit of its levels under emission deletion is nonempty. In the cylinder topology, every nonempty cylinder has positive measure.

This measure distinguishes histories that converge to the same representation. The uniform measure on terminal values would be a different readout, since it would remove the multiplicities that the construction has just preserved.

\input{inserciones_20260930/feigenbaum/antecedentes_comunes/ch_memoria_y_naturalidad.tex}

\subsection{Transport by sheets and exact accounting of outputs}
First consider the space with an orthonormal basis formed by the histories at level \(k\). The transport preserving the label of each prolongation is

\[
\mathcal U_ke_\zeta=
\sum_{\varepsilon\ \mathrm{admisible\ en}\ \zeta}
\sqrt{p_{\rm can}(\varepsilon\mid\zeta)}\,e_{\zeta\varepsilon}.
\]
The successor fibers of distinct antecedents are disjoint. Normalization of \(p_{\rm can}\) therefore proves \(\mathcal U_k^*\mathcal U_k=I\). The representation by history masses is related to this one through the diagonal change of basis incorporating \(\sqrt{\mu(\zeta)}\).

The sheet label is part of the complete register of a history. We write
\[
\widehat\zeta=((\zeta_0,s_0),\ldots,(\zeta_k,s_k)),\qquad
s_j\in\{\Sigma,\Pi\},
\]
where \(s_j\) is the active sheet at step \(j\). Prolongation adds the emission and the sheet \(s_{k+1}\) determined by the operational TPK selector, preserving the entire prefix, including the initial sheet. In the basis of these lifted histories,
\[
\mathcal U_ke_{\widehat\zeta}
=\sum_{\varepsilon\ \mathrm{admisible\ en}\ \widehat\zeta}
\sqrt{p_{\rm can}(\varepsilon\mid\widehat\zeta)}\,e_{\widehat\zeta\varepsilon},
\qquad
P_ke_{\widehat\zeta}=\mathbf1_{s_k=\Sigma}e_{\widehat\zeta},\quad
Q_k=I-P_k.
\]
Thus \(P_k\) reads the current sheet, while histories remain distinguished by their complete prefix. Two initial histories reaching the same active sheet retain distinct successor fibers. The isometry proof through disjoint collections of prolongations applies literally to this basis.
Tritic orientation is another coordinate of the record: the value \(0\) preserves the threshold readout on each sheet and does not add a third sheet to this decomposition. We now separate sheet preservation and sheet exchange:

\[
A_k=P_{k+1}\mathcal U_kP_k+Q_{k+1}\mathcal U_kQ_k,
\qquad
\eta_k=Q_{k+1}\mathcal U_kP_k+P_{k+1}\mathcal U_kQ_k.
\]
\begin{proposition}[Gram identity by sheets]
\label{hmtfg:base:prop:continuo-gram-hojas}
For a linear transport \(U\), whether isometric or not, the preceding decomposition satisfies

\[
A^*A+\eta^*\eta=PU^*UP+QU^*UQ.
\]
In particular, for \(U=\mathcal U_k\),

\[
A_k^*A_k+\eta_k^*\eta_k=I.
\]
\end{proposition}
\begin{proof}
In \(A^*A\) and \(\eta^*\eta\), mixed terms containing \(P_{k+1}Q_{k+1}\) vanish. Summing the remaining terms yields \(P_kU^*(P_{k+1}+Q_{k+1})UP_k\), and its analogue with \(Q_k\). The first identity follows from complementarity of the projectors; the second also uses \(U^*U=I\).
\end{proof}

The first identity is the general formula. Replacing its right-hand side by \(U^*U\) requires the cross blocks of that Gram matrix to vanish. Isometric preservation of history transport provides that condition here; it should not automatically be attributed to a subsequent compression.

Write \(F_{0,0}=I\), \(F_{0,n}=A_{n-1}\cdots A_0\) and \(T_N=F_{0,N}^*F_{0,N}\).

\begin{theorem}[Finite decomposition with terminal residual]
\label{hmtfg:base:thm:continuo-descomposicion-terminal}
For each horizon \(N\),

\[
I=\sum_{n=0}^{N-1}
F_{0,n}^*\eta_n^*\eta_nF_{0,n}+T_N.
\]
The sequence \(T_N\) is positive and decreasing, and has a positive strong limit \(T_\infty\). The sum of outputs converges strongly to \(I-T_\infty\).
\end{theorem}
\begin{proof}
Proposition~\ref{hmtfg:base:prop:continuo-gram-hojas} gives \(T_n-T_{n+1}=F_{0,n}^*\eta_n^*\eta_nF_{0,n}\). The telescoping sum proves the finite equality and monotonicity. For each vector, the quadratic forms decrease to a limit; polarization defines a bounded positive form and, by the representation theorem for bounded forms, an operator \(T_\infty\). Convergence is strong because, for a positive operator \(S\leq I\), \(\|Sv\|^2\leq\langle v,Sv\rangle\); this is applied to \(S=T_N-T_\infty\).
\end{proof}

The terminal residue is retained until its extinction in the domain under consideration is proved. This distinction prevents a phase return, by itself, from becoming an assumption of exhaustion of all histories.

\subsection{Oriented balance and cancellation under refinement}
A history carries its accumulated orientation \(o(h)\) and the counts of its visits to the two sheets. The corresponding local balance is

\[
b(h)=o(h)\bigl(N_\Sigma(h)-N_\Pi(h)\bigr),
\qquad
b^+=\frac{|b|+b}{2},\quad b^-=\frac{|b|-b}{2}.
\]
When a reader expresses the mean over prolongations, define \(b_c=Eb_f\). Orientation remains in the argument of \(b_f\); it is not removed before averaging.

\begin{proposition}[Exact cancellation defect]
\label{hmtfg:base:prop:continuo-defecto-cancelacion}
The quantity

\[
\kappa_{\mathrm{can}}=
\frac{E|b_f|-|Eb_f|}{2}
\]
is nonnegative and satisfies

\[
E(b_f^+)=(b_c)^++\kappa_{\mathrm{can}},
\qquad
E(b_f^-)=(b_c)^-+\kappa_{\mathrm{can}}.
\]
\end{proposition}
\begin{proof}
The weighted triangle inequality gives \(|Eb_f|\leq E|b_f|\). Writing the positive and negative parts through \((|b|\pm b)/2\) yields both equalities.
\end{proof}

The information omitted when only the balance is expressed is the common cancellation of opposite contributions. This quantity is determined from the same prolongations used by the preceding transport. The factor \(1/9\) appears only when the fiber contains nine prolongations with equal weights; conditional expectation is the rule that remains valid for general branching and weightings.

\subsection{Ternary incidence of the sheets and rectangular complex}
The two sheets determine the module \(V=\mathbb F_3e_\Sigma\oplus\mathbb F_3e_\Pi\). Its four projective directions are the lines generated by \(e_\Sigma,e_\Pi,e_\Sigma+e_\Pi,e_\Sigma-e_\Pi\). They give rise to six unordered pairs. The eight nonzero vectors are their oriented representatives. Thus, the counts \(4,6,8\) belong to distinct operations on the same sheet module.

On the six pair positions, we define

\[
(Ba)_{\{L,L'\}}=a_L+a_{L'},
\qquad
s(q)=\sum_{\{L,L'\}}q_{\{L,L'\}}.
\]
Each direction occurs in three pairs; hence \(sB=0\) in \(\mathbb F_3\). The reader \(W(q)=\mathbf1_{s(q)=0}-1/3\) is invariant under \(q\mapsto q+Ba\). On the uniform ambient space \(\mathbb F_3^6\), its moments are

\[
\mathbb EW=0,\qquad
\mathbb EW^2=\frac29,\qquad
\mathbb EW^3=\frac2{27}.
\]
Indeed, \(s\) is surjective and each fiber contains \(3^5\) elements. The reader equals \(2/3\) with probability \(1/3\) and \(-1/3\) with probability \(2/3\). These means describe the uniform ambient space; the distribution induced by a subset of histories must be calculated with its own multiplicities, already preserved in \(H_k\).

To fix the ports without losing multiplicities, at depth \(k\) retain the labeled readouts \(\lambda_\Sigma(\zeta)=(\zeta,\Sigma)\) and \(\lambda_\Pi(\zeta)=(\zeta,\Pi)\). Their sets are
\[
\mathscr P_k^\Sigma=\{\lambda_\Sigma(\zeta):\zeta\in H_k\},\qquad
\mathscr P_k^\Pi=\{\lambda_\Pi(\zeta):\zeta\in H_k\}.
\]
They are finite and nonempty because \(H_k\) is. Values, quotients and remainders are expressed on those ports; equality of values does not identify their labels. The product \(\mathscr P_k^\Sigma\times\mathscr P_k^\Pi\) is the domain of the rectangular readout, within which the support of admissible incidences is recorded.

With this level fixed, write \(I=\mathscr P_k^\Sigma\), \(J=\mathscr P_k^\Pi\) and \(A_m=\mathbb Z/3^m\mathbb Z\). The APP lexicographic order determines the minimal surviving branch and, on it, the reference ports \(i_0,j_0\). No external choice is added: changing coordinates produces a canonically isomorphic complex. The maps

\[
\kappa(u,v)_{ij}=u_i-v_j,
\qquad
(\delta w)_{ij}=w_{ij}-w_{ij_0}-w_{i_0j}+w_{i_0j_0}
\]
define the following sequence, valid over the finite ring \(A_m\):

\[
0\longrightarrow A_m\longrightarrow
A_m^I\oplus A_m^J\mathrel{\mathop{\longrightarrow}^{\kappa}}
A_m^{I\times J}\mathrel{\mathop{\longrightarrow}^{\delta}}
A_m^{(I\setminus\{i_0\})\times(J\setminus\{j_0\})}
\longrightarrow0.
\]
\begin{proposition}[Exactness of rectangular incidence]
\label{hmtfg:base:prop:continuo-exactitud-rectangular}
The preceding sequence is exact, with first arrow given by the constant diagonal.
\end{proposition}
\begin{proof}
The equality \(u_i-v_j=0\) for each pair forces all coordinates of \(u\) and \(v\) to be the same constant. Direct substitution gives \(\delta\kappa=0\). If \(\delta w=0\), take \(u_i=w_{ij_0}\) and \(v_j=w_{i_0j_0}-w_{i_0j}\). Then \(u_i-v_j=w_{ij}\). Finally, any matrix of the last term is extended with zero reference row and column; its image under \(\delta\) is the initial matrix. The proof uses only addition and subtraction, so it does not require division by a unit of the ring.
\end{proof}

Reducing from \(A_{m+1}\) to \(A_m\) while keeping \(k,\mathscr P_k^\Sigma,\mathscr P_k^\Pi,i_0,j_0\) fixed commutes with both maps. The reconstruction formulas are the same at each modular depth. This is the arithmetic compatibility proved in this section. The prolongation of histories \(k\to k+1\) preserves prefix labels; its maps between port sets are a datum distinct from the ring change \(m\to m+1\).

These maps are obtained from emission removal:
\[
\rho_k^\Sigma(\zeta\varepsilon,\Sigma)=(\zeta,\Sigma),\qquad
\rho_k^\Pi(\zeta\varepsilon,\Pi)=(\zeta,\Pi).
\]
The minimal surviving branch of the APP order just constructed has compatible prolongations. Its two labels define compatible reference ports at every level. Coordinate transport is precomposition:
\[
(L_ku)_{i'}=u_{\rho_k^\Sigma(i')},\quad
(L_kv)_{j'}=v_{\rho_k^\Pi(j')},\quad
(L_kw)_{i'j'}=w_{\rho_k^\Sigma(i'),\rho_k^\Pi(j')}.
\]
For the last module of the sequence, first extend the matrix by zero on the reference row and column, apply this same formula and restrict to the new complement of the reference ports. This rule defines \(L_k\) even when a new port has the marked port as its antecedent.

\begin{proposition}[Naturality of the rectangular complex under prolongation]
\label{hmtfg:base:iv:rectangular-natural}
The preceding maps satisfy
\(L_k\kappa_k=\kappa_{k+1}L_k\) and
\(L_k\delta_k=\delta_{k+1}L_k\).
They also commute with coefficient reduction
\(A_{m+1}\to A_m\) and compose according to prefix removal.
\end{proposition}
\begin{proof}
The first identity follows from \(u_{\rho_k^\Sigma(i')}-v_{\rho_k^\Pi(j')}\). In the second, each of the four terms defining \(\delta\) is transported by the same precomposition. Compatibility of the marked ports identifies their rows and columns; when an antecedent is the marked port, the four terms cancel and agree with extension by zero. All formulas are sums and differences with integer coefficients, so they commute with modular reduction. Composition of precompositions is precomposition with the composite prefix, which proves the last assertion.
\end{proof}

\subsection{Memory complex and compatibility of incidences}
The same accumulated memory coordinate \(c_y\in\mathbb Z\), in a residence \(y\), defines a free integer complex. Its reduction to \(A_m\) uses the same integer register before reducing it; successive depths are therefore compatible. Its generators are \(f_y\) in degree \(2\), \(e_y,b_y,r_y\) in degree \(1\), and \(g_y\) in degree \(0\):

\[
\partial f_y=3c_y e_y+b_y,\qquad
\partial e_y=g_y,\qquad
\partial b_y=-3c_y g_y,\qquad
\partial r_y=0.
\]
The identity \(\partial^2=0\) is verified on \(f_y\) through \(3c_yg_y-3c_yg_y=0\), and is immediate on the other generators.

If a trajectory \(\alpha:y\to y'\) updates memory, its transport in the complex sends the correspondingly named generators to their new representatives, except

\[
\Psi_\alpha(b_y)=b_{y'}+3(c_{y'}-c_y)e_{y'}.
\]
\begin{proposition}[Naturality of the memory complex]
\label{hmtfg:base:prop:continuo-naturalidad-memoria}
We have \(\partial\Psi_\alpha=\Psi_\alpha\partial\) and \(\Psi_{\beta\alpha}=\Psi_\beta\Psi_\alpha\). The maps extend linearly over the chosen coefficient ring.
\end{proposition}
\begin{proof}
On \(f_y\), the image of its boundary is \(3c_ye_{y'}+b_{y'}+3(c_{y'}-c_y)e_{y'} =3c_{y'}e_{y'}+b_{y'}\). On \(b_y\), the boundary of its image is \(-3c_{y'}g_{y'}+3(c_{y'}-c_y)g_{y'}=-3c_yg_{y'}\). The other two cases are immediate. In a concatenation, memory differences sum telescopically, proving the second identity.
\end{proof}

The relation between two representatives of the same residence can also be written as an explicit boundary. For an integer coordinate \(v\), reduced in the same ring where appropriate, let

\[
z^-=r_y,\qquad z^+=r_y+3c_yv e_y+v b_y,
\qquad h_*=-vf_y.
\]
Then \(\partial h_*=z^--z^+\). Under a transport that also updates \(v\), the correction \(K_\alpha=(v_{y'}-v_y)f_{y'}\) satisfies \(K_{\beta\alpha}=K_\beta+\Psi_\beta K_\alpha\). The symbols \(K_\alpha\) of this local homotopy do not denote the dodecaphasic seal \(K\): they belong to another domain and their function is specified by the preceding identity. Their function can also be checked in the equality

\[
z^+_{y'}-\Psi_\alpha z^+_y=\partial K_\alpha.
\]
Both sides equal \((v_{y'}-v_y)(3c_{y'}e_{y'}+b_{y'})\).

\begin{proposition}[Homology and memory of the local complex]
\label{hmtfg:base:prop:continuo-homologia-memoria}
Over \(\mathbb Z\), or after reduction to \(A_m\), we have

\[
H_0(\Gamma)=H_2(\Gamma)=0,
\qquad H_1(\Gamma)\cong A_m[r_y]
\]
in the modular realization, and \(H_1(\Gamma)\cong\mathbb Z[r_y]\) in the integer realization.
\end{proposition}
\Needspace{9\baselineskip}
\begin{proof}
Define the homotopy of degree \(1\) by \(h(g_y)=e_y\), \(h(b_y)=f_y\), and zero on the other generators. Let \(p\) be the projection onto the summand generated by \(r_y\) and \(\iota\) its inclusion. Checking on each generator gives

\[
\partial h+h\partial=I-\iota p.
\]
In particular, on \(b_y\) the terms \(3c_ye_y\) and \(-3c_ye_y\) cancel. The complex deformation retracts onto the module generated by \(r_y\), concentrated in degree \(1\), which proves the assertion.
\end{proof}

The transport \(\Psi_\alpha\) is a shear of determinant \(1\). Thus, this complex explicitly preserves how \(c_y\) changes at the chain level, but its local homology does not distinguish the values of that coordinate. The full memory remains in the labels and transports of the source state; it is not recovered from \(H_1\) alone. A memory-dependent numerical torsion requires the specified assembly and determinantal reading, in addition to this local complex.

\subsection{Assembly with terminal paths}
For a finite horizon, the residences and their prolongations form a path category. Each residence \(\nu\) is assigned the free module \(W(\nu)=A_m[\operatorname{Path}(\nu,\mathrm{Term})]\), which preserves each terminal continuation as a generator. An initial path acts on \(W\) by precomposition. The complex \(\Gamma(\nu)\) of the preceding section acts covariantly through \(\Psi\).

The assembly uses both transports in the same bigraded module:

\[
\mathcal B_{q,r}=
\bigoplus_{\nu_0\to\cdots\to\nu_q}
W(\nu_q)\otimes_{A_m}\Gamma_r(\nu_0).
\]
The normalized path complex is used: degenerate chains containing identities are omitted. The finite horizon then bounds the number of strict prolongations and the horizontal degree. The horizontal boundary removes an intermediate residence by composing the adjacent arrows. At the first endpoint it transports \(\Gamma\); at the last it precomposes \(W\). The alternating signs yield the path boundary \(d_{\mathrm{bar}}\). The total differential is

\[
D=d_{\mathrm{bar}}+(-1)^q\partial_\Gamma.
\]
The face identities cancel the terms of \(d_{\mathrm{bar}}^2\) in pairs. The naturality of Proposition~\ref{hmtfg:base:prop:continuo-naturalidad-memoria} implies that \(d_{\mathrm{bar}}\) commutes with \(\partial_\Gamma\); the factor \((-1)^q\) cancels the cross terms of the square. Therefore, \(D^2=0\). Compatibility is proved on path generators and is not introduced by a label of membership in a complex.

Terminal evaluation likewise does not require choosing a particular history. For each continuation \(h\), the effective composition of its updates is applied. The result preserves the pair formed by \(h\) and its reading, with its multiplicity. If \(k<\ell<N\), the continuations decompose as a disjoint union of prefixes up to \(\ell\) and their continuations up to \(N\). Proposition~\ref{hmtfg:base:prop:continuo-compatibilidad-afin} proves that evaluating directly or evaluating these two parts produces the same reading. This yields the naturality square of the terminal.

The path modules, incidences, memory transport and their differentials arise from the same prolongations. Their assembly does not consist in adding to a state five coordinates that would be preserved by definition: their relations are constructed through sums, boundaries and effective compositions, and are verified on their generators.

\subsection{Determinants, limit and arithmetic reading}
In each realization, the local free complex \(\Gamma\) admits its graded determinant line

\[
\operatorname{Det}\Gamma=
\bigotimes_r\left(\bigwedge^{\mathrm{rk}\Gamma_r}\Gamma_r\right)^{(-1)^r}.
\]
The total assembly of the preceding section, in turn, has the modules \(\operatorname{Tot}_d\mathcal B=\bigoplus_{q+r=d}\mathcal B_{q,r}\). At a finite horizon and over the same ring, they are free of finite rank and only finitely many of them are nonzero. Its line is
\[
\operatorname{Det}(\operatorname{Tot}\mathcal B)=
\bigotimes_{q,r}
\left(\bigwedge^{\operatorname{rk}\mathcal B_{q,r}}
\mathcal B_{q,r}\right)^{(-1)^{q+r}}.
\]
This formula follows from decomposition as a direct sum in each degree. It distinguishes the line of the local complex from that of the assembly with all paths; their bases are ordered through the preserved labels. Modular reduction commutes with these exterior powers because the modules and their bases come from the same free integer modules at a fixed horizon.
The line is defined by the modules and their degrees, both over \(\mathbb Z\) and after reduction to \(A_m\); the negative exponent denotes the dual module. A nonzero numerical torsion additionally requires specifying the complex over a field, the bases and the relevant homology identifications. To calculate Smith factors, the integer matrix produced before reduction is used, not an arbitrary lift of a modular matrix. A square integer block with nonzero determinant is invertible over \(\mathbb Q\); its reduction is invertible over \(A_m\) only when that determinant is a unit, that is, when it is not divisible by \(3\). Refinement of memories and incidences allows the stability of the relevant invariant factors to be formulated and checked; these conditions do not follow from the existence of the determinant line.

Passage to the limit preserves the compatible projective systems of remainders, histories and their multiplicities. The squares proved in the preceding sections remain commutative because their maps are given by the same formulas at each level. For an Archimedean coordinate, the additional step consists in expressing nonempty cylinder intervals, compatible by inclusion and with diameters tending to zero. Their nested compact closures have a one-point intersection: existence follows by compactness and uniqueness by the diameter bound. That is the value of the limit readout. If the intervals are half-open, the point may belong only to their closures; the boundary representation rule then determines its positional expression. The value does not replace the history that produced it. The other representation preserves the incidences of the same state. Both destinations share their origin without thereby receiving the same information.

% Procedencia: continuo_conjunto.tex:434--444; remisiones contextuales a
% publicaciones posteriores, no premisas de las pruebas del continuo ni de Feigenbaum.
% Se conserva la descripción positiva sin exigir etiquetas de un anfitrión concreto.
The coinductive regions, the dodecaphase record and exceptional incidence are expressed from that substrate through their respective maps. The Hilbert-space realization subsequently preserves phase labels and specifies its own unital inclusion; duality and screens act on the codomains constructed there. The joint operations brought together here transmit histories, sheets, orientation and memory to those representations. Identifying a particular functional with an external analytic form in turn requires writing its action on the generators and verifying its inner products: the conservation identities of this chapter do not replace that composition.

\subsection{Mathematical scope of the joint construction}
The state, tree and measure constructions; the correlative operations; and the terminal assembly with naturality and limit are developed in full in this article. In this chapter their balance, incidence and memory identities have been formulated and proved on the same generators. The five constructions thus remain operations of a single tree. Each specialization keeps its exact hypotheses visible: transport isometry, canonical weight law, treatment of the terminal remainder and nonsingularity of the determinantal block where applicable.

The preceding equations prove the memory cocycles, weighted cancellation, rectangular exactness and compatibility of the memory complex at every finite horizon, as well as their compatible passage to the limit. A subsequent analytic identification must compose its operator on these generators and verify the additional properties of its codomain; it neither alters nor replaces the joint construction proved here.

\iffalse
% Formulación preliminar conservada fuera del ensamblado; sustituida por la
% versión tipada y verificada que se carga al final de este bloque.
\subsection{Joint emergence of the five constructions}
\label{hmtfg:base:sec:cinco-construcciones-continuo-ch}

The preceding operations are now brought together without turning them into five branches. For a surviving history \(h\in X_k\), a horizon \(N>k\) and a modular depth \(m\geq1\), let the common base record be
\[
 \mathcal B_{k,N}(h):=
 \bigl(\mathscr T_{k,N}(h),
       (\varepsilon_\alpha,s_\alpha,o_\alpha,
        r_{\Sigma,\alpha},q_{\Sigma,\alpha},
        r_{\Pi,\alpha},q_{\Pi,\alpha},
        \operatorname{Carr}_\alpha,
        \kappa_\alpha,\partial_\alpha,
        \operatorname{Surv}_\alpha)_\alpha;
       \operatorname{src},\operatorname{tgt},\circ,\rho\bigr),
 \tag{C1}
 \label{hmtfg:base:eq:base-comun-registros}
\]
where \(\mathscr T_{k,N}(h)\) is the subtree of prolongations of \(h\) up to horizon \(N\). Each component of the record is an output of the APP divisions, TRIT orientation or TPK updates described in the preceding sections; none originates in a classical problem chosen afterward.

On this same base, define the five structured objects
\[
\begin{aligned}
 \mathcal F_{\rm sp}(\mathcal B_{k,N})
   &:=(\ell^2(X_j,\mu_j),\mathcal U_j,A_j,\eta_j,T_j)_{k\leq j\leq N},\\
 \mathcal F_{\rm sol}(\mathcal B_{k,N})
   &:=(b,\kappa_{\rm can},C_\alpha,\kappa_\alpha)_{\alpha\subset\mathscr T_{k,N}},\\
 \mathcal F_{\rm gau}(\mathcal B_{k,N})
   &:=(I_j,J_j,B_j,\kappa_j,\delta_j)_{k\leq j\leq N},\\
 \mathcal F_{\rm coh}^{(m)}(\mathcal B_{k,N})
   &:=(\Gamma_y\otimes A_m,\partial,\Psi_\alpha,H_\bullet)_{y,\alpha\subset\mathscr T_{k,N}},\\
 \mathcal F_{\rm det}^{(m)}(\mathcal B_{k,N})
   &:=(\operatorname{Tot}\mathcal B,\operatorname{Det}(\operatorname{Tot}\mathcal B))_{k,N,m}.
\end{aligned}
 \tag{C2}
\]
The discriminants \({\rm sp},{\rm sol},{\rm gau},{\rm coh},{\rm det}\) name, respectively, the spectral, solenoidal, gauge, cohomological and determinantal readouts of a single object. They are not input domains. If \(p_T\) forgets the structure added by reader \(T\) and returns the record \(\mathcal B_{k,N}(h)\), write \(\widetilde{\mathcal F}_T=(\mathcal B_{k,N},\mathcal F_T,p_T)\). The correlative output is the fiber product
\[
\begin{aligned}
 \mathcal J_{k,N}^{(m)}(h):={}&
 \widetilde{\mathcal F}_{\rm sp}
 \mathop{\times}_{\mathcal B_{k,N}(h)}
 \widetilde{\mathcal F}_{\rm sol}
 \mathop{\times}_{\mathcal B_{k,N}(h)}
 \widetilde{\mathcal F}_{\rm gau}\\
 &\mathop{\times}_{\mathcal B_{k,N}(h)}
 \widetilde{\mathcal F}_{\rm coh}^{(m)}
 \mathop{\times}_{\mathcal B_{k,N}(h)}
 \widetilde{\mathcal F}_{\rm det}^{(m)}.
\end{aligned}
 \tag{C3}
 \label{hmtfg:base:eq:constructor-correlativo}
\]
Equality of the five projections onto \(\mathcal B_{k,N}(h)\) requires preservation of the same tree, the same edges and the same record; this is why \eqref{hmtfg:base:eq:constructor-correlativo} is not a Cartesian product of five retrospectively assembled objects.

Deletion of the last layer of the tree induces \(\operatorname{Res}_{N+1,N}\). On the generators, the preceding propositions give the squares
\[
\begin{gathered}
 \operatorname{Res}_{N+1,N}\mathcal F_T(\mathcal B_{k,N+1})
 =\mathcal F_T(\mathcal B_{k,N}),
 \qquad T\in\{{\rm sp},{\rm sol},{\rm gau},{\rm coh}\},\\
 \operatorname{Det}(\operatorname{Tot}\mathcal B_{k,N+1})
 \simeq
 \operatorname{Det}(\operatorname{Tot}\mathcal K_{N+1})
 \otimes
 \operatorname{Det}(\operatorname{Tot}\mathcal B_{k,N}),\\
 (\partial\Psi_\alpha=\Psi_\alpha\partial),\qquad
 (L_k\kappa_k=\kappa_{k+1}L_k),\qquad
 (L_k\delta_k=\delta_{k+1}L_k).
\end{gathered}
 \tag{C4}
 \label{hmtfg:base:eq:naturalidad-D}
\]
The change \(A_{m+1}\to A_m\) also commutes with \(\partial\), \(\Psi\), \(\kappa\), \(\delta\) and the determinant lines. A single pro-system in \((N,m)\) is therefore obtained.

Let \(\mathcal X_\chi\) be the limit of compatible enriched histories. For \(\xi\in\mathcal X_\chi\), denote by \(\mathfrak L_N(\xi)\) the ordered prefix of its \(N\) complete emissions. The joint horizon fibers are
\[
 \mathfrak J_{N,\bullet}^{(m)}
 :=\coprod_{h\in X_0}
 \{(\mathfrak L,J):\mathfrak L\text{ is a branch of length }N,
       \ J=\mathcal J_{0,N}^{(m)}(h)\}.
 \tag{C5}
 \label{hmtfg:base:eq:fibras-conjuntas-horizonte}
\]
The mark \(\mathfrak L\) does not retain a tautological copy of \(\xi\): it is formed by the emissions and updates produced at each edge. The joint generator is
\[
\begin{aligned}
 \operatorname{Gen}_{\rm cont}:\mathcal X_\chi
   &\longrightarrow\varprojlim_{N,m}\mathfrak J_{N,\bullet}^{(m)},\\
 \xi&\longmapsto
   (\mathfrak L_N(\xi),\mathcal J_{0,N}^{(m)}(h_0))_{N,m},\\
 \mathcal C_{\rm cont}^{\rm disc}
   &:=\operatorname{im}(\operatorname{Gen}_{\rm cont}).
\end{aligned}
 \tag{C6}
 \label{hmtfg:base:eq:continuo-discreto}
\]

\begin{theorem}[Correlative emergence of the five constructions]
\label{hmtfg:base:thm:correlacion-cinco}
For every surviving history \(h\in X_k\), every horizon \(N\geq k+1\) and every depth \(m\geq1\), the fiber product \eqref{hmtfg:base:eq:constructor-correlativo} exists and is nonempty. Its five components emerge from the same APP--TRIT--TPK record, satisfy the naturality squares \eqref{hmtfg:base:eq:naturalidad-D}, base changes and passage to the limit. The image \(\mathcal C_{\rm cont}^{\rm disc}\) contains five inseparable aspects of the same continuum, not five continua or five independent maps.
\end{theorem}

\begin{proof}
Prolongability of the tree provides at least one edge at each interior vertex. The record of that edge determines isometric transport and its sheet decomposition, the oriented cocycle, the incidence ports, the memory complex and its determinant line. Thus \(\mathcal F_{\rm sp}\), \(\mathcal F_{\rm sol}\), \(\mathcal F_{\rm gau}\), \(\mathcal F_{\rm coh}^{(m)}\) and \(\mathcal F_{\rm det}^{(m)}\) are successively constructed on the same base, proving nonemptiness of \eqref{hmtfg:base:eq:constructor-correlativo}. None of the five constructions appears as an input coordinate.

The affine compatibility proposition proves memory composition; the Gram identity, rectangular exactness and naturality of the complex prove the first four lines of \eqref{hmtfg:base:eq:naturalidad-D}. Decomposition of the bicomplex structure by the last layer proves the multiplicative identity of the determinant line. All these operations are computed on the same generators and commute with modular reduction. Fiber, relation, number, orientation, carry, memory, boundary and route remain in \(\mathfrak L_N(\xi)\); multisection, terminals, and the linear and probabilistic readers remain in \(\mathcal J_{0,N}^{(m)}(h_0)\). By compatibility, the limit produces \eqref{hmtfg:base:eq:continuo-discreto} without disaggregation.
\end{proof}

Subsequent recognition of the single continuum constructed through its five spectral, solenoidal, gauge, cohomological and determinantal realizations belongs to the subsequent conventional language. It does not select its states or take part in its generation.
\fi

\input{inserciones_20260930/feigenbaum/antecedentes_comunes/continuo_cinco_construcciones_tipado.tex}

\endgroup
% Propietario reproducido por unidad demostrativa; véase PROPIETARIOS_CONTINUO_UNIVERSO.json.
% Origen: XI ES CUMPLIMIENTO_EDITORIAL_20260928, sections/ch_cierre_cardinal.tex:26--403.
% Cuerpos y pruebas conservados; sólo namespace, rutas y remisiones contextuales declaradas.
\subsection{Extension of states and transfinite semantic extension}

Let

\[
 \Sigma_{\rm HMT}
 =\bigl(\dot R_{\rm APP},\dot R_{\rm TRIT},\dot R_{\rm TPK},
 \dot R_{U},\dot R_{\Gamma_9},\dot R_{\rm Ext},\dot R_\rho,
 \dot R_{\pi},\dot R_{\varphi},\dot R_e,\dot R_{\alpha},
 \dot R_{\rm ar},\dot R_{\rm exc}\bigr)
\tag{33}\label{hmtfg:base:eq:ch-signatura}
\]

be the countable \textbf{relational} signature whose effective code is \(h\).
Each \(\dot R_S\) is the graph of the operator \(S\) on the codes where it is
defined; no external function is assumed to be total within each level. The
graphs \(\dot R_{\pi},\dot R_{\varphi},\dot R_e,\dot R_{\alpha}\) are the
output readers constructed by coinductive generation and dodecaphase closure;
they do not take the conventional values of the constants as arguments. Their
presence in \(\Sigma_{\rm HMT}\) ensures that the closure preserves the full
state: it simultaneously integrates the routes that express value and the
incidence that extends toward the exceptional realization.

At each horizon, the state code records the APP--TRIT--TPK coordinates, the
route, sheet, carry, boundary, and memory; at its representation stages it also
records \(P_N,\Phi_N,E_N,A_N\), where \(A_N\) is the prefix of
\(\alpha_{\rm HMT}\) produced by the same state. These coordinates do not
select the branch from outside: they are part of the history being encoded.
Once a primitive-recursive numbering of the finite coordinates of a state is
fixed, every inverse history
\(x_\infty=(x_N)_{N\geq1}\in X_\infty^{\rm enr}\) determines the code
\[
 \operatorname{code}(x_\infty)
 :=\{\langle N,j,c\rangle:(x_N)_j\text{ has code }c\}\subseteq\omega.
\]
The coherence equations
\(\rho_{N+1,N}(x_{N+1})=x_N\) make it possible to decode the history uniquely
from that real. We write \(m_\infty\) for the code of the realized history and
\(m_\infty(n)\) for its characteristic function under the fixed numbering.
The operator code \(h\) and this complete ledger are combined, by means of a
primitive-recursive pairing function, into the real

\[
 u_{h,m}:=
 \{\langle0,\langle n,h(n)\rangle\rangle:n<\omega\}
 \cup
 \{\langle1,\langle n,m_\infty(n)\rangle\rangle:n<\omega\}.
\tag{33a}\label{hmtfg:base:eq:ch-codigo-conjunto}
\]

The two primitive projections recover \(h\) and \(m_\infty\) from
\(u_{h,m}\).

The first extension is the metric and operator extension of the finite states.
For \(N\ge1\), let \(\widehat{\mathcal X}^{\rm enr,raw}_{6N}\) be the set of
raw enriched states of length \(6N\), and define the surviving subsystem
set-theoretically by
\[
\begin{aligned}
 X_N
 &:=\widehat{\mathcal X}^{\rm enr,surv}_{6N}\\
 &:=\Bigl\{x\in\widehat{\mathcal X}^{\rm enr,raw}_{6N}:\\[-.2ex]
 &\qquad\forall M>N\ \exists z\in
   \widehat{\mathcal X}^{\rm enr,raw}_{6M}\;\text{ such that}\\[-.2ex]
 &\qquad\operatorname{tr}_{6M,6N}(z)=x\ \wedge\
   z\text{ satisfies every intermediate TPK transition}\Bigr\}.
\end{aligned}
\]
and set
\[
 \rho_{N+1,N}:=\operatorname{tr}_{6N+6,6N}\!\upharpoonright X_{N+1},
\]
and
\[
 \operatorname{Ext}_N(x):=
 \{y\in X_{N+1}:(x,y)\in
 \operatorname{Ext}_{6N+6,6N}\}.
\]
The root level is included without any implicit exception:
\[
 X_0:=\{\ast\},\qquad
 \rho_{1,0}:X_1\to X_0\text{ is constant},\qquad
 \operatorname{Ext}_0(\ast):=X_1.
\]
The superscript \(\mathrm{surv}\) denotes the pruned tree of states possessing
compatible TPK descendants at every finite horizon. The TPK extension law
established in the preceding sections proves that the emitted seeds belong to
this tree; finite branching and Kőnig's lemma turn compatibility at every
horizon into an infinite branch. In particular,

\[
 \forall N\ge0\;\forall x\in X_N\quad
 \varnothing\ne\operatorname{Ext}_N(x)
 \subseteq\rho_{N+1,N}^{-1}(x).
\tag{34}\label{hmtfg:base:eq:ch-extension-no-vacia}
\]

By the same finitely branching tree lemma, every surviving seed has a
compatible extension

\[
 x_\infty=(x_N)_{N\ge1}\in
 X_\infty^{\rm enr}:=\varprojlim_N(X_N,\rho_{N+1,N}).
\tag{35}\label{hmtfg:base:eq:ch-limite-inverso}
\]

The second extension acts \textbf{afterward} on the complete code of that
realized history.
It is not inferred from the mere existence of the finite operator
\(\operatorname{Ext}_N\): at the set-theoretic scale, it formalizes the HMT principle that
only those objects possessing a genealogical derivation
from the state and its rules belong to the full domain. The transfinite semantic
extension is:

\[
 \begin{aligned}
  \mathfrak G_0(h,m_\infty)
  &:=\operatorname{TC}\bigl(\{\omega,u_{h,m},h,m_\infty\}\bigr),\\
  \mathfrak G_{\beta+1}(h,m_\infty)
  &:=\operatorname{Def}_{\Sigma_{\rm HMT}}
       \bigl(\mathfrak G_\beta(h,m_\infty)\bigr),\\
  \mathfrak G_\lambda(h,m_\infty)
  &:=\bigcup_{\beta<\lambda}\mathfrak G_\beta(h,m_\infty)
       \quad(\lambda\text{ limit}),\\
  \mathfrak G_{\rm HMT}(h,m_\infty)
  &:=\bigcup_{\beta\in\operatorname{Ord}}
       \mathfrak G_\beta(h,m_\infty).
 \end{aligned}
\tag{36}\label{hmtfg:base:eq:ch-clausura-transfinita}
\]

The Continuum Hypothesis and a target bijection are not introduced among the
conditions defining the successor operator; the operator is given
prospectively by

\[
 \operatorname{Def}_{\Sigma_{\rm HMT}}(M)
 :=M\cup
 \Bigl\{
   \{x\in M:\langle M,\in,
       (\dot R_S\cap M^{k_S})_{S\in\Sigma_{\rm HMT}}\rangle
       \models\varphi(x,\bar p)\}:
   \ulcorner\varphi\urcorner\in\omega,
   \ \bar p\in M^{<\omega}
 \Bigr\}.
\tag{37}\label{hmtfg:base:eq:ch-operador-def}
\]

That is, it simultaneously expresses all subsets definable by a finite formula
of the HMT relational language and parameters already produced by the
genealogy. It consults neither a target real value, a conventional constant, a
target cardinal, nor an external model. Formula~\eqref{hmtfg:base:eq:ch-operador-def}
defines the transfinite closure over the previously constructed seeds,
operators, enriched state, ledger, and double projection.

\begin{lemma}[Joint coding of the HMT program and the realized ledger]
\label{hmtfg:base:thm:ch-6}
The pair formed by the operator code \(h\) and the ledger \(m_\infty\) is
encoded by a single real \(u_{h,m}\subseteq\omega\), and every formula of
\(\Sigma_{\rm HMT}\) is uniformly translated into a formula of the language
\(\{\in,u_{h,m}\}\).

\end{lemma}
\begin{proof}
Both \(h\) and \(m_\infty\) are sequences of natural numbers: \(h\) encodes
the signature, the finite tables, and the operator programs; \(m_\infty\)
encodes one realized history, not the complete history space.
Formula~\eqref{hmtfg:base:eq:ch-codigo-conjunto} and its two decoders prove the first
assertion. Each symbol in~\eqref{hmtfg:base:eq:ch-signatura} denotes its \textbf{graph
relation encoded in \(h\)}, never the external graph of all its possible
evaluations. APP, TRIT, \(U_t\), \(\Gamma_9\), the restrictions \(\rho\), and
the finite extensions have primitive-recursive graphs on integer codes and,
therefore, \(\Delta_0\) formulas absolute between transitive levels. The two
representations are expressed as graph relations through the formulas

\[
 \begin{aligned}
 \dot R_{\rm ar}(x,y)
 &\;\Longleftrightarrow\;
 \forall k\in\omega\;\exists N\;\forall n\ge N\quad
 |\mu_n(x)-y|<2^{-k},\\
 \dot R_{\rm exc}(x,c,z)
 &\;\Longleftrightarrow\;
 \forall N\in\omega\quad z\!\upharpoonright N
   =E_N(x\!\upharpoonright N,c\!\upharpoonright N),
 \end{aligned}
\tag{37c$_0$}\label{hmtfg:base:eq:ch-grafos-publicacion}
\]

where \(\mu_n\) and \(E_N\) are the rational/finite functions whose codes
occur in \(h\). Every atomic formula containing an HMT symbol is inductively
replaced by its graph formula with parameter \(u_{h,m}\). Connectives and
quantifiers are preserved. This induction on syntactic complexity produces a
translation \(\varphi\mapsto\varphi^*\) such that, for every transitive level
\(M\) of~\eqref{hmtfg:base:eq:ch-clausura-transfinita} containing the parameters,

\[
 \langle M,\in,(\dot R_S\cap M^{k_S})_S\rangle
   \models\varphi(\bar a)
 \quad\Longleftrightarrow\quad
 M\models\varphi^*(\bar a;u_{h,m}).
\tag{37c}\label{hmtfg:base:eq:ch-eliminacion-uniforme}
\]

Thus, \eqref{hmtfg:base:eq:ch-operador-def} cannot consult an oracle, a target decimal
table, or an uncoded relation that introduces uncountably many objects in a
single step.
\end{proof}

HMT semantics is now defined \textbf{independently of identifying it with
\(\mathfrak G\)}. Let \(\mathcal T_0^{\rm HMT}\) be the set of canonical names
of the elements of
\(\mathfrak G_0(h,m_\infty)=
\operatorname{TC}(\{\omega,u_{h,m},h,m_\infty\})\). In particular, the
initial level contains the realized ledger \(m_\infty\), not every history
in \(X_\infty^{\rm enr}\) as primitive parameters. The admissible terms are
the well-founded trees obtained by

\[
 \begin{aligned}
 \mathcal T^{\rm HMT}_{\beta+1}
 &:={\mathcal T}^{\rm HMT}_\beta\cup
 \{\langle\beta,\ulcorner\varphi\urcorner,
       t_0,\ldots,t_{k-1}\rangle:
       t_i\in\mathcal T^{\rm HMT}_\beta\},\\
 \mathcal T^{\rm HMT}_\lambda&:=
       \bigcup_{\beta<\lambda}\mathcal T^{\rm HMT}_\beta,
 \qquad
 \mathcal T^{\rm HMT}:=
       \bigcup_{\beta\in\operatorname{Ord}}\mathcal T^{\rm HMT}_\beta.
 \end{aligned}
\tag{37a}\label{hmtfg:base:eq:ch-terminos}
\]

The value of a successor term is the subset of the preceding level defined by
\(\varphi\) in the induced relational structure, with the values of
\(t_0,\ldots,t_{k-1}\) as parameters. Before any set-theoretic comparison,
define

\[
 \operatorname{Ob}_{\rm sem}^{\rm HMT}(h,m_\infty)
 :=\{\llbracket t\rrbracket:t\in\mathcal T^{\rm HMT}\}.
\tag{37b}\label{hmtfg:base:eq:ch-objetos-semanticos}
\]

This is the rule of \textbf{generative exhaustiveness} adopted by HMT: every object in the entire domain must have one of those derivation trees and every admissible tree expresses an object. It is not a consequence of finite branching. The continuum hypothesis is not introduced into this rule; it is obtained through the proof developed in the article. Neither an enumeration of the reals nor a target bijection is introduced as data.

\begin{theorem}[Equivalence between term semantics and the genealogical hierarchy]
\label{hmtfg:base:thm:ch-6a}
The evaluation of the preceding terms satisfies

\[
 \operatorname{Ob}_{\rm sem}^{\rm HMT}(h,m_\infty)
 =\mathfrak G_{\rm HMT}(h,m_\infty).
\tag{37d}\label{hmtfg:base:eq:ch-equivalencia-semantica}
\]

\end{theorem}
\begin{proof}
By transfinite induction. At the initial level, the two classes coincide. If
every value of a term in \(\mathcal T^{\rm HMT}_\beta\) belongs to
\(\mathfrak G_\beta\), the semantics of a successor term is one of the subsets
in~\eqref{hmtfg:base:eq:ch-operador-def}. Conversely, every subset in
~\eqref{hmtfg:base:eq:ch-operador-def} is given by a formula code and a finite tuple of
parameters; by the induction hypothesis those parameters possess terms, and
~\eqref{hmtfg:base:eq:ch-terminos} forms the term that denotes it. At a limit, both
constructions take the same union.
\end{proof}

\Needspace{9\baselineskip}
The state encoder is compatible with truncations and sends every TPK extension
to a successor term. Let
\(\mathcal T_N^{\rm st}\subset\mathcal T^{\rm HMT}\) be the subclass of terms
encoding complete states of depth \(N\), and let
\(\tau_{N+1,N}:\mathcal T_{N+1}^{\rm st}\to\mathcal T_N^{\rm st}\) delete the
last block together with the corresponding truncation of the genealogical
ledger, carry, and memory. If \(J_N:X_N\to\mathcal T_N^{\rm st}\) is the
complete encoder, compatibility with \(\operatorname{Ext}_N\) proves the
well-typed identity

\[
 \forall x\in X_N\ \forall y\in\operatorname{Ext}_N(x)\qquad
 \tau_{N+1,N}\bigl(J_{N+1}(y)\bigr)=J_N(x).
\tag{37e$_0$}\label{hmtfg:base:eq:ch-naturalidad-codificador}
\]

Define within the semantics
\[
 \operatorname{Succ}^{\rm TPK}_{N+1}(J_N(x)):=
 \{t\in\mathcal T_{N+1}^{\rm st}:\exists y\in X_{N+1}\,
   (y\in\operatorname{Ext}_N(x)\ \wedge\ t=J_{N+1}(y))\}.
\]
Choose \(\beta\) such that \(\mathfrak G_\beta\) contains \(x\), \(X_N\),
\(X_{N+1}\), the finite graphs of
\(\operatorname{Ext}_N,J_N,J_{N+1}\), and their codes in \(h\). Then the
primitive-recursive formula for the graph of \(\operatorname{Ext}_N\) gives

\[
 J_{N+1}[\operatorname{Ext}_N(x)]
 =\operatorname{Succ}^{\rm TPK}_{N+1}(J_N(x)),
 \qquad
 J_{N+1}[\operatorname{Ext}_N(x)]
 \in\mathfrak G_{\beta+1},
\tag{37e}\label{hmtfg:base:eq:ch-sucesores-semanticos}
\]

and both sides are formed by the graph formula of the same TPK extension.
Equations~\eqref{hmtfg:base:eq:ch-naturalidad-codificador}--\eqref{hmtfg:base:eq:ch-sucesores-semanticos}
are the explicit semantic link: the subclass of successor terms is exactly the
image of the admissible extensions, and its truncation commutes. The complete
fiber of \(\rho\) is not identified with \(\operatorname{Ext}_N\), nor is
\(\operatorname{Def}\) identified with TPK merely by nomenclature.

\begin{lemma}[Transitivity, axioms, and well-ordering of the genealogical model]
\label{hmtfg:base:thm:ch-6b}
\(\mathfrak G_{\rm HMT}(h,m_\infty)\) is a transitive class, contains all
ordinals, satisfies Zermelo--Fraenkel set theory (ZF), and possesses a
definable global well-order; it therefore satisfies Choice and all axioms used
in the subsequent cardinal comparison.

\end{lemma}
\begin{proof}
Abbreviate \(B=\mathfrak G_0\) and \(u=u_{h,m}\). Induction on
~\eqref{hmtfg:base:eq:ch-clausura-transfinita} proves simultaneously that every level is
transitive and that \(\mathfrak G_\alpha\in\mathfrak G_{\alpha+1}\); if all
ordinals smaller than \(\alpha\) are present, their set is definable at the
corresponding level and expresses \(\alpha\). Hence the union contains all
ordinals. Extensionality and Foundation are inherited from the ambient
universe; Empty Set and Infinity belong to \(B\). Once a level containing the
parameters is chosen, the formulas defining \(\{a,b\}\), \(\bigcup a\), and
subsets by Separation express them at a successor level.

For Collection--Replacement, let \(F\) be a functional relation definable in
\(\mathfrak G\) on the set \(a\). The reflection scheme relative to the
definable hierarchy
\((\mathfrak G_\xi,\in,u_{h,m})_{\xi\in\mathrm{Ord}}\), applied in the ambient
metatheory to the finite list of formulas defining \(F\), produces a level
\(\mathfrak G_\theta\) containing \(a\) and the parameters and correct for
those formulas. The ranks of the witnesses in \(F[a]\) are then bounded by
\(\theta\); Separation over \(\mathfrak G_\theta\) forms the image in
\(\mathfrak G_{\theta+1}\). For Power Set, the internal set to be constructed
is not presupposed. Given \(a\in\mathfrak G_\gamma\), consider in the ambient
metatheory the already existing set \(\mathcal P^V(a)\). For every
\(b\in\mathcal P^V(a)\cap\mathfrak G\), let \(r(b)\) be its first
genealogical stage. Since \(\mathfrak G\) is a definable class, Separation in
the ambient metatheory first forms the set
\(\mathcal P^V(a)\cap\mathfrak G\). Replacement \textbf{in the ambient
metatheory}, applied to that set, bounds the ordinals \(r(b)\) by some
\(\theta\). Consequently,
\[
 \mathcal P^{\mathfrak G}(a)
 =\{b\in\mathfrak G_\theta:b\subseteq a\},
\]
and the right-hand side is expressed by Separation in
\(\mathfrak G_{\theta+1}\). It does not introduce into \(\mathfrak G\) the
subsets external to the semantics~\eqref{hmtfg:base:eq:ch-objetos-semanticos}. This is the
noncircular rank-bounding argument.

Finally, each object is assigned its first level of birth and its least
\textbf{derivation name}: a well-founded, finitely branching derivation tree
(and, by~\eqref{hmtfg:base:eq:ch-terminos}, finite for each individual name), labeled by
a level ordinal, a formula code, and the tuple of names of its parameters.
These names are ordered by transfinite recursion: first by the ordinal label,
then by the formula code, and then lexicographically by the already ordered
parameter names. The least name of each value defines a class relation
\(<_{\mathfrak G}\) that compares all objects. It is a definable global
well-order and, by taking the least element of every nonempty set, proves
Choice.
\end{proof}

\input{inserciones_20260930/feigenbaum/antecedentes_comunes/ch_reflexion_testigos.tex}



\subsection{Ordered readout of block cylinders}
\label{hmtfg:publicacion-cilindros-bloque}

For the representative \([b]_3\in\{0,1,2\}\), fix
\[
 \nu(\mathbf b)=\sum_{i=1}^6[b_i]_3\,3^{6-i}\in\{0,\ldots,728\}.
\]
In the HMT ordered completion, a prefix
\(s=(m;\mathbf b_0,\ldots,\mathbf b_{N-1})\), \(m\in\ZZ\), determines
\[
 a_N(s)=m+\sum_{j=0}^{N-1}\nu(\mathbf b_j)\,729^{-j-1},\qquad
 I_N(s)=[a_N(s),a_N(s)+729^{-N})_{\rm HMT}.
\]
The carry preserves the joining of the two representatives of a boundary. If \(C_N(s)=\overline{I_N(s)}\), the preceding relations give
\[
 I_N(s)=\bigsqcup_{\mathbf b\in\FF_3^6}I_{N+1}(s;\mathbf b),
 \qquad \operatorname{diam}C_N(s)=729^{-N}\longrightarrow0.
\]

\begin{proposition}[Coverage of the Archimedean readout]
\label{hmtfg:cobertura-arquimediana}
The map
\[
 P_{\rm ar}:\ZZ\times\Omega_\Gamma^{\rm cal}
       \longrightarrow\mathbb R_{\rm HMT},\qquad
 \{P_{\rm ar}(m,\xi)\}
    =\bigcap_N C_N(m;\beta_N(\rho_{\infty,N}\xi))
\]
is well defined, is surjective and determines
\[
 (\ZZ\times\Omega_\Gamma^{\rm cal})/\!\sim_{P_{\rm ar}}
       \ \cong\ \mathbb R_{\rm HMT}.
\]
\end{proposition}
\begin{proof}
The left endpoints form a Cauchy sequence in the HMT completion, since the intervals are nested and their diameters tend to zero. Their limit belongs to all the closed intervals and is unique. For any element of the completion, first choose the integer interval containing it and then the unique half-open subinterval of each partition containing it. The preceding proposition realizes each of those choices through an effective TPK prolongation. The compatible prefixes define \(\xi\), whose readout is the chosen element. At boundaries, the carry identifies equivalent representations after their construction. By definition and surjectivity, the quotient by equality of readout is in bijection with the completion.
\end{proof}

The construction is also carried out in the internal genealogical universe \(\mathfrak G=\mathfrak G_{\rm HMT}(h,m_\infty)\). Here \(h\) encodes the operator rules and \(m_\infty\) the compatible history; the universe is obtained through initial transitive closure, definability at each successor and union at limit stages:
\[
 G_0=\operatorname{TC}\{\omega,u_{h,m},h,m_\infty\},\qquad
 G_{\beta+1}=\operatorname{Def}_{\Sigma_{\rm HMT}}(G_\beta),\qquad
 G_\lambda=\bigcup_{\beta<\lambda}G_\beta .
\]
The relations of \(\Sigma_{\rm HMT}\) are those of the operations already constructed and their codes. The internal domain is therefore retained in all coverage claims. Denote its completion by \(\mathbb R_{\rm HMT}^{\mathfrak G}\) and the ordered realization within the same universe by \(\mathbb R^{\mathfrak G}\).

\begin{proposition}[Internal ordered chart]
\label{hmtfg:carta-interna}
There exists a unique isomorphism of complete Archimedean ordered fields
\[
 \iota:\mathbb R_{\rm HMT}^{\mathfrak G}
             \xrightarrow{\ \cong\ }\mathbb R^{\mathfrak G}
\]
that fixes the unit. It preserves rational endpoints, positive square roots and limits of convergent sequences.
\end{proposition}
\begin{proof}
The unit fixes the integers and rationals. Each element of the HMT completion is associated with the cut of rationals less than it; completeness of the ordered realization provides the unique element with that cut. Rational density proves injectivity and strict order preservation. The same procedure in the reverse direction proves surjectivity. Rational approximations from above and below preserve addition and multiplication and thereby the field structure. The positive element whose square is \(a>0\) is transported to the positive element whose square is \(\iota(a)\). If a sequence converges, for every positive rational tolerance its final terms lie in the rational neighborhood of the limit; order preservation transports that condition. All these operations are performed within \(\mathfrak G\). Uniqueness follows from preservation of rational cuts.
\end{proof}

\subsection{Descent along fibers and preservation of the selector}
\label{hmtfg:descenso-selector}

% XI REV11, ch_descenso_por_fibras.tex:10--39; prueba completa pertinente.
\begin{proposition}[Descent criterion for an operation]
\label{hmtfg:criterio-descenso}
Let \(q:\Omega\twoheadrightarrow J\) and
\(\star:D\subseteq\Omega^2\to\Omega\), with \(D\) saturated by the fibers of \(q\times q\). There exists a unique operation \(\bar\star\) on \((q\times q)(D)\) such that
\[
 q(x\star y)=q(x)\,\bar\star\,q(y)
\]
if and only if
\[
 q(x)=q(x'),\quad q(y)=q(y')
 \ \Longrightarrow\ q(x\star y)=q(x'\star y')
\]
for pairs in the domain.
\end{proposition}
\begin{proof}
If \(\bar\star\) exists, both results are the evaluation of the same pair of arguments. Conversely, choose antecedents of the visible arguments and define the result as the readout of their operation. Saturation fixes the domain; the condition on fibers makes the result independent of the chosen antecedents. Surjectivity proves uniqueness.
\end{proof}

In an application through \(P_{\rm ar}\), this criterion identifies the precise condition for transporting an operation from complete states to their readout. In the chart \(\iota\), the field structure and limits already satisfy that condition. Next, the analytic profile is constructed from the calendar and this compatibility is used to transport its iteration and roots.

\subsection{Dodecaphase readout and quadratic normalization}
\label{hmtfg:perfil-dodecafasico}

The profile, its phases and its normalization are those already constructed in \eqref{mab:rm:eq:critical-phase-reader}--\eqref{mab:rm:eq:chaos-profile}. Its closed form and moments are proved in \eqref{mab:rm:eq:chaos-closed} and \eqref{mab:rm:eq:chaos-moments24}. The finite cosine sum determines the removable extension at every zero of the closed-form denominator. The angular moment retains oriented real representatives: changing them modulo $2\pi$ changes the moment. The reader and its uniform sector are the explicit starting structure of this realization.

\subsection{Transport of iterates, roots and minima}
\label{hmtfg:transporte-raices}

% Reunión existente: TRANSVERSALIDAD_Y_ESCALADO_LOCAL_FEIGENBAUM.md, sección 14.1.
Within \(\mathbb R_{\rm HMT}^{\mathfrak G}\), the convergent series
\[
 \cos_{\rm HMT}(t)=\sum_{j=0}^{\infty}
                   \frac{(-1)^j t^{2j}}{(2j)!}
\]
defines the corresponding analytic realization. Let
\[
 u_{\rm HMT}(x)=\frac1{12}\sum_{m=1}^{12}
  \left[1-\cos_{\rm HMT}
   \left(\frac{2(3m-1)x}{\sqrt{899}_{\rm HMT}}\right)\right],
 \qquad F_a(x)=1-a\,u_{\rm HMT}(x),
\]
\[
 S_n^{\rm HMT}(a)=F_a^{\,2^n}(0),\qquad
 S_n(\lambda)=f_\lambda^{\,2^n}(0).
\]

\begin{proposition}[Ordered conjugacy of the family and the root selector]
\label{hmtfg:conjugacion-selector}
For the internal arguments of the preceding expressions,
\[
 \iota\circ F_a=f_{\iota(a)}\circ\iota,\qquad
 \iota(S_n^{\rm HMT}(a))=S_n(\iota(a)).
\]
The chart transports exactly the conditions of return, exact period and parameter order. If a set of admissible roots has a minimum, the image of that minimum is the minimum of the image set.
\end{proposition}
\begin{proof}
The chart fixes the rationals, transports the positive square root and preserves limits of the cosine partial sums. Hence, \(\iota(u_{\rm HMT}(x))=u_0(\iota(x))\). Compatibility with addition and multiplication proves the first equality; induction on the number of iterations proves the second. A root is characterized by equality to zero, which is preserved in both directions. The exact period \(2^n\) adds the inequalities \(F_a^{2^j}(0)\ne0\), \(0\le j<n\), also preserved by the injectivity of \(\iota\). If the parameter is required to exceed a preceding center, that inequality is transported by order preservation. Finally, \(a\le b\) for every \(b\) in the set is equivalent to \(\iota(a)\le\iota(b)\) for every element of its image. Surjectivity onto the sets thus defined completes the assertion about the minimum.
\end{proof}

The domain of this transport is the indicated internal continuum. The existence of each minimum and its identification with a specific dynamical branch are results of the criticality construction that follows; the chart preserves those properties once they have been established.

\subsection{Finite spectral representation of the uniform reader}
\label{hmtfg:jacobi-finito}

% Grassmannianos REV07, sections/jacobi_spectral.tex:35--136,
% mecanismo de Hankel, Gram--Schmidt, recurrencia y representación espectral.
% Especialización atómica ya reunida en la sección 14.2 del desarrollo Feigenbaum.
The squared frequencies of the profile define the positive measure
\[
 s_m=k_m^2=\frac{4(3m-1)^2}{899},\qquad
 \nu_0=\frac1{12}\sum_{m=1}^{12}\delta_{s_m},\qquad
 M_r=\int s^r\,d\nu_0(s).
\]
Its twelve nodes are distinct and positive. For \(1\le r\le12\), the Hankel matrix \(H_r=(M_{i+j})_{0\le i,j<r}\) is positive definite. Indeed, for \(p\) of degree less than \(r\),
\[
 \mathbf p^{\mathsf T}H_r\mathbf p
   =\frac1{12}\sum_{m=1}^{12}p(s_m)^2>0
\]
unless \(p=0\), because a polynomial of degree at most eleven that vanishes at twelve distinct nodes is zero. At degree twelve appears the polynomial \(\prod_m(s-s_m)\), with zero norm in this measure. The corresponding spectral representation has dimension twelve.

Let \(p_0,\ldots,p_{11}\) be the monic orthogonal polynomials,
\(h_j=\int p_j^2\,d\nu_0>0\) and
\(\psi_j=p_j/\sqrt{h_j}\). Since \(\nu_0\) is a probability measure,
\(\psi_0=1\). Multiplication by \(s\) in \(L^2(\nu_0)\) is self-adjoint and has a tridiagonal matrix \(J_{12}\) in this basis.

\begin{proposition}[Exact recurrence and spectral intertwining]
\label{hmtfg:jacobi-entrelazamiento}
With
\[
 a_j=\frac{\int s p_j(s)^2\,d\nu_0(s)}{h_j},\qquad
 \beta_j=\frac{h_j}{h_{j-1}}>0\quad(1\le j\le11),
\]
the matrix \(J_{12}\) has diagonal \(a_j\) and subdiagonal
\(\sqrt{\beta_j}\). If
\[
 Q_{mj}=\frac{\psi_j(s_m)}{\sqrt{12}}
 \quad(1\le m\le12,\ 0\le j\le11),\qquad
 D=\operatorname{diag}(s_1,\ldots,s_{12}),
\]
then
\[
 Q^{\mathsf T}Q=I,\qquad DQ=QJ_{12},\qquad
 J_{12}=Q^{\mathsf T}DQ,
\]
and the profile satisfies
\begin{equation}
 u_0(x)=e_0^{\mathsf T}
          \bigl[I-\cos(x\sqrt{J_{12}})\bigr]e_0.
 \label{hmtfg:perfil-espectral}
\end{equation}
\end{proposition}
\begin{proof}
Positivity of the Hankel systems uniquely provides the monic orthogonal polynomials. On expanding \(s p_j\) in the orthogonal basis, the coefficients of \(p_i\), \(i\le j-2\), vanish because
\(\langle sp_j,p_i\rangle=\langle p_j,sp_i\rangle=0\).
The coefficient of \(p_{j+1}\) is one by monicity, that of \(p_j\) is \(a_j\) and that of \(p_{j-1}\) is \(h_j/h_{j-1}\). At the last degree, the same equality is understood in \(L^2(\nu_0)\), where the monic polynomial of degree twelve vanishing at the nodes represents the zero vector. Normalization gives the stated positive subdiagonal.

Orthonormality is equivalent to \(Q^{\mathsf T}Q=I\). The recurrence equality at each node gives \(DQ=QJ_{12}\). Since \(Q\) is square, we also have \(QQ^{\mathsf T}=I\), and obtain the diagonalization of \(J_{12}\). Its positive eigenvalues allow its positive square root to be defined. For every function defined at those twelve nodes,
\[
 e_0^{\mathsf T}\Phi(J_{12})e_0
   =\frac1{12}\sum_{m=1}^{12}\Phi(s_m),
\]
since \(Qe_0=(1,\ldots,1)^{\mathsf T}/\sqrt{12}\).
The choice \(\Phi(s)=1-\cos(x\sqrt{s})\) proves
\eqref{hmtfg:perfil-espectral}.
\end{proof}

The representation preserves exactly the uniform functional, its moments and its profile. Its first coefficients are
\[
 a_0=M_1=2,\qquad
 \beta_1=M_2-M_1^2=\frac{2493348}{808201}.
\]
The complete genealogy remains in the TPK state from which the reader originates; \(J_{12}\) represents its specific spectral readout. The Gram--Schmidt construction and the Jacobi recurrence are applied to this atomic measure, terminating in dimension twelve.

The operator chain used in the following sections is thus fixed: APP cells and TRIT orientation feed the TPK transition; its prolongations produce cylinders and their ordered readout; the twelve oriented windows and the uniform functional produce \(u_0\); the family \(f_\lambda\) determines returns and roots; the internal chart preserves the order of their selectors. The spectral realization describes the same profile through a finite self-adjoint matrix and a cyclic vector.

% Fragmento integrable. Preambulo anfitrion: amsmath, amssymb, longtable, hyperref.
\section{Dodecaphase profile and quadratic return}
\label{hmtfg:fragmento:perfil-dodecafasico-y-retorno-cuadratico}

\subsection{Genealogy and object of the proof}
\label{hmtfg:desarrollo:1}

The initial object is the critical readout of the oriented dodecaphase wheel already constructed in the corpus. APP evaluates addition and multiplication on the two sheets of the digit lattice, preserving quotient and remainder; TRIT determines regime and orientation; TPK transport selects, transports and updates the state with its carry, boundary and memory records. The nonadic connection prolongs that state, with phase return and memory advance. Its readers act on the same discrete structure of the continuum.

The focal reader of the source is
\[
\mathcal P_{\mathrm{crit}}:U_{12}^{\mathrm{sign}}\longrightarrow
\mathbb R/2\pi\mathbb Z,\qquad
m\longmapsto\phi_m=\frac{(3m-1)\pi}{18},\quad 1\le m\le12.
\]
To evaluate the angular moment, the oriented representative of that chart is retained: replacing it with another representative modulo \(2\pi\) would change the moment. The uniform trace fixes \(w_m=1/12\). The focal chain is
\[\begin{aligned}
\text{APP–TRIT–TPK}&\longrightarrow
\text{discrete structure of the continuum}\\ &\longrightarrow U_{12}^{\mathrm{sign}}
\xrightarrow{\mathcal P_{\mathrm{crit}}}(\phi_m,w_m)
\\ &\longrightarrow u_0\longrightarrow f_\lambda
\longrightarrow\text{returns and renormalization}.
\end{aligned}\]

The abbreviated notation retains the complete state and its prolongations as antecedent; the scalar family is a readout of that antecedent. The already fixed reader, with its order, orientation and uniform sector, is taken as an explicit starting structure. The conclusions concern that particular reader; uniqueness of that reader among all possible TPK readouts lies outside these proofs.

The coordinate \(\pi_{\mathrm{HMT}}\) used by the chart belongs to the internal outputs of the corpus. Its factor cancels exactly when normalizing the second moment. The Feigenbaum values enter exclusively in subsequent recognition, never in the definition of the family or the certification of roots.


\subsection{Exact profile and quadratic criticality}
\label{hmtfg:desarrollo:2}

We retain the notation $s_m=3m-1$, $\kappa_0=s_0$, $k_m=2s_m/\sqrt{899}$ and $f_{\lambda,0}=f_\lambda$ of the profile already proved. Its entire extension has the rational series
\[
u_0(z)=\sum_{j\ge1}(-1)^{j+1}
\frac{4^j\sum_m s_m^{2j}}{12\,899^j(2j)!}z^{2j}
=z^2-\frac{715769}{2424603}z^4
+\frac{1353832978}{32695771455}z^6+\cdots.
\]
Thus the finite record of twelve sectors produces an entire profile, with quadratic criticality and scale fixed before iteration.


\subsection{Structural robustness with uniform bounds}
\label{hmtfg:desarrollo:3}

The deformation is defined in \eqref{mab:eq:pesos-dodecafasicos-armonizados}--\eqref{mab:eq:familia-feigenbaum-ponderada-armonizada}. We write $M_{2,\eta}=\sum_mw_m\phi_m^2$ and $\kappa_\eta=s_\eta=\sqrt{2/M_{2,\eta}}$, retaining the same weights, phases and profiles.
The notation is $p_m=p(\phi_m)=12\cos(3\phi_m)-\cos(4\phi_m)$ and $k_{m,\eta}=\kappa_\eta\phi_m$; in the following weighted formulae, $k_m$ abbreviates $k_{m,\eta}$, whereas $k_{m,0}$ denotes the uniform sector.

\textbf{Structural robustness theorem.} For \(|\eta|\le1/40\), the weights are positive and sum to one. The profile is even and entire, satisfies \(u_\eta''(0)=2\) and is strictly increasing on \((0,1]\). For \(0<\lambda\le2/u_\eta(1)\), the map sends \([-1,1]\) into itself, has a unique critical point in that interval and satisfies
\[
S f_{\lambda,\eta}(x)\le-\frac{640}{47647}<0,\qquad0<|x|\le1.
\]

\textbf{Proof.} The sums of cosines \(3\phi_m\) and \(4\phi_m\) vanish through cycles of lengths four and three. Since \(|p_m|\le13\),
\[
\frac9{160}\le w_m\le\frac{53}{480},\qquad
\frac{27}{40}M_{2,0}\le M_{2,\eta}\le\frac{53}{40}M_{2,0}.
\]
It follows that
\[
k_{\max,\eta}^2\le\frac{196000}{24273}<9<\pi^2,\qquad
k_{\min,\eta}^2\ge\frac{640}{47647}.
\]
For \(0<x\le1\), all sines \(\sin(k_{m,\eta}x)\) are positive, whence
\[
u_\eta'(x)=\sum_mw_mk_m\sin(k_mx)>0,\qquad
u_\eta'''(x)=-\sum_mw_mk_m^3\sin(k_mx)<0.
\]
The quotient \(u_\eta'''/u_\eta'\) is the negative of a positive mean of the \(k_m^2\). Invariance of the Schwarzian under affine changes of the value gives
\[
S f_{\lambda,\eta}
=\frac{u_\eta'''}{u_\eta'}-\frac32
\left(\frac{u_\eta''}{u_\eta'}\right)^2
\le-k_{\min,\eta}^2.
\]
Evenness covers \(x<0\), and the image of the interval is
\([1-\lambda u_\eta(1),1]\). \(\square\)

The width \(1/40\) is a choice made to obtain uniform bounds, distinct from a fundamental constant or an optimal threshold. The original sector remains \(\eta=0\).


\subsection{Global quadratic holomorphic coordinate on a strip}
\label{hmtfg:desarrollo:4}

\textbf{Holomorphic factorization theorem.} For \(|\eta|\le1/40\), in
\[
\mathcal S=\{z\in\mathbb C:|\operatorname{Re}z|<\pi/3\}
\]
there exists a holomorphic, odd and univalent function \(b_\eta\), with \(b_\eta'(0)=1\), such that
\[
\boxed{u_\eta(z)=b_\eta(z)^2,\qquad
f_{\lambda,\eta}(z)=1-\lambda b_\eta(z)^2.}
\]

\textbf{Proof.} For \(z=x+iy\) in the right half-strip,
\[
\operatorname{Re}u_\eta'(z)
=\sum_mw_mk_m\sin(k_mx)\cosh(k_my)>0.
\]
The integral of \(u_\eta'\) over the segment between two points, divided by the difference of the endpoints, has positive real part. Convexity of the half-strip proves injectivity there. Evenness gives the corresponding result on the left.

Moreover,
\[
\operatorname{Im}u_\eta(x+iy)
=\sum_mw_m\sin(k_mx)\sinh(k_my)
\]
has the sign of \(y\) for \(x>0\). On the positive real axis \(u_\eta>0\); on the imaginary axis,
\[
u_\eta(iy)=\sum_mw_m(1-\cosh(k_my))
\]
is strictly negative away from zero and decreases strictly with \(|y|\). These separations, injectivity on each half-strip and evenness prove that the fibers of \(u_\eta\) in \(\mathcal S\) are exactly \(\{z,-z\}\), except for the fiber of the origin.

The function \(u_\eta(z)/z^2\), extended with value one at zero, is holomorphic and has no zeros in the simply connected strip. It has a holomorphic logarithm \(L_\eta\) with \(L_\eta(0)=0\), which is even by uniqueness of the normalization. Define
\[
b_\eta(z)=z\exp\!\left(\frac12L_\eta(z)\right).
\]
Then \(b_\eta^2=u_\eta\), \(b_\eta\) is odd and \(b_\eta'(0)=1\). The equality \(b_\eta(z)=b_\eta(w)\) forces \(w=z\) or \(w=-z\). The second case gives equality only at the origin. This proves its univalence. \(\square\)

The twelve phases thus determine an explicit conformal deformation of the quadratic fold. The complete dynamical conjugacy preserves that deformation:
\[
b_\eta\circ f_{\lambda,\eta}\circ b_\eta^{-1}(w)
=b_\eta(1-\lambda w^2),
\]
on their composition domains. The factorization proved and a conjugacy to the entire quadratic family are distinct assertions.


\subsection{Return, rescaling and memory}
\label{hmtfg:desarrollo:5}

Let \(a=-f(1)>0\), \(H_a(x)=-ax\), \(J=H_a([-1,1])\). When \(J\subset[-1,1]\) is admissible for the return, \(f^2(J)\subset J\) and the compositions are defined,
\[
\mathcal Rf=H_a^{-1}f^2H_a,\qquad
\mathcal Rf(x)=-a^{-1}f(f(-ax)).
\]
If in addition \(f(J)\subset(0,1]\), a unique central quadratic maximum is preserved:
\[
(\mathcal Rf)(0)=1,\qquad
(\mathcal Rf)''(0)=-a f'(1)f''(0)<0.
\]
Composition of Schwarzians gives, at regular points,
\[
S(\mathcal Rf)(x)=a^2
\left[(Sf)(f(-ax))f'(-ax)^2+(Sf)(-ax)\right]<0.
\]
Domain inclusions are part of the hypotheses and are checked at every scale.

Let \(v_\eta=(u_\eta|_{[0,1]})^{-1}\). The inverse branches are
\[
\beta_\sigma(y)=\sigma v_\eta((1-y)/\lambda),\quad
\sigma\in\{-1,+1\},
\]
with the critical record \(\beta_0(1)=0\). The step
\[
(x,w)\longmapsto(f(x),w\mathbin{\|}\sigma(x))
\]
is inverted on its image by reading the last sheet and applying \(\beta_\sigma\). The word preserves branch decisions; the previous value is reconstructed through the map equation.

The return groups two transitions. From the endpoint \(x'\) and its sheets \((\sigma_0,\sigma_1)\) one recovers
\[
x=H_a^{-1}\beta_{\sigma_0}
\!\left(\beta_{\sigma_1}(H_ax')\right).
\]
The intermediate points are also reconstructed. The APP–TRIT–TPK records of sheet, carry, route, boundary and memory are concatenated while preserving their identity. The sign of the analytic branch is an additional record: it does not replace those data with a binary label.

If \(C_N\) groups histories of length \(2N\) into pairs, the restrictions satisfy
\[
\pi^{\mathrm{ret}}_{N,M}C_N
=C_M\pi^f_{2N,2M},\qquad M\le N.
\]
Both sides remove the same terminal pairs and preserve the same inverse branches. Compatible transport of finite histories is thus proved.


\subsection{Accumulated scales and derivative of the operator}
\label{hmtfg:desarrollo:6}

Let \(c_n=f^{2^n}(0)\), \(c_0=1\) and \(H_n(x)=c_nx\). As long as the compositions are admissible and \(c_n\ne0\),
\[
\boxed{\mathcal R^nf=H_n^{-1}f^{2^n}H_n},\qquad
(\mathcal R^nf)(1)=\frac{c_{n+1}}{c_n}=-a_n.
\]
Induction follows from \(H_nH_{a_n}=H_{n+1}\). The convention \(a_n>0\) requires sign alternation of \(c_n\). At a superstable root of period \(2^N\), \(c_N=0\): that normalization becomes singular and the preceding levels or a controlled limit must be used.

For \(z=-ax\), \(w=f(z)\) and a perturbation \(h\), the complete derivative is
\[
\boxed{
D\mathcal R_f[h](x)=
-\frac{h(w)+f'(w)h(z)}a+
\frac{h(1)}a
\left[\mathcal Rf(x)-xf'(w)f'(z)\right].
}
\]
It is differentiated using \(\dot a=-h(1)\), \(\dot z=xh(1)\),
\(\dot w=h(z)+xf'(z)h(1)\) and the variation of the denominator. If \(h(0)=0\), this gives \(D\mathcal R_f[h](0)=0\). The initial HMT parameter direction is \(h=-u_\eta\).

This formula includes scale variation; freezing \(a\) would give a different linearized operator.


\subsection{Certified roots and itineraries through period 256}
\label{hmtfg:desarrollo:7}

Let \(F_n(\lambda)=f_{\lambda,0}^{2^n}(0)\). Compute
\[
x_0=p_0=0,\qquad x_{j+1}=1-\lambda u_0(x_j),\qquad
p_{j+1}=-u_0(x_j)-\lambda u_0'(x_j)p_j.
\]
Then \(F_n=x_{2^n}\) and \(F_n'=p_{2^n}\).

The attached program uses rational intervals with outward integer rounding, square roots enclosed by integer square roots, and sine/cosine with an explicit Taylor remainder. The exact bound \(|u_0''|\le2\) allows centered evaluation with a quadratic remainder.

The historical approximations only propose boxes of radius \(10^{-42}\). Opposite signs at their endpoints, separation of \(F_n'\) from zero throughout the box, and exclusion of returns with periods that properly divide \(2^n\) are certified. The intermediate value theorem and monotonicity prove local existence and uniqueness; exclusion of divisors proves the least period.

\begingroup\small
\begin{longtable}{p{.12\linewidth}p{.17\linewidth}p{.62\linewidth}}
\hline
Level & Least period & Box center, rounded \\
\hline
1 & 2 & 1.345913990471832792210949 \\
2 & 4 & 1.763379787846910127290794 \\
3 & 8 & 1.850413796950136464530567 \\
4 & 16 & 1.868979756606798125150574 \\
5 & 32 & 1.872955034435659740004205 \\
6 & 64 & 1.873806367467030119412262 \\
7 & 128 & 1.873988695221365362307169 \\
8 & 256 & 1.874027744154450672897008 \\
\hline\end{longtable}
\endgroup

The first root has the exact expression \(\lambda_1=1/u_0(1)\). The complete boxes and their checks are in the JSON certificate.

The 502 preclosure states of these eight levels, all separated from zero, and their complete sign words have been enclosed. In particular,
\[
\operatorname{sign}f_{\lambda_n,0}^{2^j}(0)=(-1)^j,\qquad0\le j<n.
\]
The finite orientations of the rescalings preceding closure are established. Inclusions of restrictive intervals at every scale constitute an additional check.

For \(\delta_{\mathrm F,n}=(\lambda_{n-1}-\lambda_{n-2})/(\lambda_n-\lambda_{n-1})\), an abbreviated outer decimal interval is
\[
\boxed{\begin{gathered}
4.66921218915020415455609621588966392804<\delta_{\mathrm F,8},\\
\delta_{\mathrm F,8}<4.66921218915020415455609621588966392863.
\end{gathered}}
\]
The width of the complete interval is less than \(5.808\cdot10^{-37}\). It measures the numerical error of the finite ratio \(\delta_{\mathrm F,8}\), distinct from the error in approximating the universal limit.


\subsection{Analytic persistence of the deformed roots}
\label{hmtfg:desarrollo:8}

For each certified level, \(F_n'(\lambda_n)\ne0\). The implicit function theorem produces a local analytic branch \(\lambda_n(\eta)\) with
\[
f_{\lambda_n(\eta),\eta}^{2^n}(0)=0.
\]
By continuity, the proper returns remain separated from zero for \(\eta\) sufficiently small; the least period and itinerary are preserved. The existence of those local neighborhoods is proved. Their common width is not yet identified with the entire structural strip \(1/40\).

With \(M=M_{2,\eta}\), the profile response is
\[
v_\eta(x)=\partial_\eta u_\eta(x)
=\sum_m\frac{p_m}{12}(1-\cos(k_{m,\eta}x))
-\frac{M'}{2M}\,x u_\eta'(x).
\]
The second term preserves normalization of the second moment. If \(q_j=\partial_\eta x_j\) at fixed \(\lambda\),
\[
q_0=0,\qquad q_{j+1}=-\lambda v_\eta(x_j)-\lambda u_\eta'(x_j)q_j,
\qquad
\boxed{\lambda_n'(\eta)=-q_{2^n}/p_{2^n}}.
\]
The weights–scale–response–superstable-parameter chain is thus explicit. This finite nondegeneracy is distinct from asymptotic transversality to the stable leaf of the universal fixed point.


\subsection{Analytic tails and error propagation}
\label{hmtfg:desarrollo:9}

For \(\eta=0\), let
\[
\mu_{2j}=\frac1{12}\sum_mk_m^{2j},\quad
p_N(z)=\sum_{j=1}^N(-1)^{j+1}\frac{\mu_{2j}}{(2j)!}z^{2j},
\quad\Omega=\frac{70}{\sqrt{899}}.
\]
If \(q_N=\Omega^2r^2/[(2N+3)(2N+4)]<1\),
\[
\|u_0-p_N\|_r\le
\frac{\mu_{2N+2}r^{2N+2}}{(2N+2)!(1-q_N)}.
\]
The inequality \(\mu_{2j+2}\le\Omega^2\mu_{2j}\) bounds the tail by a geometric series. For \(0\le s\le2N+2\),
\[
\|(u_0-p_N)^{(s)}\|_r\le
\frac{\mu_{2N+2}r^{2N+2-s}}{(2N+2-s)!}
\left(1-\frac{\Omega^2r^2}{(2N+3-s)(2N+4-s)}\right)^{-1},
\]
when the denominator is positive.

To transport errors through a return, its domains are also controlled: if \(\|f-g\|\le\varepsilon\), \(a_f,a_g\ge a_*>0\), the compositions lie in a common convex domain and \(\|f'\|\le L,\ \|g\|\le M\) there, then
\[
\|\mathcal Rf-\mathcal Rg\|_r
\le\varepsilon\left[\frac{1+L+L^2r}{a_*}+\frac{M}{a_*^2}\right].
\]
The proof adds the outer perturbation, the inner perturbation, the change of argument due to \(|a_f-a_g|\le\varepsilon\) and the variation of the denominator. These bounds prevent reuse of the initial cosine sum as if all renormalized profiles retained that same form.

% Fragmento integrable. Preambulo anfitrion: amsmath, amssymb, longtable, hyperref.
\section{Nonadic depth and persistence of the itinerary}
\label{hmtfg:fragmento:profundidad-nonadica-y-persistencia-del-itinerario}

\subsection{Outer growth and inner resolution}
\label{hmtfg:prolongacion:2}

The refinement law of article XI, with integer base \(B\ge2\), is
\[
L_c(x,y)=(Bx-c,By+c),\qquad 0\le c<B.
\]
For a word \(c_1\ldots c_m\), its record is
\[
C_m=\sum_{j=1}^m c_jB^{m-j},\qquad
(x_m,y_m)=(B^mx_0-C_m,B^my_0+C_m).
\]
The outer sum is \(B^mM_0\), with \(M_0=x_0+y_0\). The normalized reading of the record belongs to the cylinder
\[
I_m(C_m)=\left[\frac{C_m}{B^m},\frac{C_m+1}{B^m}\right],
\qquad \operatorname{diam}I_m=B^{-m}.
\]
The integer pair and the scale recover the quotients, remainders and leading zeros of the word. Endpoints with two expansions retain their boundary mark. Support growth and readout refinement are thus coordinated without identifying distinct histories solely by their limit image.

The composition is exact:
\[
L_D^{[b]}\circ L_C^{[a]}=L_{B^bC+D}^{[a+b]}.
\]
By associativity, grouping a history of eighteen events into nine blocks of two or two blocks of nine produces the same record. At return level \(2^N\), the common horizon is \(9\,2^N\). The proof is an equality of composition on the same ordered history; it retains its carries and its recoverable subblocks.

The restriction of states and the prolongation of their readers are expressed through the following limits:
\[
X_\infty=\varprojlim X_m,\qquad
r^*a=a\circ r,\qquad
r^*e_x=\sum_{r(y)=x}e_y,
\qquad \mathfrak H_{\rm cyl}=\varinjlim\mathfrak H_m.
\]
States restrict; readers extend. Pullback preserves products, the unit and idempotents. An exactly compatible family induces a reader on the direct limit. Approximants that are compatible only up to an error additionally require a convergence estimate in their completion. The next section provides that estimate for the critical return.

\subsubsection{Capacity and vacancies of the record}

The capacity reader of XI compares blocks of \(729\) and \(1000\) possibilities. With \(\rho=\log_{1000}729=\log_{10}9\), it retains
\[
\mathcal C_k=\lfloor(k+1)\rho\rfloor,\qquad
\mathcal C_k-\mathcal C_{k-1}=1-z_k,
\]
\[
\mathcal C_b-\mathcal C_a+\sum_{k=a+1}^b z_k=b-a,
\qquad T_{\rm cap}(a)=\lceil a/\rho\rceil-1\quad(a\ge1).
\]
Each vacancy records a transition during which capacity is maintained; chronological length is conserved in the balance. Capacity phase and chronological phase are distinct readers.

To store a prefix of \(m\) base-nine digits, a sufficient decimal capacity is
\[
a(m)=\left\lceil\frac{m\rho}{3}\right\rceil.
\]
The quotient and remainder allow the integer \(C_m\) to be converted between bases, keeping the depth \(m\) separate. The following hold
\[
9^m\le1000^{a(m)}\le729^{T_{\rm cap}(a(m))+1}.
\]
These inequalities are also decided by integer powers, without rounding logarithms. The readout budget of the next section is thus translated into record capacity while preserving units. This count sizes the storage; the survival rules additionally determine which native histories are admissible.


\subsection{Explicit depth budget for each return}
\label{hmtfg:prolongacion:3}

\textbf{Theorem.} With \(|\eta|\le1/40\) fixed, normalize
\[
t=\lambda u_\eta(1)/2\in[0,1],\qquad z=(x+1)/2,
\]
\[
F_{t,\eta}(z)=1-t\frac{u_\eta(2z-1)}{u_\eta(1)},\qquad z\in[0,1].
\]
For every integer \(q\ge1\) and two parameters \(t,s\),
\[
\left|F_{t,\eta}^{q}(1/2)-F_{s,\eta}^{q}(1/2)\right|
\le S_q|t-s|,
\qquad S_q=\sum_{j=0}^{q-1}(6\sqrt2)^j<9^q.
\]
Consequently, a base-nine parameter cylinder of depth \(m=2^N+r\) encloses the readout of return \(q=2^N\) in an interval of diameter less than \(9^{-r}\).

\textbf{Proof.} The function \((1-\cos v)/v^2\) decreases for \(0<v<\pi\): the sign of its derivative is that of \(v\sin v-2(1-\cos v)\), which is negative because \(\tan(v/2)>v/2\). By the second moment,
\[
u_\eta(1)\ge\frac{1-\cos3}{9}\sum_jw_jk_j^2
=\frac{2(1-\cos3)}9>\frac13.
\]
The last inequality is verified, for example, by bounding \(\cos3\) by its alternating Taylor sum through degree eight, which is less than \(-1/2\). Moreover, Cauchy–Schwarz gives \(|u_\eta'(x)|\le\sqrt2\) on the real line. It follows that
\[
\operatorname{Lip}_zF_{t,\eta}\le6\sqrt2<9,\qquad
\operatorname{Lip}_tF_{t,\eta}\le1.
\]
If \(d_j\) is the difference between the two orbits, then \(d_0=0\) and \(d_{j+1}\le6\sqrt2\,d_j+|t-s|\). Induction proves the geometric sum. For \(|t-s|\le9^{-m}\), the diameter is less than \(9^{q-m}\). Substituting \(q=2^N\) and \(m=q+r\) completes the proof. \(\square\)

The bound controls parameter uncertainty with \(\eta\) fixed. An implementation separately adds and propagates cosine-evaluation and rounding errors; the preceding rational certificate provides these operations. The bound is sufficient and conservative, with a cost that grows with the horizon.

\textbf{Limit reader.} Realize \(t\) from \((x_0,y_0)=(1,0)\), \(B=9\), so that
\[
(x_m,y_m)=(9^m-C_m,C_m),\qquad t_m=C_m/9^m.
\]
The readers \(E_{q,m}=F_{t_m,\eta}^{q}(1/2)\), brought to a common level, satisfy
\[
\|E_{q,m+a}-r^*E_{q,m}\|_\infty\le S_q9^{-m}.
\]
They converge uniformly to a continuous reader \(E_q\) in the completion of the cylinder algebra. The set-theoretic version
\[
J_{q,m}(C)=\{F_{t,\eta}^{q}(1/2):t\in I_m(C)\}
\]
is exactly nested and its diameters tend to zero. This result gives the passage to the limit of the \textbf{fixed-horizon readout}, with depth control at any requested horizon.


\subsection{Doubling word and nonadic reader of chronology}
\label{hmtfg:prolongacion:4}

The eight itineraries certified in the preceding work obey
\[
W_1=+,\qquad W_{n+1}=W_n\,(-1)^n\,W_n.
\]
Its unique prefix continuation is
\[
\tau_j=(-1)^{v_2(j)},\qquad j\ge1;
\qquad \tau_{2j}=-\tau_j,\quad\tau_{2j+1}=+.
\]
The proof uses \(v_2(2^n+r)=v_2(r)\) for \(0<r<2^n\). It is a recurrence at every depth, and agreement with the eight computed returns constitutes its finite check in the family.

In the parity-lexicographic order of a map with a critical maximum, the orientation factor preceding position \(k\) is \(\prod_{j<k}(-\tau_j)\). The word \(\tau\) is strictly greater than any of its positive shifts. Indeed, the first disagreement with an odd shift occurs at \(k=1\) or \(2\); an even shift halves the comparison and doubles \(k\). The first disagreement is therefore at \(k=2^r\). There
\[
\prod_{j<2^r}(-\tau_j)=(-1)^r=\tau_{2^r},
\]
and the oriented difference with the opposite sign is positive. This argument also proves aperiodicity.

The chronological record retains
\[
K=\rho_9(K)+9q_9(K).
\]
Its dyadic reader is
\[
\ell_n(B_K)=[\rho_9(K)+9q_9(K)]_{2^n}=[K]_{2^n}.
\]
On the chart of histories in which \(\Gamma_9 B_K=B_{K+9}\),
\[
\ell_n(\Gamma_9B_K)=\ell_n(B_K)+9\pmod{2^n},\qquad
\pi_{n+1,n}\ell_{n+1}=\ell_n.
\]
Since nine is odd, advancement by nine visits the \(2^n\) positions. Its conjugacy with unit advancement is \(j\mapsto9j\). For \(0<j<2^n\),
\[
v_2(9j\bmod2^n)=v_2(j),
\]
so that nonadic resampling also preserves the oriented word in each certified cycle.

This reader uses remainder \textbf{and} quotient. The phases of \(K=1\) and \(K=10\) agree modulo nine, while their dyadic signs are opposite. Memory distinguishes the two histories. The arithmetic property of resampling holds for every odd multiplier; the HMT calendar fixes the choice of nine here.


\subsection{Compatible bilateral memory at infinite depth}
\label{hmtfg:prolongacion:5}

Let \(\mathcal H_n=\ell^2(\mathbb Z/2^n\mathbb Z)\),
\[
C_ne_j=e_{j+1},\qquad
J_ne_j=(e_j+e_{j+2^n})/\sqrt2.
\]
A check on the basis gives \(J_n^*J_n=I\) and \(C_{n+1}J_n=J_nC_n\), including the boundary terms.

Article XI provides the projectors
\[
P_0=\frac19\begin{pmatrix}8&-\sqrt8\\-\sqrt8&1\end{pmatrix},\qquad
P_1=\frac19\begin{pmatrix}1&\sqrt8\\\sqrt8&8\end{pmatrix}.
\]
They are orthogonal and complementary. Therefore
\[
\mathbb U_n=P_0\otimes I+P_1\otimes C_n
=\begin{pmatrix}T_n&D_n\\D_n&R_n\end{pmatrix}
\]
is unitary, with
\[
T_n=(8I+C_n)/9,\quad D_n=\sqrt8(C_n-I)/9,\quad R_n=(I+8C_n)/9.
\]
The naturality already proved implies
\[
\mathbb U_{n+1}(I_2\otimes J_n)=(I_2\otimes J_n)\mathbb U_n,
\quad \mathbb U_n^a=P_0\otimes I+P_1\otimes C_n^a
\quad(a\in\mathbb Z).
\]
The Hilbert limit is identified with \(L^2(\mathbb Z_2,\mu_{\rm Haar})\), taking \(e_j^{(n)}\) to \(2^{n/2}\mathbf1_{j+2^n\mathbb Z_2}\). The identities extend by density and boundedness. For every vector \(v\) and every integer time \(a\),
\[
\boxed{\|T_{\infty,a}v\|^2+\|D_{\infty,a}v\|^2=\|v\|^2,}
\]
\[
\boxed{v=T_{\infty,a}^*(T_{\infty,a}v)+D_{\infty,a}^*(D_{\infty,a}v).}
\]
Here \(T_{\infty,a}=(8I+C_\infty^a)/9\). Bilateral composition preserves the archive between steps; composing only \(T_n\) describes a different operation.

The permutation \(S_ne_j=e_{9j\bmod2^n}\) satisfies
\[
S_{n+1}J_n=J_nS_n,\qquad S_nC_nS_n^{-1}=C_n^9.
\]
Thus one obtains in the limit
\[
S_\infty C_\infty S_\infty^{-1}=C_\infty^9,
\quad (I_2\otimes S_\infty)\mathbb U_\infty
(I_2\otimes S_\infty^{-1})=\mathbb U_\infty^9.
\]
The temporal representation is faithful: \(C_\infty^a=C_\infty^b\) forces \(2^n\mid a-b\) for every \(n\), whence \(a=b\). Particular preparations may preserve temporal symmetries. The operator identity preserves this distinction.

The projective point \(([K]_{2^n})_n\) and a vector of \(L^2\) have different types: a point history induces a Dirac measure, whereas normalizable amplitudes belong to the Hilbert space. The result constructs an equivariant reader of chronology; the remaining native coordinates of the TPK accompany the complete record.


\subsection{Existence and separation of asymptotic obligations}
\label{hmtfg:prolongacion:6}

\subsubsection{Realization of the infinite word throughout the strip}
\label{hmtfg:prolongacion:6.1}

\textbf{Existence theorem.} For each \(|\eta|\le1/40\) there exists at least one
\[
\lambda_\infty(\eta)\in
\left[\frac1{2u_\eta(1)},\frac2{u_\eta(1)}\right]
\]
such that
\[
\boxed{\operatorname{sign}f_{\lambda_\infty(\eta),\eta}^{j}(0)
=(-1)^{v_2(j)}\quad\text{for every }j\ge1.}
\]
The conclusion realizes the infinite doubling itinerary within the fixed family. The existence quantifier remains separate from a unique selection and from the law of distances between superstable parameters.

\textbf{Proof.} We fix \(\eta\). We write \(K_\lambda\) for the sign word of the critical orbit; when it first returns to zero, the word ends at that zero. We use the unimodal order with a maximum, whose factor before comparing symbol \(j\) is the product of the signs \(-K_i\), \(i<j\).

At the left endpoint the image of the entire interval is in \([1/2,1]\); hence \(K_-=+++\cdots\). At the right endpoint the orbit is \(0\mapsto1\mapsto-1\mapsto-1\), whence \(K_+=+---\cdots\). Comparing the second and third symbols, respectively, gives
\[
K_-<\tau<K_+.
\]

If \(K_{\lambda_0}\) is infinite and distinct from \(\tau\), the first difference occurs at a finite position. Continuity of the corresponding iterates makes the side of the comparison locally constant. It remains to examine a superstable parameter of first period \(p\), with word \(W0\), where \(|W|=p-1\).

For \(\lambda\) near \(\lambda_0\), let \(g_\lambda=f_{\lambda,\eta}^p\) and \(a=g_\lambda(0)\). We have \(g_\lambda'(0)=0\), and a uniform local bound on the second derivative gives
\[
g_\lambda(x)=a+O(x^2).
\]
For \(|a|\) sufficiently small and nonzero, \(g_\lambda\) maps \([-2|a|,2|a|]\) into \([a-|a|/2,a+|a|/2]\). The returns retain the sign of \(a\), and the \(p-1\) intermediate positions retain \(W\). The neighboring words are therefore \((W+)^\infty\) or \((W-)^\infty\); parameters with \(a=0\) retain \(W0\).

If \(\tau\) differs from \(W\) before position \(p\), that difference fixes the same side for all three words. If it shares \(W\), put
\[
s=\tau_p,\qquad e=\prod_{j<p}(-\tau_j),\qquad A=Ws.
\]
The sign of the comparison of \(\tau\) with \(W0\) and with \((W,-s)^\infty\) is \(es\). To compare it with \(A^\infty\), let \(q\) be the first position at which \(\tau_{p+q}\ne\tau_q\). It exists by aperiodicity. Up to position \(p+q-1\), both words agree, and the orientation sign factors as \(O_pO_{q-1}\), with \(O_p=-es\). The strict maximality already proved gives
\[
O_{q-1}(\tau_q-\tau_{p+q})>0.
\]
Consequently, the comparison of \(\tau\) with \(A^\infty\) has sign \(-O_p=es\). The three local possibilities lie on the same side of \(\tau\).

If no parameter realized \(\tau\), the sets \(\{\lambda:K_\lambda<\tau\}\) and \(\{\lambda:K_\lambda>\tau\}\) would be open, disjoint, nonempty and would cover the parameter interval. Connectedness of the interval excludes that partition. The stated parameter therefore exists. \(\square\)

This itinerary-continuity argument applies after the HMT generator and its analytic family have been fixed; its hypotheses are verified directly here. The construction preserves the infinite word without introducing a parameter of the quadratic family or the universal value as an input.

\subsubsection{Margins preventing loss of the itinerary when taking limits}
\label{hmtfg:prolongacion:6.2}

Write \(c_p=f_{\lambda,\eta}^p(0)\). A sufficient uniform rational bound is
\[
|f'|,|f''|\le16,\qquad
B_p:=16^p\frac{16^p-1}{15}\ge\|(f^p)''\|_{[-1,1]}.
\]
Indeed, \(1-\cos v\ge v^2/2-v^4/24\ge v^2/8\) for \(|v|\le3\), so that \(u_\eta(1)\ge1/4\), \(\lambda\le8\) and both derivatives are bounded by \(16\), using the moments. The chain rule proves the bound \(B_p\) by induction.

The word \(\tau\) satisfies \(\tau_{2p}=-\tau_p\) for every \(p\). For any of its realizations, \(c_{2p}\) and \(c_p\) have opposite signs. Since \((f^p)'(0)=0\), Taylor gives
\[
|c_{2p}-c_p|\le\tfrac12B_p|c_p|^2.
\]
The left-hand side is greater than \(|c_p|\), and therefore
\[
\boxed{|c_p|>2/B_p>0.}
\]
For each fixed horizon there is thus a uniform margin preventing a sequence of realizations from losing its sign upon convergence. In the compact coordinates \((\eta,b)\), \(b=\lambda u_\eta(1)\in[1/2,2]\), the set of realizations of \(\tau\) is closed and compact, and its projection onto the entire \(\eta\) strip is surjective.

The same argument applies to a prefix tree: if all prefixes of \(\tau\) are realized, a subsequence in the compact set preserves every sign, because sufficiently long prefixes simultaneously include positions \(p\) and \(2p\). The existence theorem now provides the nonemptiness that this argument requires.

\subsubsection{Invariant return intervals at every level}
\label{hmtfg:prolongacion:6.3}

\textbf{Nested returns theorem.} For any of the realizations of the existence theorem in section \ref{hmtfg:prolongacion:6.1}, let \(c_j=f^j(0)\) and
\[
J_n=\operatorname{conv}\{c_{2^n},c_{2^{n+1}}\},\qquad n\ge1.
\]
Then \(0\in\operatorname{int}J_n\), \(J_{n+1}\subset J_n\),
\[
f^{2^n}(J_n)=J_n,
\]
and the \(2^n\) images \(J_n,f(J_n),\ldots,f^{2^n-1}(J_n)\) are disjoint. The normalized returns are unimodal with a maximum, with critical value one and the same word \(\tau\). These intervals are invariant return cores; a normalization requiring periodic endpoints may enlarge the core to the corresponding maximal restrictive interval and must specify that step.

\textbf{Proof of the first return.} Let \(F\) be a unimodal map with a maximum, with \(F(0)=1\) and word \(\tau\), and \(d_j=F^j(0)\). Its first signs are \(+,-,+,+,+,-,+\). The interval \(I=[d_2,1]\) is invariant: \(F(d_2)=d_3>0\), \(F(1)=d_2<0\), and the maximum is one. Put \(J=[d_2,d_4]\). Since \(F\) decreases on positive arguments and
\[
F(d_3)=d_4>0> d_6=F(d_5),
\]
we have \(d_3<d_5\), and therefore \(F(J)=[d_3,1]\).

Moreover, \(d_3>d_4\). To check this, \(F^2\) is strictly increasing between the two positive points \(d_3,d_4\), because their images \(d_4,d_5\) are positive. Its values are \(F^2(d_3)=d_5>0>d_6=F^2(d_4)\), which forces that order. Thus \(J\) and \(F(J)\) are separated by an open interval, and
\[
F^2(J)=F([d_3,1])=[d_2,d_4]=J.
\]
The map \(F^2|_J\) has a unique critical minimum: the image \(F(J)\) remains positive. The conjugacy \(h(x)=d_2x\) reverses orientation and normalizes the return as a maximum with value one. Its itinerary is
\[
\operatorname{sign}\frac{d_{2j}}{d_2}
=-\tau_{2j}=\tau_j.
\]
The same argument can be repeated at any depth.

\textbf{Inductive disjointness.} Suppose the \(p=2^n\) images of the parent \(J_n\) are disjoint, and let \(g=f^p\). The first return applied to the normalized map produces \(B=J_{n+1}\), \(A=g(B)\), with \(A\cap B=\varnothing\), \(g(B)=A\), \(g(A)=B\). For \(0\le j<p\), an intersection of \(f^j(A)\) and \(f^j(B)\) would, upon applying \(f^{p-j}\), produce an intersection of \(B\) and \(A\). Those two images are therefore disjoint. Images corresponding to distinct \(j\) belong to disjoint images of the parent. This proves the \(2p\) disjoint images and the return \(f^{2p}(B)=B\). The endpoint formula results from composing the normalizations \(x\mapsto c_{2^n}x\). \(\square\)

The compact sets
\[
\Lambda_n=\bigcup_{0\le j<2^n}f^j(J_n)
\]
are nonempty and nested. Their intersection \(\Lambda_\infty\) has a continuous reader onto the dyadic odometer: each point is assigned the label of the unique component it occupies at each level. Every compatible sequence of labels has a preimage by compactness, and advancement by \(f\) adds one to the labels. This map is a surjective semiconjugacy. Promoting it to a pointwise conjugacy requires proving that the component diameters tend to zero; that metric promotion remains separate.

% Fragmento integrable. Preambulo anfitrion: amsmath, amssymb, longtable, hyperref.
\section{Compatibility of readers and recovery of the parameter}
\label{hmtfg:fragmento:compatibilidad-de-lectores-y-restitucion-del-parametro}

\subsection{Compatibility of readers, minima and recovery of history}
\label{hmtfg:escalado:13}

The prolongation operations and ordered readers preserve the information required for root selection. The following identities specify their domains and composition.

\subsubsection{Recovered forward rule}
\label{hmtfg:escalado:13.1}

The deterministic update starts from the R36 boundary state with character \(\chi\) and internal remainder. The update is
\begin{equation}
\mathcal U_K^\chi=
\operatorname{Upd}_K^\chi\circ
\operatorname{Tra}_K^\chi\circ
\operatorname{Sel}_K^\chi,
\qquad
N_{K+1}=729N_K+d_K^\chi.
\label{hmtfg:escalado:eq:26}
\end{equation}
The digit and the update depend on the preceding state. Two forward sections with the same complete R36 state coincide by induction. The subrelation \(\operatorname{Ext}^{\chi,\rightarrow}\) is the graph of that update; the total relation \(\operatorname{Ext}\) admits multiple continuations depending on the data and characters.

Compatible prolongations transport the regional readers and the joint signature over the same prefixes.

These rules transport orientation, remainder, quotient, carry, boundary and memory. Their application to the cascade must also preserve ordered selection of the critical reader’s parameter.

\subsubsection{Independence of the antecedent constants from the free parameter}
\label{hmtfg:escalado:13.2}

Let \(B\) be the HMT state providing the wheel and its antecedent constants. Write \(\mathcal C(B)\) for the set of those readouts and \(u_B\) for the critical profile. In the domain of the critical reader, the subsequent construction is
\begin{equation}
(B,\lambda)\longmapsto f_{B,\lambda}(x)=1-\lambda u_B(x),
\qquad 0<\lambda\le 2/u_B(1).
\label{hmtfg:escalado:eq:27}
\end{equation}
For the same \(B\), the reading of the preceding constants factors through projection onto the first component:
\begin{equation}
\widetilde{\mathcal C}(B,\lambda)
=\mathcal C\circ\operatorname{pr}_B(B,\lambda)
=\mathcal C(B).
\label{hmtfg:escalado:eq:28}
\end{equation}
Consequently, two distinct values of \(\lambda\) preserve the same antecedent readouts. The equality is a property of this causal construction: the profile is fixed before \(\lambda\) is varied, and the factor \(\pi_{\mathrm{HMT}}\) cancels in its quadratic normalization.

\textbf{Lemma.} Suppose a condition \(A\) is evaluated solely on those antecedent readouts. For every fixed \(B\),
\begin{equation}
\{\lambda\in\Lambda:A(\widetilde{\mathcal C}(B,\lambda))\}
=
\begin{cases}
\Lambda,&A(\mathcal C(B))\text{ holds},\\
\varnothing,&A(\mathcal C(B))\text{ fails}.
\end{cases}
\label{hmtfg:escalado:eq:29}
\end{equation}
\textbf{Proof.} Substituting \eqref{hmtfg:escalado:eq:28} makes the predicate independent of \(\lambda\). Therefore all parameters in the domain satisfy it, or none do. This includes comparison of the antecedent constants with metrological intervals.

Uniqueness of the section \(B\) and the subsequent restriction \(f_{B,\lambda}^{2^n}(0)=0\) produce a solution fiber in \(\lambda\). The first-root rule acts on this fiber through the return and its itinerary. Selection and its prolongation are proved in the following results.

This result maintains the HMT order: metrology compares previously constructed outputs, and root selection is proved within the reader that produces it.

\subsubsection{Lemma on ordered transport of the selector}
\label{hmtfg:escalado:13.3}

Let \(P_n\) be an ordered set of genealogical candidates with minimum \(p_n^*\). Let \(Z_n\) be the corpus’s original set of candidate roots, with the real order, and let \(L_n:P_n\to Z_n\) be a monotone reader. Its image is said to be \textbf{coinitial} when
\begin{equation}
\forall z\in Z_n\quad\exists p\in P_n:
L_n(p)\le z.
\label{hmtfg:escalado:eq:30}
\end{equation}
Then
\begin{equation}
\boxed{L_n(p_n^*)=\min Z_n
\quad\Longleftrightarrow\quad
L_n(P_n)\text{ is coinitial in }Z_n.}
\label{hmtfg:escalado:eq:31}
\end{equation}
\textbf{Proof.} If the image of the minimum is the global minimum, take \(p=p_n^*\) in \eqref{hmtfg:escalado:eq:30}. Conversely, given \(z\in Z_n\), coinitiality provides \(p\) with \(L_n(p)\le z\). Monotonicity gives \(L_n(p_n^*)\le L_n(p)\le z\). Since \(L_n(p_n^*)\in Z_n\), it is the global minimum. Surjectivity of \(L_n\) is a sufficient condition, stronger than \eqref{hmtfg:escalado:eq:30}.

The ordered isomorphism of the continuum chart satisfies reader surjectivity. The concrete application to the return and its complete set of zeros is proved in the ordered-transport section; classification and isolation then determine its dynamical continuation.

\subsubsection{Minimal selection through survival at arbitrary depth}
\label{hmtfg:escalado:13.6}

Composition with survival also produces a precise global selector. The uniform family \(\eta=0\), constructed in section~\ref{hmtfg:escalado:1}, is retained, and the theorems on realization of the infinite itinerary and preservation of its signs are used. Let
\[
I=\left[\frac1{2u_0(1)},\frac2{u_0(1)}\right],\qquad
c_j(\lambda)=f_\lambda^j(0),\qquad
\tau_j=(-1)^{v_2(j)}.
\]
The set
\[
\Lambda_\tau=\{\lambda\in I:\operatorname{sign}c_j(\lambda)=\tau_j
\text{ for every }j\ge1\}
\]
is compact and nonempty. The derivative control in that source provides
\begin{equation}
B_j=\frac{16^j(16^j-1)}{15},\qquad
|c_j(\lambda)|>\frac2{B_j}\quad(\lambda\in\Lambda_\tau).
\label{hmtfg:escalado:eq:32}
\end{equation}
Define
\begin{equation}
\varepsilon_j=\frac1{B_j},\qquad
A_N=\{\lambda\in I:\tau_jc_j(\lambda)\ge\varepsilon_j,\
1\le j\le N\}.
\label{hmtfg:escalado:eq:33}
\end{equation}
Each \(A_N\) is compact by continuity of the iterates and nonempty because it contains \(\Lambda_\tau\). Moreover,
\begin{equation}
A_{N+1}\subseteq A_N,\qquad
\bigcap_{N\ge1}A_N=\Lambda_\tau.
\label{hmtfg:escalado:eq:34}
\end{equation}
The right-to-left inclusion follows from \eqref{hmtfg:escalado:eq:32}. The reverse inclusion follows from the strictly positive margins in \eqref{hmtfg:escalado:eq:33}, which fix each sign.

\textbf{Surviving-minimum theorem.} The minima \(a_N=\min A_N\) satisfy
\begin{equation}
\boxed{a_N\uparrow a_*=\min\Lambda_\tau.}
\label{hmtfg:escalado:eq:35}
\end{equation}
\textbf{Proof.} The inclusion \eqref{hmtfg:escalado:eq:34} makes the sequence of minima increasing and bounded within \(I\); it therefore has a limit \(a\). For each \(r\), all terms \(a_N\), \(N\ge r\), belong to the closed set \(A_r\); thus \(a\in A_r\). By \eqref{hmtfg:escalado:eq:34}, \(a\in\Lambda_\tau\). For any \(\lambda\in\Lambda_\tau\), we have \(a_N\le\lambda\) at every level, and on passing to the limit, \(a\le\lambda\). This proves \eqref{hmtfg:escalado:eq:35}.

In the HMT cylindrical reader, the object \(a_*\) has compatible prefixes at every depth, with orientation, remainder and memory preserved. Its choice uses the order of the parameter reading and complete survival. Compactness proves the existence and convergence of these minima.

The parameters \(a_N\) solve sign constraints with a margin. Their boundaries may satisfy \(\tau_jc_j=\varepsilon_j\), whereas the superstable roots satisfy \(c_{2^n}=0\). Therefore \eqref{hmtfg:escalado:eq:35} establishes a global selector for infinite realization and preserves, as a distinct object, the corpus’s first-root selector.

\subsubsection{Recovery of the parameter from the complete critical history}
\label{hmtfg:escalado:13.7}

The second term of the critical history provides an exact inverse reader:
\begin{equation}
c_1(\lambda)=1,\qquad
c_2(\lambda)=1-\lambda u_0(1),\qquad
\boxed{\lambda=\frac{1-c_2(\lambda)}{u_0(1)}.}
\label{hmtfg:escalado:eq:36}
\end{equation}
The positivity \(u_0(1)>0\), proved for the profile, makes division valid. Consequently, two identical complete critical histories have the same parameter. The sign word \(\tau\) is a reduced readout of that history; it preserves its combinatorics and omits the magnitude of \(c_2\).

Formulas \eqref{hmtfg:escalado:eq:35}–\eqref{hmtfg:escalado:eq:36} unite survival-based selection and recovery of the parameter from the critical history. Equality of its limit with the stable crossing is proved by first exit, and identification of the centers by analytic transport of their hybrid classes.

% Fragmento integrable. Preambulo anfitrion: amsmath, amssymb, longtable, hyperref.
\section{Ordered transport and Jacobi representation}
\label{hmtfg:fragmento:transporte-ordenado-y-representacion-de-jacobi}

\subsection{Complete transport from the continuum and spectral realization of the profile}
\label{hmtfg:escalado:14}

The ordered transport of the continuum and the spectral representation of the reader are composed before parameter selection.

\subsubsection{Coverage of all roots and transport of their minimum}
\label{hmtfg:escalado:14.1}

The ordered construction of the continuum developed in the preceding works produces the surjective Archimedean readout from compatible cylinders and their ordered quotient. Within its declared universe \(\mathfrak G\), it proves the isomorphism of complete ordered fields
\begin{equation}
\iota:\mathbb R_{\rm HMT}^{\mathfrak G}
   \xrightarrow{\cong}\mathbb R^{\mathfrak G}.
\label{hmtfg:escalado:eq:37}
\end{equation}
The name \(\eta\) used in the source work is changed here to \(\iota\), to distinguish it from the weighting parameter, fixed at zero. The second projection preserves incidence, boundary and memory over the same antecedent; the evaluation quotient and the complete history retain their respective domains.

Let us construct in the HMT field, before selecting a root,
\[
\cos_{\rm H}(t)=\sum_{k=0}^{\infty}\frac{(-1)^kt^{2k}}{(2k)!},
\qquad
u_{\rm H}(x)=\frac1{12}\sum_{m=1}^{12}
\left[1-\cos_{\rm H}\!\left(\frac{2(3m-1)x}{\sqrt[\rm H]{899}}\right)\right],
\]
\begin{equation}
F_a(x)=1-a\,u_{\rm H}(x),\qquad S_n^{\rm H}(a)=F_a^{\,2^n}(0).
\label{hmtfg:escalado:eq:38}
\end{equation}
The integers and the second moment are those already derived from the dodecaphase reader. The square root is the positive one in the ordered field. The convergent series analytically realizes that reader; it preserves the profile and the normalizing coefficient of section~\ref{hmtfg:escalado:1}.

\textbf{Return transport theorem.} For every parameter in the internal domain of \eqref{hmtfg:escalado:eq:38},
\begin{equation}
\iota(F_a(x))=f_{\iota(a)}(\iota(x)),\qquad
\iota(S_n^{\rm H}(a))=S_n(\iota(a)).
\label{hmtfg:escalado:eq:39}
\end{equation}
\textbf{Proof.} The ordered isomorphism fixes the rationals and preserves the positive square root. It preserves convergent limits: every neighborhood of positive rational radius is transported to another of the same radius. Applied to the partial sums of \eqref{hmtfg:escalado:eq:38}, it transports the cosine and \(u_{\rm H}\). The first equality follows from preservation of addition and multiplication; the second follows by induction on the number of iterations.

With a previous center \(a_{n-1}\) fixed, define \(Z_n^{\rm H}\) by zero return, exact critical period \(2^n\) and \(a>a_{n-1}\), within the declared parameter domain. Let \(Z_n\) be the set defined by the same conditions for \(f_\lambda\), to the right of \(\iota(a_{n-1})\). Then
\begin{equation}
\boxed{\iota[Z_n^{\rm H}]=Z_n,\qquad
\iota(\min Z_n^{\rm H})=\min Z_n.}
\label{hmtfg:escalado:eq:40}
\end{equation}
The second equality is asserted when the minimum exists; the first also preserves the nonexistence of candidates.

\textbf{Proof.} \eqref{hmtfg:escalado:eq:39} preserves the zero return and the earlier returns determining the exact period. Order preserves position relative to the previous center. Surjectivity of \eqref{hmtfg:escalado:eq:37}, applied to any root parameter, proves reverse coverage. If \(a_*\) is the minimum, every \(z\in Z_n\) has the form \(\iota(a)\), with \(a\in Z_n^{\rm H}\), whence \(\iota(a_*)\le z\). The converse assertion uses \(\iota^{-1}\).

Surjectivity covers every parameter in the chart, including those whose roots have not been isolated by a finite computation. The isomorphism preserves the declared universe of interpretation and the selection order.

\subsubsection{Jacobi representation of the complete dodecaphase profile}
\label{hmtfg:escalado:14.2}

Define the atomic measure of the reader,
\begin{equation}
s_m=\frac{4(3m-1)^2}{899},\qquad
\nu_0=\frac1{12}\sum_{m=1}^{12}\delta_{s_m},\qquad
M_r=\int s^r\,d\nu_0(s).
\label{hmtfg:escalado:eq:41}
\end{equation}
The twelve nodes are distinct and positive. For a polynomial \(p\ne0\) of degree less than \(r\le12\),
\begin{equation}
\sum_{i,j=0}^{r-1}p_iM_{i+j}p_j
=\frac1{12}\sum_{m=1}^{12}p(s_m)^2>0.
\label{hmtfg:escalado:eq:42}
\end{equation}
A polynomial of degree less than twelve does not vanish at all twelve nodes. At degree twelve, \(\prod_m(s-s_m)\) appears, with zero norm.

Orthogonalization with respect to this measure produces orthonormal polynomials \(\psi_0,\ldots,\psi_{11}\), with \(\psi_0=1\). Let
\[
Q_{mj}=\psi_j(s_m)/\sqrt{12},\quad
D=\operatorname{diag}(s_m),\quad J_{12}=Q^{\mathsf T}DQ.
\]
The following hold
\begin{equation}
Q^{\mathsf T}Q=I,\quad DQ=QJ_{12},\qquad
\boxed{u_0(x)=e_0^{\mathsf T}
\bigl[I-\cos(x\sqrt{J_{12}})\bigr]e_0.}
\label{hmtfg:escalado:eq:43}
\end{equation}
\textbf{Proof.} Orthogonality follows from \eqref{hmtfg:escalado:eq:42}. Multiplication by \(s\) produces a three-term recurrence: entries separated by more than one position vanish by orthogonality and degree. Choosing positive leading coefficients makes the subdiagonals positive. By functional calculus,
\(\cos(x\sqrt{J_{12}})=Q^{\mathsf T}\cos(x\sqrt D)Q\).
The vector \(Qe_0\) is constant, with value \(1/\sqrt{12}\); substituting it recovers the sum in section~\ref{hmtfg:escalado:1}.

The article on Grassmannians develops this mechanism for marked measures and refinements. Here it is applied to \eqref{hmtfg:escalado:eq:41}; rank twelve belongs to this reader. The genealogical marks are retained alongside the reader, in accordance with the marked faithfulness of the source work. \(J_{12}\) represents the analytic profile; the complete state additionally preserves routes, orientation and memory.

The verifier for this extension checks, with exact fractions, the twelve degrees of positive norm, vanishing at degree twelve, the intertwiner \(VT=DV\) in the monic basis and weighted self-adjointness \(HT=T^{\mathsf T}H\). It checks moments of degree zero through thirty; the identity for every degree follows from the intertwiner. Its first coefficients are
\begin{equation}
a_0=2,\qquad \beta_1=\frac{2493348}{808201}.
\label{hmtfg:escalado:eq:44}
\end{equation}

\subsubsection{Oriented sensitivity of the return}
\label{hmtfg:escalado:14.3}

For an exact return of period \(p=2^n\), let
\[
c_j=f_\lambda^j(0),\quad q_j=\partial_\lambda c_j,\quad
D_j=\prod_{k=1}^j f_\lambda'(c_k),\quad D_0=1.
\]
Before the critical return, the factors of \(D_{p-1}\) are nonzero. Differentiation and telescoping give
\begin{equation}
q_{j+1}=-u_0(c_j)+f_\lambda'(c_j)q_j,\qquad
\boxed{\frac{q_p}{D_{p-1}}
=-\sum_{j=1}^{p-1}\frac{u_0(c_j)}{D_j}.}
\label{hmtfg:escalado:eq:45}
\end{equation}
In the canonical word, \(\operatorname{sign}c_j=(-1)^{v_2(j)}\). The sign of \(f_\lambda'(c_j)\) is opposite; therefore,
\begin{equation}
\operatorname{sign}D_j=(-1)^{j+\sum_{k=1}^jv_2(k)}
=(-1)^{s_2(j)},
\label{hmtfg:escalado:eq:46}
\end{equation}
using \(\sum_{k=1}^jv_2(k)=j-s_2(j)\), where \(s_2(j)\) counts the binary ones. The normalized sensitivity is
\begin{equation}
\mathcal T_n=-\sum_{j=1}^{2^n-1}(-1)^{s_2(j)}t_j,\qquad
t_j=\frac{u_0(c_j)}{|D_j|}>0.
\label{hmtfg:escalado:eq:47}
\end{equation}
Nodal positivity of \eqref{hmtfg:escalado:eq:42} determines \(u_0\); \eqref{hmtfg:escalado:eq:47} additionally retains the orbital products and their orientations.

The additional moment representation would require a positive measure \(\rho\) on \([0,1]\) with \(t_j=\int z^{j-1}\,d\rho\). Under that condition,
\begin{equation}
\mathcal T_n=\int_0^1
\frac{1-\prod_{r=0}^{n-1}(1-z^{2^r})}{z}\,d\rho(z)>0
\label{hmtfg:escalado:eq:48}
\end{equation}
for \(\rho\ne0\), extending the integrand by its value \(1\) at zero. The identity is obtained by expanding the binary product.

The rational check at the previously certified period-four root obtains
\[
t_1\approx0.40425224267600139730,\quad
t_2\approx0.04925249167432646403,\quad
t_3\approx0.15740870098294833737
\]
and the rigorous outer interval
\begin{equation}
\begin{gathered}
0.296096033367379523967764444238832345360107<\mathcal T_2,\\
\mathcal T_2<0.296096033367379523967764444238832345360112.
\end{gathered}
\label{hmtfg:escalado:eq:49}
\end{equation}
It also certifies \(t_3>t_2\). A positive measure on \([0,1]\) would impose \(t_3\le t_2\); that specific lifting of orbital weights to positive moments is ruled out. The positive sign of \eqref{hmtfg:escalado:eq:49} remains certified. This check distinguishes the failure of a proposed proof from the failure of the property it sought to prove.

% Fragmento integrable. Preambulo anfitrion: amsmath, amssymb, longtable, hyperref.
\section{Classification and continuation of the first-root selector}
\label{hmtfg:fragmento:clasificacion-y-continuacion-del-selector-de-primeras-raices}

\subsection{Object and provenance}
\label{hmtfg:clasificacion:1}

This appendix preserves the family rendered by the HMT dodecaphase reader:

\[
u_0(x)=\frac1{12}\sum_{m=1}^{12}\left(1-\cos\frac{2(3m-1)x}{\sqrt{899}}\right),
\qquad f_\lambda(x)=1-\lambda u_0(x).
\]

The family originates in the operator and analytic antecedents presented in this development. The existence and compactness of its realizations of the infinite itinerary were proved in the subsections on realization of the word and preservation of its signs.

The following argument is original and uses only those properties and the return operations of the fixed family. It does not use a parameter from another family or the value of a universal constant as an input. The sheets of the itineraries used here are the two monotone branches of the critical reader; the remaining coordinates and HMT memory retain their residence in the state from which the reader originates.

We write
\[
c_j(\lambda)=f_\lambda^j(0),\qquad
\tau_j=(-1)^{v_2(j)},\qquad
W_n=\tau_1\cdots\tau_{2^n-1}.
\]
The itinerary of a critical orbit that returns to the origin for the first time ends in the zero symbol. The itinerary order is the unimodal order with a maximum: before comparing the symbol of index \(j\), the orientation factor is
\[
O_{j-1}(K)=\prod_{i=1}^{j-1}(-K_i).
\]
The ordinary order of the symbols is \(-1<0<1\).


\subsection{Classification before the infinite word}
\label{hmtfg:clasificacion:2}

\textbf{Classification theorem.} Let \(F:I\to I\), with \(I=[\ell,1]\), \(\ell<0\), be a continuous map, strictly increasing to the left of zero and strictly decreasing to its right, whose maximum satisfies \(F(0)=1\). Suppose zero has exact period \(2^n\), with \(n\ge1\), and its critical itinerary \(K\) satisfies \(K<\tau\). Then
\[
\boxed{K=W_n0.}
\]
The conclusion concerns the itinerary and allows distinct parameters in a family to have that same itinerary.

\textbf{Proof.} Write \(d_j=F^j(0)\). For period two the itinerary is \(+0=W_10\). Consider an exact period \(p=2^n\ge4\).

We have \(d_2<0\). Indeed, \(d_2=0\) would give period two. If \(d_2>0\), monotonicity of the positive branch gives
\[
F([d_2,1])=[d_2,d_3]\subset[d_2,1],
\]
and the critical orbit remains in a strictly positive interval, incompatible with a return to zero.

The initial signs are then \(+,-\). The orientation factor for comparing the third sign is \(-1\); hence \(K<\tau\) forces \(d_3>0\). A zero at that position would produce period three.

Nor can \(d_4<0\) occur. Under this hypothesis,
\[
d_2<d_4<0,
\qquad F([d_2,d_4])=[d_3,d_5]\subset(0,1],
\]
since \(d_4=F(d_3)>F(1)=d_2\) and \(d_5=F(d_4)>F(d_2)=d_3\). The positive branch is decreasing, so that
\[
F^2([d_2,d_4])=[d_6,d_4]\subset[d_2,d_4].
\]
Here \(d_6\ge d_2\), because \(d_5\le1\), and \(d_6<d_4\), because \(d_5>d_3\). This strictly negative interval traps the even returns of the critical orbit and excludes a return to zero. Therefore \(d_4\ge0\). If \(p=4\), one obtains precisely \(+-+0=W_20\).

Now suppose \(p\ge8\). Then \(d_4>0\). The orientation factors for the fifth and sixth signs are, respectively, \(-1\) and \(+1\). Since \(\tau_5=+1\) and \(\tau_6=-1\), the inequality \(K<\tau\) forces
\begin{equation}
\operatorname{sign}(d_1,\ldots,d_6)=(+,-,+,+,+,-).
\label{hmtfg:clasificacion:eq:1}
\end{equation}
Zero signs at those positions are excluded by the exact period.

We construct the return interval
\[
J=[d_2,d_4].
\]
The inequality \(F(d_3)=d_4>0>d_6=F(d_5)\), together with \(d_3,d_5>0\), implies \(d_3<d_5\); therefore
\[
F(J)=[d_3,1].
\]
Moreover, \(d_4<d_3\). To check this, the image under \(F\) of the interval with endpoints \(d_3,d_4\) is strictly positive, since its endpoint images are \(d_4,d_5>0\). On that interval \(F^2\) is strictly increasing. Its endpoint values satisfy
\[
F^2(d_3)=d_5>0>d_6=F^2(d_4),
\]
which forces \(d_3>d_4\). Consequently,
\begin{equation}
J\cap F(J)=\varnothing,
\qquad F^2(J)=F([d_3,1])=[d_2,d_4]=J.
\label{hmtfg:clasificacion:eq:2}
\end{equation}

The map \(F^2|_J\) has a unique minimum at zero. Conjugacy by \(h(x)=d_2x\), which reverses orientation, defines
\begin{equation}
G=h^{-1}\circ F^2\circ h:
\left[\frac{d_4}{d_2},1\right]\longrightarrow
\left[\frac{d_4}{d_2},1\right].
\label{hmtfg:clasificacion:eq:3}
\end{equation}
This map satisfies the same maximum hypotheses as \(F\), with \(G(0)=1\), and its critical point has exact period \(p/2\). Its itinerary is
\begin{equation}
K'_j=-K_{2j}.
\label{hmtfg:clasificacion:eq:4}
\end{equation}
All odd-indexed symbols of \(K\) are positive, because the even positions belong to \(J\) and the odd positions to \(F(J)\subset(0,1]\).

The transformation \eqref{hmtfg:clasificacion:eq:4} preserves the comparison order with \(\tau\). Indeed, \(\tau_{2j-1}=+1\) and \(\tau_{2j}=-\tau_j\). If \(q\) is the first position at which \(K'\) and \(\tau\) differ, the first position of difference between \(K\) and \(\tau\) is \(2q\), and
\begin{equation}
O_{2q-1}(K)=-O_{q-1}(K'),
\qquad K_{2q}-\tau_{2q}=-(K'_q-\tau_q).
\label{hmtfg:clasificacion:eq:5}
\end{equation}
The product determining the order is identical. The identities remain valid when the first distinct symbol is the final zero. Thus \(K'<\tau\).

Induction applied to \(G\) gives \(K'=W_{n-1}0\). The positive odd positions and \eqref{hmtfg:clasificacion:eq:4} reconstruct exactly \(K=W_n0\). The theorem is proved. \(\square\)

The theorem also proves that each classified map has the restrictive doubling returns preceding its superstable return, with the domains of \eqref{hmtfg:clasificacion:eq:2}–\eqref{hmtfg:clasificacion:eq:3}. It does not presuppose monotonicity of a family parameter.

\textbf{Symmetric normalization of the fixed family.} When \(F\) is even and defined on \([-1,1]\), the preceding inequalities additionally give
\[
F(d_4)=d_5>d_3=F(d_2)=F(-d_2).
\]
Since both arguments \(d_4,-d_2\) are positive and \(F\) decreases on that branch,
\begin{equation}
0<d_4<-d_2<1.
\label{hmtfg:clasificacion:eq:5a}
\end{equation}
Therefore \(F^2([d_2,-d_2])=[d_2,d_4]\), and the same conjugacy \(h(x)=d_2x\) defines on all of \([-1,1]\) an even map with a maximum whose image is \([d_4/d_2,1]\subset[-1,1]\). It agrees exactly with the operator \(RF=-F(F(-ax))/a\), \(a=-F(1)\), used in the corpus. This symmetric extension preserves the critical return, its period and its itinerary.


\subsection{First parameter with infinite itinerary}
\label{hmtfg:clasificacion:3}

We work on the compact interval already used in the existence proof,
\[
L=\left[\frac1{2u_0(1)},\frac2{u_0(1)}\right].
\]
Let
\[
\Lambda_\tau=\{\lambda\in L:
\operatorname{sign}c_j(\lambda)=\tau_j\text{ for every }j\ge1\}.
\]
By the existence and sign-separation proofs of the prolongation, this set is nonempty and compact. Therefore there exists
\begin{equation}
a=\min\Lambda_\tau.
\label{hmtfg:clasificacion:eq:6}
\end{equation}

\textbf{Comparison lemma.} For every \(\lambda\in[1/(2u_0(1)),a)\), we have \(K_\lambda<\tau\).

\textbf{Proof.} At the left endpoint, the itinerary is \(+^\infty<\tau\). The separation argument for realization of the infinite word proves that the strict comparison sets with \(\tau\) are open: in an infinite itinerary distinct from \(\tau\), the first finite difference persists by continuity; at a superstable parameter, the two neighboring periodic words and the word with a final zero lie on the same side of \(\tau\), by its strict maximality. By definition, the interval preceding \(a\) contains no parameters with itinerary \(\tau\). Its connectedness maintains the initial comparison throughout that interval. \(\square\)


\subsection{Existence of each first new root}
\label{hmtfg:clasificacion:4}

We retain the original selector. The first root is
\[
\lambda_1=\frac1{u_0(1)},
\]
and, for \(n\ge2\), \(\lambda_n\) is the first root of \(c_{2^n}\) to the right of \(\lambda_{n-1}\) satisfying
\begin{equation}
c_{2^j}(\lambda_n)\ne0\quad(1\le j<n).
\label{hmtfg:clasificacion:eq:7}
\end{equation}
Since every first period dividing \(2^n\) is a power of two, \eqref{hmtfg:clasificacion:eq:7} is equivalent to requiring exact critical period \(2^n\).

\textbf{Existence lemma.} All roots of this selector exist and satisfy
\begin{equation}
\lambda_1<\lambda_2<\cdots<\lambda_n<a.
\label{hmtfg:clasificacion:eq:8}
\end{equation}

\textbf{Proof.} The sign \(c_2(a)<0\) implies \(\lambda_1<a\). Suppose \(\lambda_{n-1}<a\) has been constructed, and put \(p=2^{n-1}\). The analytic function \(c_p(\lambda)\) is not identically zero, since \(c_p(a)\ne0\). Its zero set in \([\lambda_{n-1},a]\) is finite, nonempty and does not contain \(a\). Let \(z<a\) be its last zero.

On \((z,a]\), the sign of \(c_p\) is constant and equal to \(s=\tau_p\). For \(g_\lambda=f_\lambda^p\), we have
\[
g_\lambda(0)=c_p(\lambda),\qquad g'_\lambda(0)=0.
\]
A uniform local bound on the second derivative gives
\begin{equation}
c_{2p}(\lambda)
=g_\lambda(c_p(\lambda))
=c_p(\lambda)+O(c_p(\lambda)^2).
\label{hmtfg:clasificacion:eq:9}
\end{equation}
As \(\lambda\) approaches \(z\) from the right, the first-order term dominates; hence \(c_{2p}(\lambda)\) has sign \(s\). At \(a\), its sign is
\[
\tau_{2p}=-\tau_p=-s.
\]
The intermediate value theorem provides a zero \(\beta\in(z,a)\) of \(c_{2p}\). Since \(c_p\ne0\) on that interval, the critical period at \(\beta\) is exactly \(2p\).

Analyticity and \(c_{2p}(a)\ne0\) imply that the zeros of \(c_{2p}\) in \([\lambda_{n-1},a]\) are finite. Among its zeros of exact period \(2p\) to the right of \(\lambda_{n-1}\), there is therefore a first one. It is the original selector \(\lambda_n\), and it satisfies \(\lambda_n\le\beta<a\). \(\square\)

Use of the last auxiliary zero \(z\) belongs solely to the existence proof. The selected parameter remains the \textbf{first new root} \(\lambda_n\).


\subsection{Convergence of the global selector to the first infinite itinerary}
\label{hmtfg:clasificacion:5}

\textbf{Global continuation theorem.} For the fixed HMT family and the original selector, all selected itineraries are canonical and
\begin{equation}
\boxed{K_{\lambda_n}=W_n0,\qquad
\lambda_n\nearrow\min\Lambda_\tau.}
\label{hmtfg:clasificacion:eq:10}
\end{equation}

\textbf{Proof.} The comparison and existence lemmas place each selected parameter below \(a\), with an itinerary less than \(\tau\). The classification theorem identifies that itinerary with \(W_n0\). The sequence is increasing and bounded by \(a\); let \(b\le a\) be its limit.

Fix \(p\ge1\). For \(n\) sufficiently large, \(2p<2^n\). Then the signs of \(c_p(\lambda_n)\) and \(c_{2p}(\lambda_n)\) agree with \(\tau_p\) and \(\tau_{2p}=-\tau_p\), respectively. The uniform bound on second derivatives,
\[
\|(f_\lambda^p)''\|\le B_p,
\qquad B_p=16^p\frac{16^p-1}{15},
\]
and the identity \((f_\lambda^p)'(0)=0\) imply
\[
|c_{2p}(\lambda_n)-c_p(\lambda_n)|
\le\tfrac12B_p|c_p(\lambda_n)|^2.
\]
Since the two terms on the left-hand side have opposite signs,
\begin{equation}
|c_p(\lambda_n)|>2/B_p.
\label{hmtfg:clasificacion:eq:11}
\end{equation}
By continuity, \(c_p(b)\) retains the sign \(\tau_p\) and satisfies \(|c_p(b)|\ge2/B_p>0\). This holds for each \(p\), so that \(b\in\Lambda_\tau\). Minimality of \(a\) gives \(b\ge a\), and therefore \(b=a\). \(\square\)

The complete proof preserves the global selector and produces its itineraries and its limit. The information from the first eight levels is compatible with the result, but the proof of \eqref{hmtfg:clasificacion:eq:10} does not extrapolate those eight levels: it uses induction, invariant intervals and uniform separation of every sign.


\subsection{Logical checks and falsifiers}
\label{hmtfg:clasificacion:7}

\par\noindent\textbf{1.}\enspace The classification theorem requires a strict unimodal maximum and an invariant interval. If either branch loses monotonicity, the intervals \eqref{hmtfg:clasificacion:eq:2} may cease to be restrictive and the classification must be checked again.
\par\noindent\textbf{2.}\enspace The condition \(K<\tau\) is essential. Beyond the limit word, centers of period \(2^n\) with other combinatorics may exist; the theorem does not identify them with \(W_n0\).
\par\noindent\textbf{3.}\enspace Analyticity in the parameter provides isolation and finiteness of the zeros on the compact set. For a merely continuous family, the existence of a strictly subsequent first root requires an additional argument.
\par\noindent\textbf{4.}\enspace A limit of critical roots may, in general, lose signs upon reaching a return of lower period. Inequality \eqref{hmtfg:clasificacion:eq:11} excludes precisely that phenomenon for this sequence and each fixed horizon.
\par\noindent\textbf{5.}\enspace Convergence of the first roots to the first parameter with itinerary \(\tau\) is not equivalent to uniqueness of all parameters with that itinerary and does not by itself determine a metric rate. Transverse crossing and branch identification provide those subsequent conclusions.

% Fragmento integrable. Preambulo anfitrion: amsmath, amssymb, longtable, hyperref.
\section{Transverse entry and local scaling}
\label{hmtfg:fragmento:entrada-transversal-y-escalado-local}

\subsection{Provenance and result}
\label{hmtfg:escalado:1}

APP evaluates addition and multiplication on the two sheets of the lattice, preserving quotient and remainder. TRIT determines regime and orientation. TPK transport selects, transports and updates the state with its carry, boundary and memory records. Its readers act on the same discrete structure of the continuum. The prolongation \texttt{w6→w12→w18→w24→w30→R36→G9} preserves this provenance and distinguishes phase return from memory increment.

The critical reader of the dodecaphase wheel, constructed in the operator antecedents and transported by compatible refinement, produces
\[
\phi_j=\frac{(3j-1)\pi_{\mathrm{HMT}}}{18},\qquad
w_j=\frac1{12},\qquad
k_j=\sqrt{\frac2{\sum_{\ell=1}^{12}w_\ell\phi_\ell^2}}\,\phi_j
=\frac{2(3j-1)}{\sqrt{899}}.
\]
The coordinate \(\pi_{\mathrm{HMT}}\) retains its APP–TRIT–TPK origin; its common factor cancels in this second moment. The resulting reader is
\begin{equation}
u_0(x)=\frac1{12}\sum_{j=1}^{12}(1-\cos k_jx),\qquad
f_\lambda(x)=1-\lambda u_0(x).
\label{hmtfg:escalado:eq:1}
\end{equation}
The focal sector is the uniform one, \(\eta=0\). The preceding extension separately retains the weighted family \(|\eta|\le1/40\). The new numerical bounds correspond to \eqref{hmtfg:escalado:eq:1}.

The family is even and entire, with \(f_\lambda(0)=1\), a quadratic critical point and second moment \(\sum_jw_jk_j^2=2\). The parameter \(\lambda\) ranges over a family downstream of the HMT generator. The target Feigenbaum constant remains outside the inputs of \eqref{hmtfg:escalado:eq:1}.

\textbf{Main result.} There exists a unique
\begin{equation}
\lambda_*\in[1.874038,1.874039]
\label{hmtfg:escalado:eq:2}
\end{equation}
for which the fourth normalized return of \eqref{hmtfg:escalado:eq:1} belongs to the local stable leaf of the real doubling fixed point. The crossing of the family with that leaf is transverse, with quantified separation. For every \(n\ge0\),
\begin{equation}
\boxed{\|R^{4+n}f_{\lambda_*}-g\|_{\mathcal A}
<0.001666\,(0.83996)^n.}
\label{hmtfg:escalado:eq:3}
\end{equation}
The local superstable cascade associated with this crossing has parameters \(\mu_j\) and constants \(C\ne0\), \(0<\theta<1\), such that
\begin{equation}
\mu_j-\lambda_*=C\delta_{\mathrm F}^{-j}\bigl[1+O(\theta^j)\bigr],
\qquad
\frac{\mu_{j-1}-\mu_{j-2}}{\mu_j-\mu_{j-1}}\longrightarrow\delta_{\mathrm F}.
\label{hmtfg:escalado:eq:4}
\end{equation}
Here \(\delta_{\mathrm F}>1\) is the unstable eigenvalue of the doubling operator. Selection of \(\mu_j\) is performed through the local returns of section~\ref{hmtfg:escalado:8}. Identification with the global selector is proved in sections~\ref{hmtfg:escalado:15}--\ref{hmtfg:escalado:17}: all global terms retain the combinatorics and the centers agree at every sufficiently high level when their periods are matched.


\subsection{Analytic chart and effective return}
\label{hmtfg:escalado:2}

The chart of even functions is used
\[
s=x^2,\qquad z=\frac{s-1}{5/2},\qquad
f(x)=1-x^2V(z),\qquad V(z)=\sum_{j\ge0}v_jz^j.
\]
The coordinates of the Banach space are
\begin{equation}
u=10v_0,\qquad \nu=(v_1,v_2,\ldots),\qquad
\|(\Delta u,\Delta\nu)\|_{\mathcal A}
=|\Delta u|+\sum_{j\ge1}|\Delta v_j|.
\label{hmtfg:escalado:eq:5}
\end{equation}
In particular, the constant coefficient of \(V\) has weight ten. This normalization agrees with the Hertling–Spandl chart \(u/10\), section 2 of the source.

Let \(a=-f(1)=v_0-1\). For admissible real returns,
\begin{equation}
(Rf)(x)=-\frac{f(f(-ax))}{a}.
\label{hmtfg:escalado:eq:6}
\end{equation}
The program evaluates an exact factorization of \eqref{hmtfg:escalado:eq:6}. Defining
\[
S=\frac{a^2s-1}{5/2},\qquad b=V(S),\qquad
h=1-a^2sb,\qquad t=\frac{h^2-1}{5/2},
\]
\[
\operatorname{shift}V(z)=\sum_{j\ge0}v_{j+1}z^j,
\]
one obtains
\begin{equation}
\boxed{V_R=a\,b\,(h+1)\left[V(t)+\frac25\operatorname{shift}V(t)\right].}
\label{hmtfg:escalado:eq:7}
\end{equation}
To verify it, one uses \(v_0=1+a\),
\(h^2-1=-a^2sb(h+1)\) and
\(V(t)-v_0=t\operatorname{shift}V(t)\). Thus,
\(a+1-h^2V(t)=a^2sb(h+1)[V(t)+(2/5)\operatorname{shift}V(t)]\),
which is precisely \(asV_R\).

The initial moments are rational:
\begin{equation}
\frac1{12}\sum_jk_j^{2m}
=\frac1{12}\left(\frac4{899}\right)^m
\sum_{j=1}^{12}(3j-1)^{2m}.
\label{hmtfg:escalado:eq:8}
\end{equation}
Therefore both \eqref{hmtfg:escalado:eq:7} and the initial construction use rational intervals; the initial trigonometric functions are represented by their series with an outer remainder bound.


\subsection{External bounds used as subsequent recognition}
\label{hmtfg:escalado:3}

Once \eqref{hmtfg:escalado:eq:1} has been constructed, Lanford’s hyperbolicity theorem is used as an analytic certificate for the operator \eqref{hmtfg:escalado:eq:6}. Let \(g_0\) be the published central polynomial and \(g\) the fixed point. The estimates used are
\[
\|g-g_0\|_{\mathcal A}<\varepsilon_0=4\,10^{-5},
\]
and, in the ball \(\|f-g_0\|_{\mathcal A}<0.01\), the blocks of \(DR\), written \(R=(U,V)\), satisfy
\begin{equation}
U_u\ge a_0=4.521,\quad
\|U_\nu\|\le b_0=0.560,\quad
\|V_u\|\le c_0=0.756,\quad
\|V_\nu\|\le d_0=0.719.
\label{hmtfg:escalado:eq:9}
\end{equation}
The fixed point has a unique unstable direction, with positive real eigenvalue \(\delta_{\mathrm F}>1\); the rest of the spectrum lies inside the unit disk. The figures in \eqref{hmtfg:escalado:eq:9} come from the published estimates, not from a new reproduction of Lanford’s proof in this dossier. The coefficients of \(g_0\) are explicitly incorporated into the program’s central reference polynomial, in accordance with table 2 of Hertling–Spandl.

Sources: \href{https://repo-archives.ihes.fr/A_GARDER/20180409/Test/I_Prepublications/LANFORD/1968-2013/P_81_17/P_81_17_web.pdf}{Lanford, hyperbolicity estimates}; \href{https://arxiv.org/pdf/1410.3277}{Hertling–Spandl, analytic chart and table 2}.


\subsection{Rational entry and direction certificate}
\label{hmtfg:escalado:4}

The renormalization entry certificate divides \eqref{hmtfg:escalado:eq:2} into sixteen equal rational intervals. It propagates four returns of \eqref{hmtfg:escalado:eq:7} and the exact derivative with respect to \(\lambda\), using differentiation of the composition and outward rounding to one hundred decimal places.

Each series consists of a polynomial of degree 32 with interval coefficients and an arbitrary analytic error \(E\), \(\|E\|_1\le\epsilon\). This error may also modify low-order coefficients. For \(\|q\|_1<1\),
\begin{equation}
\|E\circ q\|_1\le\epsilon,\qquad
\|E'\circ q\|_1\le\frac{\epsilon}{(1-\|q\|_1)^2},\qquad
\|\operatorname{shift}E\|_1\le\epsilon.
\label{hmtfg:escalado:eq:10}
\end{equation}
The products include polynomial–error and error–error terms. The initial tail is bounded by its first absolute term and a geometric ratio, using \(\|1+(5/2)z\|_1=7/2\). Each composed argument remains strictly inside the unit ball of this algebra.

For \(\Gamma(\lambda)=R^4f_\lambda=(u_\lambda,\nu_\lambda)\), the certificate gives the following conservative bounds, uniform in \eqref{hmtfg:escalado:eq:2}:

\begingroup\small
\begin{longtable}{p{0.465\linewidth}p{0.465\linewidth}}
\hline
Quantity & Certified outer bound \\
\hline
\(\|\Gamma(\lambda)-g_0\|_{\mathcal A}\) & \(<0.004993128\) \\
\(\|\nu_\lambda-\nu_{g_0}\|_1\) & \(<0.001396077\) \\
\(u'_\lambda\) & \(>3821.289115\) \\
\(\|\nu'_\lambda\|_1\) & \(<551.757000\) \\
\(\|\nu'_\lambda\|_1/u'_\lambda\) & \(<0.144386<1/4\) \\
Analytic error of the function & \(<8.482\,10^{-19}\) \\
Analytic error of the tangent & \(<4.569\,10^{-15}\) \\
\hline\end{longtable}
\endgroup

The complete decimals, the sixteen subintervals and every intermediate return are retained in the rational entry certificate.

The cone \(\|\dot\nu\|_1\le|\dot u|/4\) is invariant under \eqref{hmtfg:escalado:eq:9}:
\[
|\dot u_{\rm nuevo}|\ge(4.521-0.560/4)|\dot u|
=4.381|\dot u|,
\]
\begin{equation}
\|\dot\nu_{\rm nueva}\|_1\le(0.756+0.719/4)|\dot u|
=0.93575|\dot u|<\tfrac14(4.381)|\dot u|.
\label{hmtfg:escalado:eq:11}
\end{equation}
The certified parameter direction therefore enters a uniformly expanding cone of the operator.


\subsection{Quantitative construction of the stable leaf}
\label{hmtfg:escalado:5}

Translate \(g\) to the origin and write the coordinates as \((x,y)=(u-u_g,\nu-\nu_g)\). Fix
\begin{equation}
L=0.16,\qquad r=0.004,\qquad
A=a_0-Lc_0=4.40004,\qquad q=d_0+c_0L=0.83996.
\label{hmtfg:escalado:eq:12}
\end{equation}
The cylinder \(|x|\le Lr\), \(\|y\|_1\le r\) lies inside the ball of \eqref{hmtfg:escalado:eq:9}, since
\((1+L)r+\varepsilon_0=0.00468<0.01\).

Consider the complete space of functions \(\sigma:B_r\to\mathbb R\) with \(\sigma(0)=0\) and Lipschitz constant at most \(L\), equipped with the uniform distance. For each \(y\), define \((\mathcal G\sigma)(y)\) as the root in \([-L\|y\|_1,L\|y\|_1]\) of
\begin{equation}
F_\sigma(x,y)=U(x,y)-\sigma(V(x,y)).
\label{hmtfg:escalado:eq:13}
\end{equation}
On that interval, \(\|V(x,y)\|_1\le q\|y\|_1<r\). As \(x\) increases, \eqref{hmtfg:escalado:eq:9} shows that \(F_\sigma\) grows with slope at least \(A\). At its endpoints,
\[
F_\sigma(L\|y\|_1,y)
\ge[AL-b_0-Ld_0]\|y\|_1,
\]
and the opposite inequality holds at the negative endpoint. The margin is
\begin{equation}
AL-b_0-Ld_0=0.0289664>0.
\label{hmtfg:escalado:eq:14}
\end{equation}
A unique root is obtained. Comparing \eqref{hmtfg:escalado:eq:13} for two values of \(y\),
\[
\operatorname{Lip}(\mathcal G\sigma)
\le\frac{b_0+Ld_0}{A}
=\frac{0.67504}{4.40004}<0.16.
\]
Comparing two graphs,
\begin{equation}
\|\mathcal G\sigma-\mathcal G\widetilde\sigma\|_\infty
\le A^{-1}\|\sigma-\widetilde\sigma\|_\infty,
\qquad A^{-1}<0.228.
\label{hmtfg:escalado:eq:15}
\end{equation}
The contraction theorem produces a unique fixed graph \(x=\sigma_*(y)\). On it,
\begin{equation}
\|y_n\|_1\le q^n\|y_0\|_1,
\qquad
\|(x_n,y_n)\|_{\mathcal A}\le(1+L)q^n\|y_0\|_1.
\label{hmtfg:escalado:eq:16}
\end{equation}
Moreover, the vertical distance \(d(x,y)=x-\sigma_*(y)\) satisfies
\[
|d(R(x,y))|\ge A|d(x,y)|
\]
as long as the orbit remains in the cylinder. Every orbit that remains there indefinitely belongs to the graph. This graph identifies exactly the local stable leaf and provides estimate \eqref{hmtfg:escalado:eq:16}. Its local regularity agrees with the analytic stable manifold of the hyperbolic fixed point.


\subsection{Existence, uniqueness and transversality of the parameter}
\label{hmtfg:escalado:6}

For \(\lambda\) in \eqref{hmtfg:escalado:eq:2}, define
\[
H(\lambda)=u_\lambda-u_g-
\sigma_*(\nu_\lambda-\nu_g).
\]
The certificate ensures \(\|\nu_\lambda-\nu_g\|_1<0.001436077<r\). For \(\lambda_2>\lambda_1\),
\begin{equation}
H(\lambda_2)-H(\lambda_1)
\ge\int_{\lambda_1}^{\lambda_2}
\bigl[u'_t-L\|\nu'_t\|_1\bigr]dt
>3733.007995\,(\lambda_2-\lambda_1).
\label{hmtfg:escalado:eq:17}
\end{equation}
To determine the endpoint signs, \(\|g-g_0\|_{\mathcal A}<\varepsilon_0\) and \(|\sigma_*(y)|\le L\|y\|_1\) are used. Rational point evaluation produces
\begin{equation}
H(1.874038)<-0.0011856110945,
\qquad H(1.874039)>0.0021866053399.
\label{hmtfg:escalado:eq:18}
\end{equation}
By continuity, a root exists; \eqref{hmtfg:escalado:eq:17} proves its uniqueness and the transversality of the crossing. At that root, \eqref{hmtfg:escalado:eq:16} and
\((1+L)(0.001396077+0.00004)<0.001666\)
prove \eqref{hmtfg:escalado:eq:3}.

The bound is uniform in the number of returns: it converts increasing depth into an explicit geometric operator error.


\subsection{Real admissibility of all returns}
\label{hmtfg:escalado:7}

The four initial returns are accompanied, on each rational subinterval, by the inequalities
\begin{equation}
0<a=-f(1)<1,\qquad b=f(a)>a,\qquad f(b)<a.
\label{hmtfg:escalado:eq:19}
\end{equation}
The program stores \(a,b,f(b)\) before each application of \eqref{hmtfg:escalado:eq:7}. Unimodality, evenness and the quadratic critical point are preserved under these normalized returns.

For subsequent iterations, a uniform check in the ball of \eqref{hmtfg:escalado:eq:9} suffices. If \(E=V-V_{g_0}\), the norm \eqref{hmtfg:escalado:eq:5} gives, for \(-0.4\le z\le0\),
\[
|E(z)|\le0.004,\qquad |E'(z)|\le0.01.
\]
The second inequality uses \(\sup_{j\ge1}j(0.4)^{j-1}=1\). Rational evaluation of the central coefficients provides
\begin{equation}
\frac65<V(z)<\frac85,\qquad |V'(z)|<\frac12,
\qquad 0.39<a<0.41.
\label{hmtfg:escalado:eq:20}
\end{equation}
Consequently,
\[
f'(x)=-2x\left[V(z)+\frac25x^2V'(z)\right],
\]
and the bracketed expression exceeds one for \(0\le x\le1\). The map is unimodal and preserves \([-1,1]\). Moreover,
\[
b=f(a)>1-(0.41)^2\frac85=\frac{4569}{6250}>a,
\]
\begin{equation}
f(b)<1-\left(\frac{4569}{6250}\right)^2\frac65
=\frac{35028967}{97656250}<0.39<a.
\label{hmtfg:escalado:eq:21}
\end{equation}
The orbits of the stable leaf satisfy \eqref{hmtfg:escalado:eq:19} at every depth. Thus the analytic iterations of the operator correspond to real doubling returns, with their domains and orientations preserved.


\subsection{Scaling of the local superstable cascade}
\label{hmtfg:escalado:8}

The metric conclusion requires two crossings: that of the family with the stable leaf, proved in section~\ref{hmtfg:escalado:6}, and that of the unstable manifold with the superstable target. The second is proved by Eckmann–Wittwer for
\(\Sigma_2=\{\psi:\psi(1)=0\}\), corresponding to the period-two superstable cycle. Normalization \eqref{hmtfg:escalado:eq:6} agrees with that of the even real operator used there.

The general theorem of Collet–Eckmann–Lanford, section 6 of the cited source, proposition 6.1 and theorems 6.2–6.3, applies to a Banach-space map \(C^2\), a fixed point with a unique unstable eigenvalue \(\delta_{\mathrm F}>1\), a curve transverse to its stable leaf and a target transverse to its unstable manifold. The local preimages of the target intersect the curve at points satisfying
\begin{equation}
\delta_{\mathrm F}^j(\mu_j-\lambda_*)\longrightarrow C\ne0.
\label{hmtfg:escalado:eq:22}
\end{equation}
The index \(j\) counts local returns. Since the initial curve here is \(R^4f_\lambda\), the physical period is \(2^{j+5}\) when the target is taken in \(\Sigma_2\); any fixed number of steps used to bring the target into the local neighborhood is incorporated into a fixed index shift and into \(C\).

The real conditions of section~\ref{hmtfg:escalado:7} and the target’s real doubling branch fix the exact period and itinerary. The sequence is selected by this oriented local continuation. Taking differences in \eqref{hmtfg:escalado:eq:22} gives the ratio limit of \eqref{hmtfg:escalado:eq:4}.

Sources of the two results used: \href{https://link.springer.com/article/10.1007/BF01013368}{Eckmann–Wittwer, \emph{A complete proof of the Feigenbaum conjectures}}; \href{https://people.math.harvard.edu/~knill/history/lanford/papers/ColletEckmannLanford.pdf}{Collet–Eckmann–Lanford, section 6}.

\subsubsection{Power-law asymptotic remainder}

The analyticity of our family allows \eqref{hmtfg:escalado:eq:22} to be made more precise. In the normal-coordinate construction of Collet–Eckmann–Lanford, only the stable manifold, the unstable manifold and the target hypersurface are flattened. These three submanifolds admit local \(C^2\) changes of coordinates with bounded derivatives. The smooth cutoff used is applied in the scalar unstable coordinate; the complete foliation is then obtained through the constructed change of coordinates.

In those coordinates, the correction is written as a series of increments \(w_j\) whose norm \(C^1\) satisfies
\[
\|w_j\|_{C^1}\le C_1\rho^j,\qquad 0<\rho<1.
\]
The chain rule for the cut-off map, with uniformly bounded first- and second-order derivatives, provides
\[
\|w_j\|_{C^2}\le C_2B^j
\]
for some \(B\ge1\). The weighted norm of the construction controls the ordinary \(C^1\) norm on the domain considered. By interpolation,
\[
[Dw_j]_\beta\le C_3
\left(\rho^{1-\beta}B^\beta\right)^j.
\]
Choose
\begin{equation}
0<\beta<\min\left\{1,
\frac{-\log\rho}{\log B-\log\rho}\right\}.
\label{hmtfg:escalado:eq:23}
\end{equation}
The series converges in \(C^{1,\beta}\). Let \(z\) be the resulting normal coordinate, normalized so that the preimages of the target satisfy \(z=\delta_{\mathrm F}^{-j}\), absorbing the fixed index shift. Since our curve is analytic and transverse,
\[
z(\Gamma(\lambda))=c(\lambda-\lambda_*)
+O(|\lambda-\lambda_*|^{1+\beta}),\qquad c\ne0.
\]
Local inversion gives
\begin{equation}
\mu_j-\lambda_*=C\delta_{\mathrm F}^{-j}
+O\left(\delta_{\mathrm F}^{-(1+\beta)j}\right).
\label{hmtfg:escalado:eq:24}
\end{equation}
This proves the remainder in \eqref{hmtfg:escalado:eq:4}, with \(\theta=\delta_{\mathrm F}^{-\beta}\). The statement asserts the existence of the constants in the parameter remainder. The explicit operator rate \(0.83996\) of \eqref{hmtfg:escalado:eq:3}, for its part, controls convergence of the renormalized maps.

\subsubsection{Conversion of the remainder into ratio error}

If \(\mu_j-\lambda_*=C\delta_{\mathrm F}^{-j}(1+e_j)\), with \(|e_j|\le M\theta^j\), write
\[
\mu_{j-1}-\mu_j=C(\delta_{\mathrm F}-1)\delta_{\mathrm F}^{-j}(1+r_j),
\quad
r_j=\frac{\delta_{\mathrm F} e_{j-1}-e_j}{\delta_{\mathrm F}-1}.
\]
With \(K=M(\delta_{\mathrm F}+\theta)/(\delta_{\mathrm F}-1)\), we have \(|r_j|\le K\theta^{j-1}\). Therefore,
\begin{equation}
\boxed{
\left|\frac{\mu_{j-2}-\mu_{j-1}}{\mu_{j-1}-\mu_j}-\delta_{\mathrm F}\right|
\le\frac{\delta_{\mathrm F} K(1+\theta)\theta^{j-2}}
{1-K\theta^{j-1}}}
\label{hmtfg:escalado:eq:25}
\end{equation}
when \(K\theta^{j-1}<1\). This formula separates the evaluation precision of each root from the truncation error of the cascade.


\subsection{Certified continuation of the eight initial roots}
\label{hmtfg:escalado:9}

The finite continuation certificate preserves the family \eqref{hmtfg:escalado:eq:1} and the previous eight-root certificate. For \(2\le n\le8\), it studies
\[
S_n(\lambda)=f_\lambda^{2^n}(0)
\]
from the certified box of \(\lambda_{n-1}\) to \(1.874039\). The complete rational covering contains exactly the inherited root \(\lambda_{n-1}\) and the new root \(\lambda_n\). In particular, \(\lambda_n\) is the unique new root in \((\lambda_{n-1},1.874039]\).

Boxes are excluded by the sign of the return or by a local derivative of certified sign. Anchor boxes retain existence and uniqueness through endpoint signs and local monotonicity. Their endpoints cover each interval exactly; adjacencies are checked. Evaluation uses the moments \eqref{hmtfg:escalado:eq:8}, order-32 Horner evaluation, outer geometric tails and \(|u_0''|\le2\). The uniform tail of the potential is less than \(1.294\,10^{-58}\) on \(|x|\le3/2\).

The final certificate contains 650 processed boxes and 332 covering leaves. The first level has the exact expression \(\lambda_1=1/u_0(1)\). The following approximations are shown solely for guidance; their complete boxes are retained in the certificates:

\begingroup\small
\begin{longtable}{p{0.465\linewidth}p{0.465\linewidth}}
\hline
Level & Superstable parameter \\
\hline
1 & 1.345913990471832792210949095… \\
2 & 1.763379787846910127290794030… \\
3 & 1.850413796950136464530566778… \\
4 & 1.868979756606798125150574309… \\
5 & 1.872955034435659740004205128… \\
6 & 1.873806367467030119412262311… \\
7 & 1.873988695221365362307168780… \\
8 & 1.874027744154450672897007573… \\
\hline\end{longtable}
\endgroup

% Fragmento integrable. Preambulo anfitrion: amsmath, amssymb, longtable, hyperref.
\section{Isolation of the limit and scaling of the global selector}
\label{hmtfg:fragmento:aislamiento-del-limite-y-escalado-del-selector-global}

\subsection{Indefinite continuation of the first root and exclusion of the intermediate interval}
\label{hmtfg:escalado:15}

The complete classification proved above preserves the family and the first-new-root selector. Ordered transport preserves the zeros and their minimum; the return word preserves the nonadic chronological reading.

\subsubsection{Classification and existence at every depth}
\label{hmtfg:escalado:15.1}

Write \(W_n=\tau_1\cdots\tau_{2^n-1}\), with \(\tau_j=(-1)^{v_2(j)}\), and \(a_*=\min\Lambda_\tau\). The existence and compactness of \(\Lambda_\tau\) are proved in the prolongation development, \ref{hmtfg:prolongacion:6.1}–\ref{hmtfg:prolongacion:6.2}; its minimum exists by compactness and order.

\textbf{Classification.} Every strictly unimodal map with a maximum, on an invariant interval, whose critical point has exact period \(2^n\) and whose itinerary \(K\) satisfies \(K<\tau\), has itinerary \(W_n0\).

The induction proceeds through the effective return. For period at least eight, comparison with \(\tau\), together with the intervals that would trap orbits with incompatible signs, forces the prefix \((+,-,+,+,+,-)\). This yields
\[
J_1=[d_2,d_4],\quad F(J_1)=[d_3,1],\quad
d_4<d_3,\quad F^2(J_1)=J_1.
\]
The return normalized by \(h(x)=d_2x\) has maximum one, half the period, and itinerary \(K'_j=-K_{2j}<\tau\). Order is preserved because
\[
O_{2q-1}(K)=-O_{q-1}(K'),\qquad
K_{2q}-\tau_{2q}=-(K'_q-\tau_q).
\]
The initial cases are \(+0\) and \(+-+0\); induction reconstructs \(W_n0\). For even maps, also \(0<d_4<-d_2<1\), so the return extends to \([-1,1]\) and coincides with \eqref{hmtfg:escalado:eq:6}.

Before \(a_*\), all itineraries are smaller than \(\tau\): the strict-comparison sets are open by the separation proved in \ref{hmtfg:prolongacion:6.1}; the interval preceding \(a_*\) is connected and begins with \(+^\infty<\tau\).

\textbf{Existence of each first root.} Suppose \(\lambda_{n-1}<a_*\) has been constructed, and let \(p=2^{n-1}\). The zeros of \(c_p\) in \([\lambda_{n-1},a_*]\) form a nonempty finite set: the function is analytic and \(c_p(a_*)\ne0\). Let \(z<a_*\) be its last zero. In \((z,a_*]\), \(c_p\) has sign \(\tau_p\). For \(g_\lambda=f_\lambda^p\),
\[
g_\lambda'(0)=0,\qquad
c_{2p}(\lambda)=c_p(\lambda)+O(c_p(\lambda)^2).
\]
Near \(z\), on the right, \(c_{2p}\) has sign \(\tau_p\); at \(a_*\), it has the opposite sign. There is a zero of \(c_{2p}\) in \((z,a_*)\), and its period is exactly \(2p\), since there \(c_p\ne0\). Finiteness of the zeros provides the first new root in \((\lambda_{n-1},a_*)\). The last auxiliary zero proves existence; the effective choice remains the first root. The base case is \(\lambda_1=1/u_0(1)<a_*\).

\subsubsection{Limit of the original sequence}
\label{hmtfg:escalado:15.2}

\textbf{Theorem.} The original selector exists at every level and satisfies
\begin{equation}
\boxed{K_{\lambda_n}=W_n0,\qquad
\lambda_n\nearrow a_*=\min\Lambda_\tau.}
\label{hmtfg:escalado:eq:50}
\end{equation}
\textbf{Proof.} The preceding existence result gives an increasing sequence bounded by \(a_*\). The classification fixes all its itineraries. Let \(b\le a_*\) be its limit. For each \(p\), sufficiently late terms preserve the opposite signs \(\tau_p,\tau_{2p}\). Taylor’s formula and the uniform bound \(B_p=16^p(16^p-1)/15\) give
\[
|c_{2p}(\lambda_n)-c_p(\lambda_n)|
\le\tfrac12 B_p|c_p(\lambda_n)|^2,\qquad
|c_p(\lambda_n)|>2/B_p.
\]
Passing to the limit preserves \(\operatorname{sign}c_p(b)=\tau_p\), with absolute value at least \(2/B_p>0\). For every \(p\), this implies \(b\in\Lambda_\tau\); minimality gives \(b=a_*\).

The result establishes indefinite existence, preservation of the symbolic cascade, and selection of its first limiting parameter for the original global rule. The HMT completion transports this sequence and its limit while preserving order and reading.

\subsubsection{Certificate excluding the intermediate interval}
\label{hmtfg:escalado:15.3}

The bridge program certifies
\begin{equation}
\Lambda_\tau\cap[\lambda_{8,\rm inf},1.874038]=\varnothing,
\label{hmtfg:escalado:eq:51}
\end{equation}
where
\[
\lambda_{8,\rm inf}
=1.874027744154450672897007573595092408435133414523352839687903472138975
\]
is the rational lower endpoint of the certified interval for \(\lambda_8\).

The cover contains 374 exactly adjacent intervals and no unresolved region: 319 with \(c_{512}>0\), 10 with \(c_{1024}<0\), two with \(c_{2048}>0\), all opposite to the sign of \(\tau\), and 43 near centers of periods 256, 512 or 1024, excluded by their second return.

In the latter case, \(g=f_\lambda^p\), \(d=g(0)\) and a uniform bound \(|g''|\le M\) between zero and \(d\) satisfy \(M|d|<2\). Then
\[
|g(d)-d|\le\tfrac12M|d|^2<|d|\quad(d\ne0).
\]
Both returns have the same sign; for \(d=0\), both are zero. Both situations exclude \(\tau_{2p}=-\tau_p\) with strict signs.

The evaluation uses rational moments, geometric tails of the potential and its second derivative, and \(|u_0'''|<6\). Parameter centering intersects the natural enclosure with
\[
c_j(\lambda_{\rm medio})+
\partial_\lambda c_j(I)\,(I-\lambda_{\rm medio});
\]
the derivative is propagated over the full box. The self-map property of \([-1,1]\) is proved before restricting the enclosures to it.

The complete certificate retains the cover and checks its exact tiling, together with every sign or Taylor inequality.

\subsubsection{Joint localization}
\label{hmtfg:escalado:15.4}

The local crossing \(\lambda_*\) realizes \(\tau\) through its real returns at every depth, so that \(a_*\le\lambda_*\). By \eqref{hmtfg:escalado:eq:50}, \(\lambda_8<a_*\); together with \eqref{hmtfg:escalado:eq:51},
\begin{equation}
\boxed{1.874038<a_*\le\lambda_*<1.874039.}
\label{hmtfg:escalado:eq:52}
\end{equation}
The localization preserves the canonical itinerary and places its first limiting parameter in the same interval as the stable crossing.

\subsection{First-exit isolation and coincidence of the limits}
\label{hmtfg:escalado:16}

This section completes the isolation identified in section~\ref{hmtfg:escalado:15.4}. It preserves the HMT family \eqref{hmtfg:escalado:eq:1}, the minimal selection \eqref{hmtfg:escalado:eq:50}, the chart \eqref{hmtfg:escalado:eq:5} and the stable sheet of sections~\ref{hmtfg:escalado:5}--\ref{hmtfg:escalado:7}. The proof uses the order \(a_*\le\lambda_*\): it suffices to exclude a separation toward the negative side of the expanding cone.

\subsubsection{Upper and lower bounds for transverse transport}
\label{hmtfg:escalado:16.1}

Besides \eqref{hmtfg:escalado:eq:9}, Lemma 2.1 of \href{https://arxiv.org/pdf/1410.3277}{Hertling–Spandl, section 2} provides \(\|D\Phi\|<0.9\) in the same ball of radius \(0.01\), where
\[
\Phi_u(u,\nu)=u-\frac{U(u,\nu)-u}{3.669}.
\]
Therefore,
\begin{equation}
\left|1-\frac{U_u-1}{3.669}\right|<0.9,\qquad
U_u<1+1.9(3.669)=7.9711.
\label{hmtfg:escalado:eq:53}
\end{equation}
The number \(3.669\) belongs to the published preconditioner of this analytic certificate; the generated family \eqref{hmtfg:escalado:eq:1} retains its prior coefficients.

Take \(\kappa=0.21\). For two points in the ball joined by a secant with
\(\Delta u<0\), \(\|\Delta\nu\|_1\le\kappa|\Delta u|\), integration of the blocks of \(DR\) along the segment gives
\begin{equation}
4.4034|\Delta u|
\le|\Delta u'|\le8.0887|\Delta u|,\qquad \Delta u'<0,
\label{hmtfg:escalado:eq:54}
\end{equation}
\begin{equation}
\|\Delta\nu'\|_1
\le(0.756+0.719\kappa)|\Delta u|
=0.90699|\Delta u|
<\kappa(4.4034)|\Delta u|.
\label{hmtfg:escalado:eq:55}
\end{equation}
Here \(4.4034=4.521-0.560\kappa\) and \(8.0887=7.9711+0.560\kappa\). Thus, the secant cone preserves orientation and expands while the segment remains in the ball.

\subsubsection{Reduction of an infinite separation to a finite region}
\label{hmtfg:escalado:16.2}

Suppose \(a_*<\lambda_*\), and set
\[
f_n=R^{4+n}f_{a_*},\qquad
s_n=R^{4+n}f_{\lambda_*},\qquad
(\Delta u_n,\Delta\nu_n)=f_n-s_n.
\]
The positive bound for \(u'\) and \(\|\nu'\|_1/u'<0.144386<\kappa\) on \(\Gamma(J)\) place the initial secant in the negative cone. The stable sheet is written relative to \(g\), with
\begin{equation}
s_n=g+e_n,\quad
\|e_n\|_{\mathcal A}<0.001666(0.83996)^n,\quad
|(e_n)_u|\le0.16\|(e_n)_\nu\|_1.
\label{hmtfg:escalado:eq:56}
\end{equation}
The entry checks imply
\begin{equation}
|\Delta u_0|
<0.004993128+0.00004
 +0.16(0.001396077+0.00004)
=0.00526290032< h,\qquad h=0.006.
\label{hmtfg:escalado:eq:57}
\end{equation}
While \(|\Delta u_n|<h\), \eqref{hmtfg:escalado:eq:55}–\eqref{hmtfg:escalado:eq:56} give
\begin{equation}
\|f_n-g_0\|_{\mathcal A}
<(1+\kappa)h+0.001666+0.00004
=0.008966<0.01.
\label{hmtfg:escalado:eq:58}
\end{equation}
The point \(s_n\) and the segment between the two also remain in that convex ball. The lower expansion \eqref{hmtfg:escalado:eq:54} forces a finite first exit \(N\), and the upper bound applied to the preceding step provides
\begin{equation}
0.006\le-\Delta u_N<0.0485322,\qquad
\|\Delta\nu_N\|_1\le0.21|\Delta u_N|.
\label{hmtfg:escalado:eq:59}
\end{equation}
The inequality \(8.0887h=0.0485322\) is used; the full jump, not only the first-exit surface, is included.

Consequently, \(f_N\) belongs to the following region:
\begin{equation}
\begin{split}
\mathcal C_-=\{\,g+(d_u,d_\nu)+e:\;&
-0.0485322\le d_u\le-0.006,\\
&\|d_\nu\|_1\le0.21|d_u|,\quad
|e_u|+\|e_\nu\|_1\le0.001666,\\
&|e_u|\le0.16\|e_\nu\|_1\,\}.
\end{split}
\label{hmtfg:escalado:eq:60}
\end{equation}
Since \(a_*\) realizes \(\tau\), all its normalized returns are unimodal self-maps of \([-1,1]\) with the same itinerary. This property follows from return induction, independently of the distance from \(f_N\) to \(g_0\). This is why the following certificate treats \(\mathcal C_-\) intersected with that class of self-maps.

\subsubsection{Certified exclusion of the first-exit region}
\label{hmtfg:escalado:16.3}

The rational first-exit exclusion certificate, reproduced in the computational supplement, covers the interval of \eqref{hmtfg:escalado:eq:59} exactly with 64 adjacent rational bands. All are excluded; no additional subdivision or unresolved region enters the result.

Of the 64 bands, 54 display an opposite sign and ten a contracting return of period 16 or 32. The largest Lipschitz bound for those returns is less than \(0.475146491\); the maximum exclusion horizon is iterate 64.

The evaluation retains the first forty perturbation coefficients as shared variables, together with an explicit bound on their infinite tail. For each band, the center is \(g_0\) shifted in coordinate \(u\). If \(a,b\ge0\) bound the dual sensitivities in \(u,\nu\), the support of the uncertainty in the fixed point and the stable residual is
\begin{equation}
0.00004\max(a,b)
+0.001666\max\left(b,\frac{0.16a+b}{1.16}\right).
\label{hmtfg:escalado:eq:61}
\end{equation}
The second expression is obtained by maximizing \(a|e_u|+b\|e_\nu\|_1\) under the two constraints of \eqref{hmtfg:escalado:eq:56}. This uses the orientation of the stable sheet proved previously, as well as its norm.

Let \(x_j\) be the center’s orbit and \(L_j\) its linear sensitivity to all shared coefficients. The Taylor model has the form
\[
c_j=x_j+L_j(\Delta v)+r_j,\qquad |r_j|\le E_j.
\]
If \(\rho_j\) bounds \(|L_j(\Delta v)|+E_j\), the outward recurrence used is
\begin{equation}
E_{j+1}\le |f_0'(x_j)|E_j
+\tfrac12\sup|f_0''|\rho_j^2
+\sup|\Delta f'|\rho_j.
\label{hmtfg:escalado:eq:62}
\end{equation}
The derivatives are bounded along the segment between the two states. The central orbit is verified within \([-1,1]\); the other orbits belong to that interval by the self-map property included in the certified domain. The calculations use outward rational rounding, including for the geometric tails.

Each band is excluded by one of two inequalities:

\par\noindent\textbf{1.}\enspace A critical value has a sign strictly opposite to \(\tau_j\).
\par\noindent\textbf{2.}\enspace For \(p=2^k\), the return \(H=f^p\) has Lipschitz constant \(L_p<1\) on the segment between \(0\) and \(H(0)\). Then
\[
|H(H(0))-H(0)|\le L_p|H(0)|.
\]
If \(H(0)\ne0\), the two values have the same sign; if \(H(0)=0\), both vanish. In both cases the strict itinerary \(\tau_{2p}=-\tau_p\) is excluded.

The constant \(L_p\) is obtained by propagating the initial segment \([-r,r]\), \(r\ge|H(0)|\), alongside the critical-orbit boxes and multiplying the outward derivative bounds. The maps of the full band are bounded, retaining \eqref{hmtfg:escalado:eq:61}.

\subsubsection{Coincidence theorem for the accumulation parameter}
\label{hmtfg:escalado:16.4}

\textbf{Theorem.} For the uniform HMT family \eqref{hmtfg:escalado:eq:1}, the limit of the original first-root selection coincides with the certified stable crossing:
\begin{equation}
\boxed{\lambda_n\nearrow a_*=\lambda_*,
\qquad 1.874038<\lambda_*<1.874039.}
\label{hmtfg:escalado:eq:63}
\end{equation}

\textbf{Proof.} \eqref{hmtfg:escalado:eq:50}–\eqref{hmtfg:escalado:eq:52} give \(\lambda_n\uparrow a_*\le\lambda_*\), with both parameters in \(J\). If the inequality were strict, \eqref{hmtfg:escalado:eq:54}–\eqref{hmtfg:escalado:eq:59} would produce a return \(f_N\in\mathcal C_-\) realizing \(\tau\). The complete cover of section~\ref{hmtfg:escalado:16.3} excludes precisely that possibility. Hence \(a_*=\lambda_*\).

The proof controls every depth through a finite first exit and a uniform cover of its region. The identity \eqref{hmtfg:escalado:eq:63} links the global selector’s limit to the transverse crossing.


\subsection{Identification of the original centers and global scaling}
\label{hmtfg:escalado:17}

The final step preserves the selection of \eqref{hmtfg:escalado:eq:50}. The local continuation and the original parameters are compared through their exact period and their hybrid class; coincidence is obtained through uniqueness in a fixed parameter neighborhood.

\subsubsection{Degree-two analytic realization and transversality}
\label{hmtfg:escalado:17.1}

The fixed point \(g\) of section~\ref{hmtfg:escalado:3} is even, analytic, with a quadratic critical point and stationary doubling combinatorics. Theorem 2.1 of \href{https://annals.math.princeton.edu/wp-content/uploads/annals-v164-n3-p01.pdf}{de Faria–de Melo–Pinto}, applied to that single combinatorics, gives a one-element limit set, with a proper holomorphic extension of degree two between nested topological disks. Its attraction on the interval identifies that element with \(g\), since \(Rg=g\).

Lemma 3.2 of the same source extends a fixed number of returns to a full analytic neighborhood of \(g\), with degree-two images and positive modulus. The norm \eqref{hmtfg:escalado:eq:5} controls the uniform norm on a sufficiently small complex neighborhood of \([-1,1]\), because \(|(x^2-1)/2.5|<1\) there. The convergence \eqref{hmtfg:escalado:eq:3} enters that neighborhood; analytic dependence of the returns provides, for some fixed integer \(d\), a family
\begin{equation}
\mathcal F_\lambda=R^d f_\lambda
\label{hmtfg:escalado:eq:64}
\end{equation}
of degree-two maps, analytic in a complex neighborhood of \(\lambda_*\). The disks and the number of returns are fixed before varying \(\lambda\) in that neighborhood.

The parameter tangent belongs to the expanding cone, whereas the stable sheet satisfies the slope bound of section~\ref{hmtfg:escalado:5}. The returns preserve that separation. The degree-two realization locally identifies the stable sheet with the hybrid class of the infinitely renormalizable map; hence \eqref{hmtfg:escalado:eq:64} is transverse to that class.

The change of analytic space is justified in detail in \ref{hmtfg:clasificacion:8.2}, devoted to the transport of transversality. A direction tangent to the hybrid class would produce decreasing renormalized tangents, by uniform convergence on a hybrid disk and Cauchy’s estimate. The certified direction, by contrast, satisfies \(|\dot f(1)|=|\dot u|/10\) and grows under the returns by the expanding cone. This common evaluation of both representatives proves transversality without assuming equivalence of their norms.

\subsubsection{Uniqueness in a fixed neighborhood}
\label{hmtfg:escalado:17.2}

For \(n>d\), the classification of section~\ref{hmtfg:escalado:15} gives
\begin{equation}
K(\mathcal F_{\lambda_n})=W_{n-d}0.
\label{hmtfg:escalado:eq:65}
\end{equation}
Straightening a degree-two map preserves the critical orbit; \eqref{hmtfg:escalado:eq:65} identifies its hybrid class with that of the canonical quadratic center of period \(2^{n-d}\). The a priori bounds for the real doubling sequence are those used in Lyubich’s universality theorem.

\href{https://www.maths.tcd.ie/EMIS/journals/Annals/149_2/lyubich.pdf\#page=69}{Lyubich’s Theorem 7.4, pp. 387–388} gives a single intersection with each of those classes, at sufficiently high levels, on the analytic transversal and within a fixed neighborhood of the accumulation parameter.

By \eqref{hmtfg:escalado:eq:63}, \(\lambda_n\) eventually enters that neighborhood. The local continuation of section~\ref{hmtfg:escalado:8}, indexed by the same period, belongs to the same hybrid class. Uniqueness implies
\begin{equation}
\boxed{\lambda_n=\widehat\mu_n\quad(n\ge n_0),}
\label{hmtfg:escalado:eq:66}
\end{equation}
where \(\widehat\mu_n\) denotes the local center of physical period \(2^n\). This is a reindexing of \(\mu_j\) by a fixed shift determined by the returns and the target of section~\ref{hmtfg:escalado:8}. The original selection of \(\lambda_n\) remains intact.

The integer \(n_0\) is existential in this proof. Continuation at every level and the word \(W_n0\) are already proved in \eqref{hmtfg:escalado:eq:50}; \eqref{hmtfg:escalado:eq:66} provides the metric identity needed for the limit.

\subsubsection{Scaling theorem for the first HMT roots}
\label{hmtfg:escalado:17.3}

\textbf{Theorem.} For the uniform family \eqref{hmtfg:escalado:eq:1}, let \(\lambda_n\) be the successive first new roots of exact critical period \(2^n\). There exist \(C>0\), \(0<\theta<1\) and \(M<\infty\) such that, for every sufficiently large \(n\),
\begin{equation}
\lambda_*-\lambda_n=C\delta_{\mathrm F}^{-n}(1+e_n),
\qquad |e_n|\le M\theta^n,
\label{hmtfg:escalado:eq:67}
\end{equation}
and
\begin{equation}
\boxed{\lim_{n\to\infty}
\frac{\lambda_{n-1}-\lambda_{n-2}}
{\lambda_n-\lambda_{n-1}}=\delta_{\mathrm F}.}
\label{hmtfg:escalado:eq:68}
\end{equation}

\textbf{Proof.} \eqref{hmtfg:escalado:eq:66} transports the local scaling \eqref{hmtfg:escalado:eq:24} to the original roots; the finite index shift is absorbed into \(C,M\). Its sign is positive because \(\lambda_n<\lambda_*\).

The power-law remainder also admits a direct justification through the holonomy of hybrid classes. Lyubich’s Lemma 7.3 gives \(C^{1+\beta}\) regularity, with nonzero derivative, on the transversal at the accumulation point. The one-dimensional unstable dynamics is analytic and has multiplier \(\delta_{\mathrm F}>1\); its linearizing coordinate places successive centers at distances proportional to \(\delta_{\mathrm F}^{-n}\). Composing with the inverse holonomy gives a nonzero linear term and a remainder \(O(\delta_{\mathrm F}^{-(1+\beta)n})\). Thus one may take \(\theta=\delta_{\mathrm F}^{-\beta}\).

Writing
\[
r_n=\frac{\delta_{\mathrm F} e_{n-1}-e_n}{\delta_{\mathrm F}-1},
\qquad K=\frac{M(\delta_{\mathrm F}+\theta)}{\delta_{\mathrm F}-1},
\]
one obtains
\[
\lambda_n-\lambda_{n-1}
=C(\delta_{\mathrm F}-1)\delta_{\mathrm F}^{-n}(1+r_n),\qquad
|r_n|\le K\theta^{n-1}.
\]
Therefore,
\begin{equation}
\left|
\frac{\lambda_{n-1}-\lambda_{n-2}}
{\lambda_n-\lambda_{n-1}}-\delta_{\mathrm F}
\right|
\le
\frac{\delta_{\mathrm F} K(1+\theta)\theta^{n-2}}
{1-K\theta^{n-1}},
\quad K\theta^{n-1}<1.
\label{hmtfg:escalado:eq:69}
\end{equation}
The right-hand side tends to zero, proving \eqref{hmtfg:escalado:eq:68}.

\subsubsection{Chain of results and scope of the closure}
\label{hmtfg:escalado:17.4}

The effective composition is: APP–TRIT–TPK dodecaphase reader; profile \eqref{hmtfg:escalado:eq:1}; ordered transport of returns \eqref{hmtfg:escalado:eq:38}–\eqref{hmtfg:escalado:eq:40}; classification and existence of the first roots \eqref{hmtfg:escalado:eq:50}; localization \eqref{hmtfg:escalado:eq:52}; first-exit isolation \eqref{hmtfg:escalado:eq:63}; eventual identity of centers \eqref{hmtfg:escalado:eq:66}; scaling \eqref{hmtfg:escalado:eq:67}–\eqref{hmtfg:escalado:eq:69}. The Jacobi representation \eqref{hmtfg:escalado:eq:43} preserves this reader’s full profile and its marks remain in the antecedent.

The spatial expansion of the HMT reader and nonadic refinement provide the genealogy and evaluation at arbitrary depth. Scaling between parameters is proved through the return dynamics, transversality and isolation that we have just composed. The cited analytic proofs act on the family already constructed; their comparison values do not select its phases or coefficients.

The closure concerns the original global selector of the uniform sector \(\eta=0\), with the listed theorem dependencies. The asymptotic law retains its existential quantifiers. The width of a finite-root box measures evaluation precision; the estimate \eqref{hmtfg:escalado:eq:69} expresses convergence toward the limit.


The programs and certificates in the reproducibility supplement make it possible to reproduce the finite parts of the proof. Continuation at every depth rests on the inductions, the uniform first-exit argument and the identification theorems, with their explicit domains.

% Fragmento integrable. Preambulo anfitrion: amsmath, amssymb, longtable, hyperref.
\section{Analytic identification of the original centers}
\label{hmtfg:fragmento:identificacion-analitica-de-los-centros-originales}

\subsection{Degree-two realization and metric identification of the centers}
\label{hmtfg:clasificacion:8}

The equality between the first parameter of the infinite itinerary and the stable crossing, proved by first exit, provides
\begin{equation}
\boxed{a_*:=\min\Lambda_\tau=\lambda_*.}
\label{hmtfg:clasificacion:eq:H}
\end{equation}

The family, dodecaphase profile, returns and selector remain exactly as in sections~\ref{hmtfg:clasificacion:1}--\ref{hmtfg:clasificacion:5}. The external results used below recognize the dynamics of that already constructed object. Their universal values do not enter the selection of coefficients or parameters of the HMT family.

\textbf{Metric transfer theorem.} The identification \eqref{hmtfg:clasificacion:eq:H}, together with the previously proved entry, admissibility and transversality bounds, implies that the original sequence of first new roots satisfies, for certain constants \(C>0\) and \(0<\theta<1\),
\begin{equation}
\boxed{\lambda_*-\lambda_n=C\delta_{\mathrm F}^{-n}
\bigl(1+O(\theta^n)\bigr),}
\label{hmtfg:clasificacion:eq:14}
\end{equation}
where \(\delta_{\mathrm F}>1\) is the unstable eigenvalue of the doubling operator at the recognized fixed point. In particular,
\begin{equation}
\frac{\lambda_{n-1}-\lambda_{n-2}}
     {\lambda_n-\lambda_{n-1}}
=\delta_{\mathrm F}+O(\theta^n).
\label{hmtfg:clasificacion:eq:15}
\end{equation}
The sequence eventually coincides with the centers of the local continuation, once both indices refer to the same physical period. The proof retains the first-root selector of section~\ref{hmtfg:clasificacion:4}.

\subsubsection{Obtaining a common quadratic-like domain}
\label{hmtfg:clasificacion:8.1}

Write \(g\) for the renormalization fixed point and
\(\Gamma(\lambda)=R^4f_\lambda\). The analytic chart used in the certificate is
\[
f(x)=1-x^2V\!\left(\frac{x^2-1}{2.5}\right),\qquad
V(z)=\frac{u}{10}+\sum_{k\ge1}\nu_kz^k,
\]
with difference norm \(|\Delta u|+\|\Delta\nu\|_1\). The admissibility inequalities prove that \(g\) is even, analytic, unimodal, quadratic and of doubling type; moreover, \(g(x)=\varphi(x^2)\) with \(\varphi'<0\) in \([0,1]\).

Theorem 2.1 of de Faria–de Melo–Pinto, applied to the combinatorial set consisting solely of doubling, provides a one-point limit set with a quadratic-like extension. Since \(Rg=g\), the real convergence in that theorem identifies this point with \(g\). Lemma 3.2 provides, in an analytic neighborhood of \(g\), a finite renormalization iterate with a quadratic-like domain and uniformly positive modulus. The locations are sections 2.1–2.3, pp. 736–738, and section 3.1, Lemma 3.2, p. 745, of \href{https://annals.math.princeton.edu/wp-content/uploads/annals-v164-n3-p01.pdf}{de Faria–de Melo–Pinto, \emph{Global hyperbolicity of renormalization for \(C^r\) unimodal mappings}}.

To check membership in the neighborhood of that lemma, let \(\Omega_\epsilon\) be a sufficiently small complex neighborhood of \([-1,1]\), with
\[
\rho_\epsilon:=\sup_{x\in\Omega_\epsilon}
\left|\frac{x^2-1}{2.5}\right|<1,
\qquad M_\epsilon:=\sup_{x\in\Omega_\epsilon}|x|^2<\infty.
\]
The restriction of the certified chart satisfies
\begin{equation}
\|\Delta f\|_{\Omega_\epsilon}
\le M_\epsilon\left(\frac{|\Delta u|}{10}
          +\rho_\epsilon\|\Delta\nu\|_1\right)
\le C_\epsilon\|\Delta f\|_{\mathcal A}.
\label{hmtfg:clasificacion:eq:16}
\end{equation}
The certificate gives
\[
\|R^m\Gamma(\lambda_*)-g\|_{\mathcal A}
<0.001666(0.83996)^m.
\]
A fixed integer \(m\) is therefore chosen to place this image in the neighborhood of Lemma 3.2. Continuity of the finite composition provides an interval \(I_*\) around \(\lambda_*\) where the same operations are defined. If \(N\) is the iterate supplied by that lemma, we set
\begin{equation}
d=4+m+N,\qquad
\mathcal F_\lambda=R^df_\lambda,\quad \lambda\in I_*.
\label{hmtfg:clasificacion:eq:17}
\end{equation}
The family \(\mathcal F\) has quadratic-like representatives of degree two, on a common domain and with uniformly positive modulus. Its real-analytic dependence admits a local complexification of the parameter. The number \(d\) is fixed and counts the doublings performed before applying the parameter theorem.

This construction fixes a common domain and a finite number of returns \(d\), sufficient to apply the parameter theorem.

\subsubsection{Transport of certified transversality}
\label{hmtfg:clasificacion:8.2}

The direction of \(\Gamma\) satisfies the cone condition
\[
\|\dot\nu\|_1\le\tfrac14|\dot u|,\qquad \dot u\ne0.
\]
The certified differential bounds imply invariance of this cone and
\begin{equation}
|\dot u_{k+1}|\ge4.381|\dot u_k|
\label{hmtfg:clasificacion:eq:18}
\end{equation}
along the orbit of \(\Gamma(\lambda_*)\). All its points remain in the certified domain. In particular, the direction transported through the finite steps of \eqref{hmtfg:clasificacion:eq:17} retains a nonzero expanding component.

The hybrid class of the fixed point coincides with its stable manifold in the space of quadratic-like germs: \href{https://www.maths.tcd.ie/EMIS/journals/Annals/149_2/lyubich.pdf}{Lyubich, \emph{Feigenbaum–Coullet–Tresser universality and Milnor's hairiness conjecture}, Theorem 6.1, pp. 371–372}. Transversality remains valid on passing to that space, as can be checked without identifying different Banach norms. Indeed, if \(\partial_\lambda\mathcal F_{\lambda_*}\) were tangent to the hybrid class, it would be the tangent of an analytic disk contained in it. Uniform exponential convergence of renormalization on a compact subdisk, followed by Cauchy’s estimate in its parameter, would make the evaluation at \(x=1\) of its renormalized tangents converge to zero. However, in the preceding chart,
\[
\dot f(1)=-\dot u/10,
\]
and \eqref{hmtfg:clasificacion:eq:18} makes its absolute value grow. Both derivatives represent the same germ and have the same real evaluation. The contradiction proves
\begin{equation}
\partial_\lambda\mathcal F_{\lambda_*}
\notin T_{\mathcal F_{\lambda_*}}\mathcal H_{c_\infty}.
\label{hmtfg:clasificacion:eq:19}
\end{equation}
Thus, \(S=\{\mathcal F_\lambda\}\) is a complex-analytic transversal to the hybrid class of the doubling limit.

\subsubsection{Eventual identification of the first roots by their period}
\label{hmtfg:clasificacion:8.3}

Denote by \(c_r\) the canonical quadratic center of period \(2^r\), obtained through \(r\) successive doublings of the period-one center. This symbol denotes the parameter of the quadratic recognition family; it is distinct from the orbital functions \(c_j(\lambda)=f_\lambda^j(0)\) of the preceding sections. We write \(c_r^{\rm quad}\) to preserve the distinction.

The global continuation theorem and \eqref{hmtfg:clasificacion:eq:H} give \(\lambda_n\to\lambda_*\), hence \(\lambda_n\in I_*\) for sufficiently large \(n\). The classification and restrictive intervals of section~\ref{hmtfg:clasificacion:2} show that, when \(n>d\),
\begin{equation}
\mathcal F_{\lambda_n}=R^df_{\lambda_n}
\quad\text{has exact critical period }2^{n-d}
\quad\text{and itinerary }W_{n-d}0.
\label{hmtfg:clasificacion:eq:20}
\end{equation}
Its critical orbit remains in the invariant real interval, contained in the quadratic-like domain; its filled Julia set is therefore connected. Straightening preserves the postcritical orbit and the real branches. Following its restrictive returns gives \(n-d\) canonical doublings and, finally, a fixed critical point. The latter straightens to the center \(0\); the successive inverses of the straightening maps of the doubling copies give the unique center \(c_{n-d}^{\rm quad}\). This is the concrete application of the combinatorial parameterization and the straightening homeomorphism of the copies described in Lyubich, sections 5.1–5.3, pp. 362–364. We conclude
\begin{equation}
\mathcal F_{\lambda_n}\in\mathcal H_{c_{n-d}^{\rm quad}}.
\label{hmtfg:clasificacion:eq:21}
\end{equation}

Lyubich’s Theorem 7.4, pp. 387–388, gives a unique local intersection of \(S\) with each class \(\mathcal H_{c_r^{\rm quad}}\) for sufficiently large \(r\). Its assumptions of a priori complex bounds hold for the real quadratic doubling cascade, according to Theorem 5.6, p. 367, of the same work. Write \(\zeta_r\) for the parameter of that intersection. Equations \eqref{hmtfg:clasificacion:eq:20}–\eqref{hmtfg:clasificacion:eq:21}, convergence to the crossing and local uniqueness imply
\begin{equation}
\boxed{\lambda_n=\zeta_{n-d}\quad\text{for every }n\text{ sufficiently large}.}
\label{hmtfg:clasificacion:eq:22}
\end{equation}
The uniqueness used here concerns \textbf{each canonical hybrid class} in a fixed neighborhood of the crossing. It is not inferred merely from the existence of a local cascade or from uniqueness of the crossing with the stable sheet. Equation \eqref{hmtfg:clasificacion:eq:22} identifies the original selector, already defined globally, with the local intersections; it does not introduce a replacement selector.

In particular, let \(\mu_j\) be the local cascade of section \ref{hmtfg:escalado:8}, and let \(s\) be the fixed integer determined by its exact physical period:
\[
\operatorname{per}_{\rm crit}(f_{\mu_j})=2^{j+s}.
\]
In the presentation with curve \(R^4f_\lambda\) and period-two target, \(s=5\); the additional fixed steps used to localize the target modify \(s\). By \eqref{hmtfg:clasificacion:eq:21} and the same uniqueness,
\begin{equation}
\mu_j=\zeta_{j+s-d},\qquad
\boxed{\lambda_n=\mu_{n-s}\quad(n\text{ sufficiently large}).}
\label{hmtfg:clasificacion:eq:23}
\end{equation}
The shift \(d\) in the quadratic-like preparation and the shift \(s\) in the target periods are distinct objects. Identification is made through the physical period, not by equating those two integers.

\subsubsection{Power-law remainder through transverse holonomy}
\label{hmtfg:clasificacion:8.4}

Lyubich’s Lemma 7.3, p. 384, with proof on pp. 384–387, provides \(C^{1+\beta}\)-conformal regularity, for some \(\beta>0\), of the holonomy between transversals to the hybrid class of the Feigenbaum point. Its definition on p. 384 concerns the connectedness loci on the transversals; these contain, in particular, the centers considered here. We may reduce the exponent so that \(0<\beta\le1\).

Let \(h\) be the holonomy from \(S\) to the unstable manifold, and let \(z\) be an analytic coordinate linearizing the one-dimensional restriction of renormalization:
\begin{equation}
z(g)=0,\qquad z(Rq)=\delta_{\mathrm F} z(q).
\label{hmtfg:clasificacion:eq:24}
\end{equation}
The existence of this coordinate follows by applying the linearization of a repelling holomorphic germ, with multiplier \(\delta_{\mathrm F}>1\), to the analytic unstable manifold. The expanding direction of \eqref{hmtfg:clasificacion:eq:18}, cone invariance and the stationary return identify that multiplier with the positive unstable eigenvalue used in the certificate.

If \(q_r\) is the center of the class \(\mathcal H_{c_r^{\rm quad}}\) lying on the unstable manifold, compatibility of straightening and renormalization gives \(Rq_r=q_{r-1}\). This identity is the one used in the proof of Theorem 7.4, p. 388. By \eqref{hmtfg:clasificacion:eq:24}, there exists \(b\ne0\), absorbing any fixed initial index, such that
\begin{equation}
z(q_r)=b\delta_{\mathrm F}^{-r}
\label{hmtfg:clasificacion:eq:25}
\end{equation}
for every sufficiently large \(r\).

The regularity of \(h\) and transversality \eqref{hmtfg:clasificacion:eq:19} give, on the connectedness locus near \(\lambda_*\),
\begin{equation}
H(\lambda):=z\bigl(h(\mathcal F_\lambda)\bigr)
=A(\lambda-\lambda_*)
+O\!\left(|\lambda-\lambda_*|^{1+\beta}\right),
\qquad A\ne0.
\label{hmtfg:clasificacion:eq:26}
\end{equation}
Since \(H(\zeta_r)=b\delta_{\mathrm F}^{-r}\), the lower estimate
\(|H(\lambda)|\ge |A|\,|\lambda-\lambda_*|/2\) in a sufficiently small neighborhood first implies
\(|\zeta_r-\lambda_*|=O(\delta_{\mathrm F}^{-r})\). Substituting this bound into the remainder of \eqref{hmtfg:clasificacion:eq:26} yields
\begin{equation}
\zeta_r-\lambda_*=\frac bA\delta_{\mathrm F}^{-r}
+O\!\left(\delta_{\mathrm F}^{-(1+\beta)r}\right).
\label{hmtfg:clasificacion:eq:27}
\end{equation}
We apply \eqref{hmtfg:clasificacion:eq:22} with \(r=n-d\). The finite shift modifies the leading coefficient and the remainder constant. Since \(\lambda_n<\lambda_*\), we obtain \eqref{hmtfg:clasificacion:eq:14}, with
\[
C=-\frac bA\delta_{\mathrm F}^d>0,
\qquad \theta=\delta_{\mathrm F}^{-\beta}\in(0,1).
\]
Finally,
\[
\lambda_n-\lambda_{n-1}
=C(\delta_{\mathrm F}-1)\delta_{\mathrm F}^{-n}\bigl(1+O(\theta^n)\bigr),
\]
which proves \eqref{hmtfg:clasificacion:eq:15}. \(\square\)


\subsubsection{Quantitative scope of the theorem}
\label{hmtfg:clasificacion:8.5}

The identification of limits proved by first exit activates metric transfer for the global first roots of the uniform sector. The operator rate \(0.83996\) controls the renormalized maps on the stable sheet. The parameter rate \(\theta=\delta_{\mathrm F}^{-\beta}\) follows from the regularity of transverse holonomy. Both estimates retain their variable and domain.

The theorem establishes the existence of \(\beta,\theta,C\), the neighborhood \(I_*\), the finite number of returns \(d\) and an index \(n_0\) beyond which the centers coincide. Their individual numerical values are not assigned in this asymptotic statement. This formulation completes the limit and its power-law remainder; numerical evaluation of those witnesses is a distinct quantitative specialization.

\section{Return, memory and universality of the dodecaphase cascade}
\label{sec:hmt-feig-interpretacion}

\subsection{The generating structure and the critical reading}

The Feigenbaum construction in Triadic Modular Holography links an
oriented arithmetic organization to an asymptotic scaling law.
Its starting point is the effective composition of APP, TRIT and TPK. APP evaluates
sum and product on the two sheets of the digit lattice and preserves the
quotients and remainders of those operations. TRIT determines the local regime and
the orientation of its reading. TPK selects and transports the states,
updates their records and prolongs their routes with carry, boundary and memory.
The joint discrete structure of the continuum provides the common residence
of those states and of the readers acting on them.

The prolongation
\[
w_6\longrightarrow w_{12}\longrightarrow w_{18}
\longrightarrow w_{24}\longrightarrow w_{30}
\longrightarrow R_{36}\longrightarrow G_9
\]
organizes this persistence. The initial seed is lifted while preserving its
prefixes; successive lifts change the regime without losing orientation
or memory; the coinductive boundary allows the construction to continue, and the
nine charts bring together its local readings. Nonadic holonomy expresses
phase return accompanied by an increase in memory. Thus, a visible
turn preserves its position within the complete history. From the outset,
the return operation acquires two inseparable components:
reiteration of a rule and preservation of its provenance.

The critical reader fixed on the oriented dodecaphase wheel turns this
structure into twelve phases with a uniform trace. Their representatives are
\[
\phi_m=\frac{(3m-1)\pi_{\mathrm{HMT}}}{18},
\qquad w_m=\frac1{12},\qquad 1\le m\le12.
\]
The choice of oriented representative matters because the second moment
uses its magnitude. The angular quotient records the phase, while the
representative retains the information needed to normalize it. Coordinate
\(\pi_{\mathrm{HMT}}\) comes from the coinductive readers of the same
APP--TRIT--TPK core. Within its correlated image, coordinates
\(\pi,\varphi,e\) and the dodecaphase closure of \(\alpha_{\mathrm{HMT}}\)
retain their constructive order and their own domains. The critical reader
considered here uses the angular coordinate already generated.

Normalization is determined through an integer sum:
\[
\sum_{m=1}^{12}(3m-1)^2=5394=6\cdot899,
\qquad k_m=\frac{2(3m-1)}{\sqrt{899}}.
\]
The common factor \(\pi_{\mathrm{HMT}}\) cancels when the second moment is fixed.
This cancellation preserves the provenance of the phases and produces a
particularly precise representation of their normalized reading. We obtain
\[
u_0(x)=\frac1{12}\sum_{m=1}^{12}\bigl(1-\cos(k_mx)\bigr),
\qquad f_\lambda(x)=1-\lambda u_0(x),
\qquad \sum_mw_mk_m^2=2.
\]
The profile is even and entire; it satisfies \(u_0(0)=u_0'(0)=0\) and
\(u_0''(0)=2\). For \(0<\lambda\le2/u_0(1)\), the family maps
\([-1,1]\) into itself and has a quadratic critical maximum with
\(f_\lambda''(0)=-2\lambda\). Parameter \(\lambda\) ranges over this family
after the reader, orientation, weights and scale have been fixed. The universal
number recognized at the end is thus separated from the selection
of the coefficients.

\subsection{Twelve phases, moments and spectral representation}

The profile admits a second expression bringing its phases together in a finite
operator. The nodes
\[
s_m=k_m^2=\frac{4(3m-1)^2}{899}
\]
determine the positive measure \(\mu_0=12^{-1}\sum_m\delta_{s_m}\). Its
moments \(M_r=\int s^r\,d\mu_0(s)\) define strictly positive Hankel forms
up to rank twelve. This positivity follows from evaluating
a polynomial at twelve distinct nodes: a polynomial of degree less than
twelve retains at least one nonzero evaluation. At degree twelve, the product
of the twelve nodal factors annihilates the norm and determines the reader’s
rank exactly.

Orthogonalization of these moments produces a Jacobi matrix
\(J_{12}\), real symmetric, with positive subdiagonal entries. Its functional calculus
recovers the complete profile:
\[
u_0(x)=e_0^{\mathsf T}
\bigl[I-\cos(x\sqrt{J_{12}})\bigr]e_0.
\]
The cosine sum and the Jacobi operator are two realizations of the same
reader. The moments retain all frequencies and their weights; the
orientation, route and memory marks remain in the state that
generates them. Rank twelve identifies the spectral dimension of this
specific reading. The depth of its prolongations belongs to another dimension
of the construction: the length of compatible histories and the resolution
at which they are evaluated.

The analytic geometry is also made explicit. For the weighted
deformation \(u_\eta\), in the strip \(|\eta|\le1/40\), the weights remain
positive and the second moment remains normalized. In
\(|\operatorname{Re}z|<\pi/3\) there exists an odd, univalent holomorphic function
\(b_\eta\), normalized by \(b_\eta'(0)=1\), such that
\(u_\eta=b_\eta^2\). The wheel therefore produces a conformal deformation
of the quadratic fold. Its dynamical conjugacy preserves that deformation
through \(b_\eta(1-\lambda w^2)\). The parameter-scaling theorem
developed here concerns the uniform member \(\eta=0\); the
structural properties of the weighted strip and the local persistence of
its roots retain their specific quantifiers.

\subsection{Reading depth and recovery of histories}

Nonadic depth refines the reading of a parameter, whereas dynamical
doubling changes the return horizon. Refinement with
carry expresses the first operation through
\[
L_c(x,y)=(Bx-c,By+c),\qquad
C_m=\sum_{j=1}^m c_jB^{m-j},
\]
\[
(x_m,y_m)=(B^mx_0-C_m,B^my_0+C_m).
\]
The outer support grows as \(B^m\), and the normalized record cylinder
has diameter \(B^{-m}\). The integer pair, depth and boundary
mark preserve reconstruction of quotients, remainders and prefixes.
Outer growth and inner resolution thus describe compatible
operations on the same history.

The exact composition
\(L_D^{[b]}L_C^{[a]}=L_{B^bC+D}^{[a+b]}\) allows a history to be grouped into
blocks without altering its data. A doubling return combines two
transitions; the nonadic calendar combines nine. Both groupings are
compared over a common horizon of \(9\,2^N\) events, preserving
the order of the subblocks and their carries.

The contribution of memory can be observed directly in the profile’s
inverse branches. If \(v=(u_0|_{[0,1]})^{-1}\), the two branches are
\[
\beta_\sigma(y)=\sigma v\!\left(\frac{1-y}{\lambda}\right),
\qquad \sigma\in\{-1,+1\}.
\]
The final value, together with the last sheet, reconstructs the preceding value. A
sequence of sheets allows all intermediate steps to be recovered. The branch
word accompanies the APP--TRIT--TPK records of route, boundary, carry and
memory; each record retains its function. This concrete recovery
explains what information is required to invert the critical reading.

The resolution of each return also has an explicit bound. With
\(t=\lambda u_\eta(1)/2\), \(z=(x+1)/2\) and
\[
F_{t,\eta}(z)=1-t\frac{u_\eta(2z-1)}{u_\eta(1)},
\]
one obtains
\[
\left|F_{t,\eta}^{q}(1/2)-F_{s,\eta}^{q}(1/2)\right|
\le S_q|t-s|,\qquad S_q=\sum_{j=0}^{q-1}(6\sqrt2)^j<9^q.
\]
A base-nine cylinder of depth \(m=2^N+r\) therefore allows the
return of horizon \(2^N\) to be read with diameter less than \(9^{-r}\). This
estimate gives quantitative content to arbitrary depth: for every
requested horizon and precision, it determines a sufficient depth.
Evaluation of the analytic functions and rounding have their own
remainders, which propagate together with the parameter uncertainty.

Vacancies enter the capacity balance of the record.
Comparison between blocks of 729 and 1000 possibilities distinguishes accumulated
capacity from chronological length, preserving the steps at which the former
remains constant. This accounting determines the space required to
preserve a reading and its depth. The admissibility of histories
follows, in turn, from the survival rules. Both operations
contribute to the construction and retain distinct mathematical objects.

\subsection{Nonadic chronology and doubling combinatorics}

The critical doubling word has an expression at every depth:
\[
\tau_j=(-1)^{v_2(j)},\qquad
\tau_{2j}=-\tau_j,\qquad \tau_{2j+1}=+1.
\]
Its prefixes satisfy \(W_{n+1}=W_n(-1)^nW_n\). The rule determines how
the two copies of the return are oriented and where their separation appears.
The nonadic reading of chronology preserves the complete integer
\(K=\rho_9(K)+9q_9(K)\). Reducing it modulo \(2^n\), the holonomy
that advances nine events induces translation by nine on that cycle. Since
nine is odd, this translation visits all its positions. The
permutation \(j\mapsto9j\pmod{2^n}\), which conjugates it to unit advance,
preserves the dyadic valuation of each nonzero position. The chronological remainder and quotient
jointly support that compatibility.

Bilateral memory also has an operator realization compatible
with refinement. On cycles of length \(2^n\), the shift operator
\(C_n\) is incorporated into
\[
\mathbb U_n=P_0\otimes I+P_1\otimes C_n,
\]
with the complementary projectors of the nonadic reader. In the limit, its
visible and complementary components satisfy
\[
\|T_{\infty,a}v\|^2+\|D_{\infty,a}v\|^2=\|v\|^2,
\qquad
v=T_{\infty,a}^*T_{\infty,a}v+D_{\infty,a}^*D_{\infty,a}v.
\]
These identities describe preservation and recovery within the declared space
and representation. Contraction of the renormalized maps,
which will appear later, measures another operation: the approach of their profiles
to a stable form. The history preserved by the lift and the
convergence of a functional reading are articulated at different levels.

\subsection{The ordered chart and global root selection}

The Archimedean readout of the HMT continuum transports the operations
through an isomorphism of complete ordered fields, within the declared
universe \(\mathfrak G\):
\[
\iota:\mathbb R_{\mathrm{HMT}}^{\mathfrak G}
\longrightarrow\mathbb R^{\mathfrak G}.
\]
The internal construction of the cosine through its series, the positive square root of 899,
the profile and its iterates precedes zero selection. Preservation
of sum, product, order and convergent limits gives
\[
\iota\bigl(S_n^{\mathrm H}(a)\bigr)=S_n\bigl(\iota(a)\bigr),
\qquad S_n(\lambda)=f_\lambda^{2^n}(0).
\]
Surjectivity covers all roots in the domain of that chart; order
preserves preceding returns, exact period and the condition of
lying to the right of the preceding root. In this way,
the minimum of each candidate set is also transported. The genealogy reaches
the effective parameter selection, with the second projection and its
marks preserved on the same antecedent.

The selector is fixed by \(\lambda_1=1/u_0(1)\) and by the first
new root of exact critical period \(2^n\) lying to the right of
\(\lambda_{n-1}\). The itinerary classification proves that each selected
root has word \(W_n0\). Restrictive returns halve its
period while preserving unimodal order, and induction restores
the complete word. Analyticity of returns that are not identically zero
isolates their zeros in each compact set;
comparison of signs before the first infinite-itinerary parameter
guarantees a new root at every level.

Let \(\Lambda_\tau\) be the nonempty compact set of parameters realizing
the infinite word, and let \(a_*\) be its minimum. It is proved that
\[
\lambda_n\nearrow a_*.
\]
Passage to the limit preserves each sign through a uniform bound at a fixed
horizon. If \(c_p=f_\lambda^p(0)\), the opposite signs of \(c_p\) and
\(c_{2p}\), together with \((f_\lambda^p)'(0)=0\), imply
\[
|c_p|>\frac2{B_p},\qquad
B_p=16^p\frac{16^p-1}{15}.
\]
Thus, the prefixes of increasing periods preserve the infinite word as
they accumulate. The proof uses the original sequence of first roots from
its global definition.

\subsection{The stable sheet and identification of the selected limit}

The normalized return
\[
(Rf)(x)=-\frac{f(f(-ax))}{a},\qquad a=-f(1)>0,
\]
combines two transitions, changes the scale and preserves the orientation needed
to restore a maximum of value one. The restrictive domains establish
its interpretation as a real return. The derivative of this operator includes
the variation of scale \(a\); that contribution determines the parameter
sensitivity of the complete transformation.

In the chart \(f(x)=1-x^2V((x^2-1)/(5/2))\), with coordinates \(u=10v_0\)
and \(\nu=(v_1,v_2,\ldots)\), the fourth return of the family enters a
quantified neighborhood of the fixed point \(g\). The curve crosses
its stable sheet transversely at a unique local parameter
\[
1.874038<\lambda_*<1.874039,
\]
and satisfies
\[
\|R^{4+n}f_{\lambda_*}-g\|_{\mathcal A}
<0.001666(0.83996)^n.
\]
The stable sheet is a graph of slope at most \(0.16\). The
parameter secants belong to a transverse expanding cone. Both
properties distinguish the decrease in stable deformations from the
growth of separation between parameters.

Identifying \(a_*\) with \(\lambda_*\) brings this local geometry together
with the global minimum. First, the interval between the eighth root
and the lower endpoint of the local box is excluded. Then the order
\(a_*\le\lambda_*\) allows a hypothetical separation to be studied through
the negative cone. While its transverse coordinate has absolute value less
than \(0.006\), the secant remains in the analytic domain and grows by a
factor between \(4.4034\) and \(8.0887\). A persistent separation
would therefore produce a first exit with amplitude between
\(0.006\) and \(0.0485322\).

The stable-orbit residual retains information more precise than its
norm: its components satisfy \(|e_u|\le0.16\|e_\nu\|_1\). This
relation allows their effects on the critical orbit to be bounded jointly.
The cover of the entire first-exit region shows in each case a
sign incompatible with \(\tau\), or a contracting return forcing two
successive dyadic critical values to have the same sign. The infinite
word requires opposite signs. The contradiction identifies the parameters:
\[
\lambda_n\nearrow a_*=\lambda_*.
\]
The proof combines an induction valid at every depth with a
uniform finite exclusion of the region that any separation
would reach. Its infinite scope rests on that composition.

\subsection{Universal scaling and the meaning of precision}

The degree-two holomorphic realization of the returns and the established
transversality allow their hybrid classes to be compared. For each sufficiently
high period, the canonical class has a unique local intersection
with the family. The global first roots already converge to the crossing and
preserve the doubling combinatorics; hence they eventually coincide
with the local centers when their physical periods are matched. The proof
transports the original selection to the universal metric law.

There exist \(C>0\), \(M<\infty\) and \(0<\theta<1\) such that
\[
\lambda_*-\lambda_n=C\delta_{\mathrm F}^{-n}(1+e_n),
\qquad |e_n|\le M\theta^n,
\]
and, consequently,
\[
\lim_{n\to\infty}
\frac{\lambda_{n-1}-\lambda_{n-2}}
{\lambda_n-\lambda_{n-1}}=\delta_{\mathrm F}.
\]
The factor \(\delta_{\mathrm F}\) is the expanding eigenvalue of the doubling operator
recognized at the end of the construction. The specific family preserves its
twelve phases, moments and provenance history; the limit of its
parameter ratios belongs to the same universality class of quadratic
criticality. This relation jointly explains the specificity of the
generator and the universality of the output.

Precision has three residences here. The root intervals
control their finite evaluation. The estimate
\(0.001666(0.83996)^n\) controls operator convergence on the stable
sheet. The term \(M\theta^n\) controls the transfer of that geometry to
the parameter scale through transverse holonomy. Each estimate
corresponds to its variable and operation. The eighth quotient is enclosed
around \(4.669212189150204154556096215889663928\), with width less
than \(5.808\cdot10^{-37}\); that width measures the evaluation of \(\delta_{\mathrm F,8}\).
The distance from \(\delta_{\mathrm F,8}\) to the limit requires quantifying the constants
of the parameter remainder, whose existence is already proved.

The notation keeps separate \(\delta_{\mathrm F}\), the Feigenbaum parameter-scaling
factor; \(\alpha_{\mathrm F}\), the spatial constant associated with its
renormalization problem; and \(\alpha_{\mathrm{HMT}}\), the corpus’s fine-structure
constant. The present result identifies \(\delta_{\mathrm F}\)
through the first roots of the uniform sector \(\eta=0\). The weighted
strip retains its structural and existence theorems, and its quantified
metric extension constitutes a distinct specific statement.

The contribution of the chain lies in continuity between generation,
reading, memory, selection and limit. Oriented arithmetic determines the
profile; the spectral representation preserves its twelve components; the
cylinders allow it to be resolved with arbitrary precision at each horizon;
history transport restores its paths; the ordered chart
preserves all roots and their minima; classification prolongs the
selector; and transverse expansion, together with stable contraction,
determines its asymptotic law. Universality emerges as a proved
property of this complete construction and its analytic realization.

\endgroup
